\chapter{Connected plant surrogates and operating optimization}
\label{ch:plant}

\section{Rationale}

Activated sludge process optimization concerns the response of a connected treatment system rather than an isolated effluent prediction. Biological reactors transform material, secondary clarification separates solids, and recycle streams return material to earlier treatment locations. Mixed liquor contains both wastewater constituents and suspended biomass. Clarifier overflow carries treated water from the system, while underflow supplies concentrated sludge for return and wasting. The concentrations and flow rates of these streams determine how much material is discharged, circulated, retained, and removed. Aeration, recycle, and wasting therefore influence both final effluent quality and the solids inventory maintained within the plant \citep{Henze2008,HartelPopel1992}.

Mechanistic plant models provide the physical basis for representing these connections. The benchmark of \citet{Alex2008} links biological reactors, secondary clarification, internal recycle, return sludge, and wasting within one treatment system. The component-transport analysis of \citet{JeppssonDiehl1996} further shows why clarifier outlets must retain a consistent biological composition. At the reactor scale, the structured prediction of \citet{Selerio2026} provides an inspectable approximation at a single treatment boundary. The first gap addressed in this chapter is therefore how to extend component prediction and physical enforcement from one reactor to shared streams and stored solids across a connected plant. Accurate or physically checked predictions at individual units do not, by themselves, establish consistency among those units.

The return-sludge connection illustrates this problem directly. A predicted underflow concentration re-enters the mixer through the return flow. If the mixer prediction implies a different contribution from the same stream, the two outputs represent different amounts of the same transported material. Both predictions can remain non-negative while their connection is physically inconsistent. Clarifier storage introduces an additional requirement because retained solids contribute to the plant inventory and determine solids retention time. A connected plant response must therefore retain internal concentrations, separator outlets, and stored solids rather than reduce the treatment system to final-effluent totals alone \citep{Alex2008,JeppssonDiehl1996}.

Joint projection provides a way to reconcile these quantities simultaneously. Its effect, however, must be evaluated at the treatment locations where the predictions will actually be used. An adjustment that reconciles recycle and separator balances may also alter internal biomass, nutrient concentrations, and final effluent quality. Location-specific errors, physical residuals, and computational effort are therefore all needed to assess the connected prediction. At the same time, satisfaction of the selected flow, inventory, and concentration constraints does not establish agreement with every nonlinear reaction or settling equation. Physical enforcement guarantees only the relations that are explicitly imposed \citep{KircherVotsmeier2025,Selerio2026}.

The second gap concerns the operating decision obtained from this connected response. Surrogates can reduce the computational effort required for repeated process evaluations \citep{McBride2019,BhosekarIerapetritou2018}. The simulation-based search of \citet{Fang2011} and the resource-recovery framework of \citet{Durkin2024} illustrate how learned models can support wastewater treatment decisions. Yet agreement over a holdout design does not guarantee agreement at the operating settings preferred by an optimizer. Search naturally concentrates evaluations in regions that appear favorable. If treatment is predicted too optimistically in such a region, the resulting error can influence the selected capacity, aeration, recycle, or wasting settings even when average prediction error over the broader domain is small.

Independent evaluation of selected controls is therefore a separate requirement. The surrogate comparisons of \citet{Pedrozo2025} distinguish predictive accuracy from optimization behavior, while the approximation-management methods of \citet{Hameed2026} retain checks against the original model. These contributions motivate evaluating the mechanistic plant response at the controls selected by each optimization route. The relevant evidence includes the resulting effluent concentrations, the objective calculated on a common reference basis, and satisfaction of the original engineering requirements. The availability of a candidate decision and the numerical certification of the search are also different outcomes. Failed searches must remain visible when the practical usefulness of a route is assessed.

The meaning of an operating comparison also depends on the priorities embedded in its objective. \citet{Araujo2013} examine constrained economic operation, while \citet{PadronPaez2020} and \citet{Aboagye2021} connect treatment performance with resource use. Two optimization routes can therefore prefer different controls because one accepts greater resource demand to obtain better treatment. A meaningful comparison must use common priorities, common control bounds, and common mechanistic evaluation. It should also separate treatment-quality and resource contributions so that the source of a lower objective remains visible. A lower objective is a result under the declared preferences; it does not imply that every effluent quantity or every resource demand is lower.

Capacity selection is particularly important in the present study. At fixed fresh-influent flow, changing HRT changes the biological reactor volume. The resulting problem therefore combines capacity selection with aeration, recycle, and wasting decisions. It is not equivalent to adjusting the operation of an installed plant whose tanks are already fixed. The resource terms used below are engineering proxies for comparing the declared alternatives rather than measured electricity bills or full facility costs. These boundaries define the practical meaning of the optimization results and prevent a numerical objective from being interpreted as a demonstrated field saving.

This chapter addresses the fourth research question in Section~\ref{sec:research_questions}. The connected surrogate is termed Extended ICSOR. The direct route solves a smooth nonlinear programming (NLP) problem and is termed Smooth NLP. Their selected controls are compared through independent mechanistic evaluation and a common engineering objective. The contribution is a joint component-state formulation that connects physical enforcement with explicit decision verification. The assessment separates network consistency, prediction accuracy, selected treatment outcomes, search reliability, and computational effort. This separation identifies where rapid approximation can support plant analysis and where mechanistic evaluation remains necessary.

\section{Theory}

\subsection{Hydraulic connections and reactor inventories}

\subsubsection{Flow through the treatment train}

Let $Q_0$ denote the fresh-influent flow in m$^3$ d$^{-1}$ and let $c_{\rm in}\in\mathbb R_+^F$ denote its component concentrations. The plant has $n_r$ completely mixed reactors and an $L$-layer clarifier. Reactor concentrations are $c_1,\ldots,c_{n_r}$, the mixer concentration is $m$, and clarifier overflow and underflow concentrations are $c_E$ and $c_U$. Because the overflow is the final treated-water outlet, $c_{\rm out}=c_E$.

The dimensionless ratios $r_I$, $r_R$, and $w$ describe mixed-liquor recycle, return activated sludge, and waste activated sludge flow relative to $Q_0$. Mixed-liquor recycle leaves the final reactor before the clarifier. Return sludge and waste sludge both leave through the clarifier underflow and therefore share the same underflow concentration. Each ratio is an actual flow divided by $Q_0$. Denote the reactor-train, clarifier-feed, underflow, and overflow flows by $Q_P,Q_C,Q_U,Q_E$. The mixer, reactor splitter, underflow splitter, and external water balance give
\begin{align}
\equationname{Reactor-train flow including both recycle streams}
Q_P&=Q_0+r_IQ_0+r_RQ_0,\\
\equationname{Clarifier-feed flow after mixed-liquor recycle withdrawal}
Q_C&=Q_P-r_IQ_0,\\
\equationname{Clarifier underflow split into return sludge and waste sludge}
Q_U&=r_RQ_0+wQ_0,\\
\equationname{Effluent flow after waste withdrawal}
Q_E&=Q_0-wQ_0.
\end{align}
The second equation removes the mixed-liquor recycle before clarification, while the third collects both destinations of the clarifier underflow. Dividing each balance by $Q_0$ gives
\begin{equation}
\equationname{Plant flow ratios relative to fresh-influent flow}
q_P=1+r_I+r_R,\qquad q_C=1+r_R,\qquad
q_U=r_R+w,\qquad q_E=1-w.
\label{eq:plant_flows}
\end{equation}
Here $q_P$ is the reactor-train flow ratio and $q_C$ is the clarifier-feed flow ratio. The identity $q_C=q_E+q_U$ confirms the clarifier water balance. The adopted domain requires $r_I\ge0$, $r_R\ge0$, $0\le w<1$, and $q_U>0$, which ensures that the outlet concentrations are defined by positive outlet flows.

Figure~\ref{fig:plant_configuration} places these balances on the physical connections of a plant with $n_r$ reactors and $L$ clarifier layers. The case-study counts and aeration assignments are introduced in Methodology. The two recycle streams return material to the mixer, whereas treated water and waste sludge cross the external plant boundary.

\begin{figure}[htbp]
\centering
\resizebox{\linewidth}{!}{%
\begin{tikzpicture}[>=Latex,
unit/.style={draw=plantblue,fill=plantblue!9,rounded corners,minimum width=13mm,minimum height=12mm,align=center,font=\small,thick},
flow/.style={->,thick,draw=plantblue},lab/.style={font=\scriptsize,fill=white,inner sep=2pt}]
\node[draw=plantteal,fill=plantteal!12,circle,minimum size=13mm,thick,font=\small] (mix) at (0,0) {Mixer};
\node[unit] (r1) at (2.1,0) {$1$\\$c_1$};
\node[unit] (ri) at (4.0,0) {$i$\\$c_i$};
\node[unit] (rn) at (5.9,0) {$n_r$\\$c_{n_r}$};
\node[font=\large,text=plantblue] at (3.05,0) {$\cdots$};
\node[font=\large,text=plantblue] at (4.95,0) {$\cdots$};
\coordinate (split) at (7.0,0);
\draw[draw=plantteal,fill=plantteal!8,rounded corners,thick] (8.5,-1.0) rectangle (10.8,1.0);
\node[font=\small,text=plantteal!70!black] at (9.65,1.28) {$L$-layer clarifier};
\draw[draw=plantteal!55] (8.5,.45)--(10.8,.45);
\draw[draw=plantteal!55] (8.5,-.4)--(10.8,-.4);
\node[font=\scriptsize] at (9.65,.72) {Layer $1$: $s_1=X_E$};
\node[font=\small,text=plantteal] at (9.65,.05) {$\vdots\quad s_\ell\quad\vdots$};
\node[font=\scriptsize] at (9.65,-.70) {Layer $L$: $s_L=X_U$};
\draw[flow] (-1.65,0)--(mix.west) node[midway,above,lab] {$Q_0,c_{\rm in}$};
\draw[flow] (mix.east)--(r1.west) node[midway,above,lab] {$q_PQ_0,m$};
\draw[flow] (r1.east)--(2.80,0);
\draw[flow] (3.30,0)--(ri.west);
\draw[flow] (ri.east)--(4.70,0);
\draw[flow] (5.20,0)--(rn.west);
\draw[thick,draw=plantblue] (rn.east)--(split);
\draw[flow] (split)--(8.5,0) node[midway,above,lab] {$q_CQ_0,c_{n_r}$};
\fill[plantpurple] (split) circle (1.5pt);
\draw[->,thick,draw=plantpurple] (split)--(7.0,2.1)--(0,2.1)--(mix.north)
node[pos=.58,above,lab,text=plantpurple] {Mixed-liquor recycle: $r_IQ_0,c_{n_r}$};
\draw[flow] (10.8,.72)--(12.75,.72) node[midway,above,lab] {$q_EQ_0,c_E$};
\node[font=\scriptsize,text=plantblue] at (11.60,.32) {Treated-water outlet};
\coordinate (underflow) at (9.65,-1.62);
\draw[thick,draw=plantorange] (9.65,-1.0)--(underflow) node[midway,right,lab] {$q_UQ_0,c_U$};
\fill[plantorange] (underflow) circle (1.5pt);
\draw[->,thick,draw=plantorange] (underflow)--(9.65,-2.7)
node[below,lab,text=plantorange] {Waste sludge: $wQ_0,c_U$};
\draw[->,thick,draw=plantorange] (underflow)--(0,-1.62)--(mix.south)
node[pos=.48,below,lab,text=plantorange] {Return sludge: $r_RQ_0,c_U$};
\end{tikzpicture}}
\caption{Connected plant boundary with $n_r$ biological reactors and $L$ clarifier layers. Blue arrows show the treatment path, purple the mixed-liquor recycle, and orange the clarifier underflow split into return and waste sludge. Intermediate reactors and layers are omitted for clarity. All flow ratios use fresh-influent flow as their basis. The operating case study uses five reactors and ten layers; its aeration duties are specified in Methodology.}
\label{fig:plant_configuration}
\end{figure}

The configuration also explains why the reactor-train and clarifier-feed flows differ. Mixed-liquor recycle is withdrawn before the clarifier, whereas return sludge originates from the clarifier underflow. The two recycle streams therefore carry concentrations from different locations. The mixer balance must use each stream's own flow and concentration rather than assigning a single recycle concentration to both.

The mixer balance is written component by component rather than in terms of an aggregate water-quality quantity. For component $f$, incoming and outgoing mass rates satisfy
\begin{equation}
\equationname{Scalar mixer equation for one component}
Q_Pm_f=Q_0c_{{\rm in},f}+r_IQ_0c_{n_r,f}+r_RQ_0c_{U,f}.
\end{equation}
When the component basis is g m$^{-3}$, each product has units of g d$^{-1}$. Dividing first by $Q_0$ and then by $q_P$ gives
\begin{equation}
\equationname{Mixer concentration as a flow-weighted component average}
m_f=\frac{c_{{\rm in},f}+r_Ic_{n_r,f}+r_Rc_{U,f}}{1+r_I+r_R}.
\end{equation}
For two components, the same calculation is performed independently in each row,
\begin{equation}
\equationname{Illustrative two-component mixer matrix equation}
\begin{bmatrix}q_P&0\\0&q_P\end{bmatrix}
\begin{bmatrix}m_1\\m_2\end{bmatrix}
=\begin{bmatrix}c_{{\rm in},1}+r_Ic_{n_r,1}+r_Rc_{U,1}\\
c_{{\rm in},2}+r_Ic_{n_r,2}+r_Rc_{U,2}\end{bmatrix}.
\end{equation}
The complete component vector therefore satisfies
\begin{equation}
\equationname{Vector mixer equation for the complete component state}
q_Pm=c_{\rm in}+r_Ic_{n_r}+r_Rc_U.
\label{eq:plant_mixer}
\end{equation}
A high mixer concentration can consequently arise from internal circulation rather than from creation of an external pollutant load. The component-resolved framework of \citet{Zhang2026} similarly treats recirculation and discharge streams explicitly.

A hypothetical example makes the flow weighting and recycle effect concrete. Let $Q_0=10{,}000$ m$^3$ d$^{-1}$, $r_I=2$, $r_R=0.5$, and $w=0.02$. The resulting physical flows $Q_P,Q_C,Q_U,Q_E$ are 35{,}000, 15{,}000, 5{,}200, and 9{,}800 m$^3$ d$^{-1}$, respectively. For one particulate component, let the fresh-influent, final-reactor, and underflow concentrations be 200, 1{,}000, and 2{,}000 g m$^{-3}$. The mixer calculation is
\begin{align}
\equationname{Illustrative component inventory entering the mixer}
q_Pm_f&=200+2(1{,}000)+0.5(2{,}000)=3{,}200,\\
\equationname{Illustrative mixer concentration after flow normalization}
m_f&=3{,}200/3.5=914.29\ \mathrm{g\,m^{-3}}.
\end{align}
The three incoming physical mass rates are 2{,}000, 20{,}000, and 10{,}000 kg d$^{-1}$, whose sum equals the 32{,}000 kg d$^{-1}$ leaving the mixer. The mixer concentration exceeds the fresh-influent concentration because both recycle streams return concentrated material. The example therefore illustrates how internal circulation changes concentration without creating external mass.

The nominal stage HRT is defined on the fresh-flow basis. Because volume divided by flow has units of time, $V_i/Q_0$ is measured in days and multiplication by 24 converts it to hours. The actual flow through each reactor is $Q_P=Q_0q_P$, so the residence time during one passage through reactor $i$ is $V_i/Q_P$ \citep{Henze2008,Alex2008}. The inverse of this residence time is the dilution rate. With reactor volume $V_i$,
\begin{equation}
\equationname{Stage retention times and recycle-adjusted dilution rates}
h_i=\frac{24V_i}{Q_0},\qquad H=\sum_{i=1}^{n_r}h_i,\qquad
d_i=\frac{Q_0q_P}{V_i}=\frac{24q_P}{h_i}.
\label{eq:plant_residence}
\end{equation}
The quantities $h_i$ and $H$ are in hours, while $d_i$ is in d$^{-1}$. The actual residence time during one passage through stage $i$ is $h_i/q_P$. For a hypothetical plant with $H=24$ h, $r_I=2$, and $r_R=0.5$, the reactor-train flow ratio is 3.5, so the train residence time per passage is approximately 6.86 h. With equal reactor volumes, each stage then has a residence time per passage of about 1.37 h. The example shows why nominal HRT and recycle flow must be interpreted together.

Fixed stage fractions can be represented by $h_i=f_i^V H$, where $f_i^V>0$ and $\sum_i f_i^V=1$. The case study uses equal fractions. Varying $H$ at fixed $Q_0$ therefore varies the total biological reactor volume. The optimization consequently combines capacity selection with operating decisions. In an existing plant with fixed tank volumes, these volumes would instead be parameters and $H$ would no longer appear in the free-control vector.

Let $\vartheta$ collect only the settings that vary freely. A constant setting remains a model parameter because treating it as a varying regression coordinate would create a zero-variance feature. A carbon or chemical addition may also enter the control vector. If an addition is represented only by a source term, its carrier flow is assumed negligible. If the carrier flow is appreciable, it must be represented as an additional hydraulic stream and the downstream dilution rates must be revised accordingly.

\subsubsection{Biological conversion and non-negativity}

The reaction network uses the component and process model developed in Chapter~\ref{ch:foundations}. Its stoichiometric matrix is $\nu\in\mathbb R^{P\times F}$ and its process-rate vector in reactor $i$ is $\rho_i\in\mathbb R_+^P$. For component $f$, the dynamic balance is
\begin{equation}
\equationname{Dynamic reactor-stage equation for one component}
V_i\frac{dc_{i,f}}{dt}=Q_Pc_{i-1,f}-Q_Pc_{i,f}
+V_i\sum_{p=1}^{P}\nu_{pf}\rho_{i,p}+V_i\omega_i(e_{S_O})_f+V_i\xi_{i,f}.
\end{equation}
The reaction sum combines the contributions of all modeled processes. A negative $\nu_{pf}$ consumes component $f$, while a positive entry produces it. The oxygen selector equals one for the dissolved-oxygen component and zero elsewhere. At steady state, accumulation vanishes. Dividing the remaining equation by $V_i$ gives
\begin{equation}
\equationname{Steady-state reactor-stage equation for one component}
0=d_i(c_{i-1,f}-c_{i,f})+\sum_{p=1}^{P}\nu_{pf}\rho_{i,p}
+\omega_i(e_{S_O})_f+\xi_{i,f}.
\end{equation}
Collecting all $F$ component balances gives the vector reactor equation, with $c_0=m$,
\begin{equation}
\equationname{Vector reactor-stage equation with aeration and external additions}
0=d_i(c_{i-1}-c_i)+\nu^\top\rho_i(c_i,\vartheta)
+k_La_i(a_i)(S_O^*-c_{i,S_O})e_{S_O}+\xi_i(\vartheta).
\label{eq:plant_reactor}
\end{equation}
The vector $e_{S_O}$ selects dissolved oxygen, $S_O^*$ is the oxygen saturation concentration, and $\xi_i$ represents declared stage additions. Every term has units of component concentration per day. The treatment-model framework of \citet{Henze2006} separates hydraulic transport, stoichiometric conversion, and external sources in this manner.

Non-negativity is a physical property of the dynamic state. At a boundary where one component concentration is zero, the combined reaction and external-source terms must not direct that component into negative values. With non-negative inflow and a locally Lipschitz balance function defined on the boundary, this condition preserves the non-negative state space of the dynamic model \citep{Haddad2010}. Numerical acceptance therefore checks concentrations and reaction rates independently of the solver's stopping message.

The same requirement does not arise automatically in a statistical surrogate. A fitted quadratic equation may predict a negative concentration even when every training state is non-negative because polynomial regression imposes no sign restriction on new outputs. The work of \citet{KircherVotsmeier2025} addresses simultaneous mass conservation and non-negativity in kinetic surrogates. Both properties therefore require explicit treatment in the connected plant surrogate.

\subsubsection{Stage invariants and the external plant balance}

The role of mass conservation can first be seen in a simple two-component reaction. Suppose one component is converted into another without changing their combined amount. The corresponding stoichiometric and inventory rows are
\begin{equation}
\equationname{Illustrative stoichiometric matrix and normalized invariant row}
\nu=\begin{bmatrix}-1&1\end{bmatrix},\qquad
A=\frac1{\sqrt2}\begin{bmatrix}1&1\end{bmatrix}.
\end{equation}
The factor $1/\sqrt2$ normalizes the inventory row to unit length. The reaction contribution disappears under multiplication by $A$,
\begin{equation}
\equationname{Illustrative cancellation of the reaction contribution to an invariant}
A\nu^\top=\frac1{\sqrt2}\left[1(-1)+1(1)\right]=0.
\end{equation}
At an illustrative dilution rate of 2 d$^{-1}$, let the upstream concentrations be $(12,3)$ and the downstream concentrations be $(8,7)$ g m$^{-3}$, with a reaction rate of 8 g m$^{-3}$ d$^{-1}$. The two component balances are
\begin{align}
\equationname{Illustrative steady-state equation for the consumed component}
0&=2(12-8)-8,\\
\equationname{Illustrative steady-state equation for the produced component}
0&=2(3-7)+8.
\end{align}
Adding the balances removes the unknown reaction rate. Both upstream and downstream totals are 15 g m$^{-3}$, and the normalized invariant is that total divided by $\sqrt2$. The full reactor model applies the same cancellation principle to a larger set of components and processes.

Let $A\in\mathbb R^{K\times F}$ contain independent inventory rows satisfying
\begin{equation}
\equationname{Reaction, oxygen-forcing, and rank identities of the plant invariant operator}
A\nu^\top=0,\qquad Ae_{S_O}=0,\qquad
\operatorname{rank}(A)=K.
\label{eq:plant_invariant_operator}
\end{equation}
The rows are normalized to unit Euclidean length, although this normalization does not make them mutually orthogonal. Physically named inventories and an orthonormal basis can span the same mass-conservation subspace while using different coordinates. The stoichiometric fixed-layer formulation of \citet{SturmWexler2022} illustrates the usefulness of separating conserved inventories from the reaction rates that redistribute them.

For inventory row $k$, multiply component equation $f$ by $A_{kf}$ and sum over $f$. The reaction contribution becomes
\begin{equation}
\equationname{Scalar cancellation of reaction rates in an invariant equation}
\sum_{f=1}^{F}A_{kf}\sum_{p=1}^{P}\nu_{pf}\rho_{i,p}
=\sum_{p=1}^{P}\left(\sum_{f=1}^{F}A_{kf}\nu_{pf}\right)\rho_{i,p}=0.
\end{equation}
The inner sum is entry $(k,p)$ of $A\nu^\top$. Oxygen transfer likewise disappears because $Ae_{S_O}=0$. The remaining scalar relation is
\begin{equation}
\equationname{Scalar reactor-stage invariant equation with external additions}
d_i\left(\sum_fA_{kf}c_{i-1,f}-\sum_fA_{kf}c_{i,f}\right)
+\sum_fA_{kf}\xi_{i,f}=0.
\end{equation}
Rearranging the transport difference and dividing by the positive $d_i$ gives
\begin{equation}
\equationname{Invariant inventory change across one reactor stage}
A(c_i-c_{i-1})=d_i^{-1}A\xi_i(\vartheta).
\label{eq:plant_invariant_transition}
\end{equation}
When no stage addition is present, $Ac_i=Ac_{i-1}$. Thus nitrification may transfer nitrogen from ammonium to oxidized forms, and denitrification may transfer it to dissolved nitrogen gas, while the complete nitrogen inventory remains unchanged. A reported quantity that excludes dissolved nitrogen gas can decrease even though the mass-conservation inventory remains constant. The nutrient-form interpretation in \citet{USEPA2009} therefore supports examining individual species as well as aggregate reporting quantities.

The same example also shows why external additions must appear in the invariant relation. If component 1 receives an addition of 2 g m$^{-3}$ d$^{-1}$ while the same reaction rate is retained, the downstream concentrations become $(9,7)$. The two balances are $2(12-9)-8+2=0$ and $2(3-7)+8=0$. Their total increases by $2/2=1$ g m$^{-3}$. An invariant relation that forced zero change would incorrectly remove the material introduced by the dose.

Internal recycle streams disappear only after the complete plant balance is assembled. First sum the reactor-stage transitions. Every intermediate $Ac_i$ then appears once with a positive sign and once with a negative sign, giving
\begin{align}
\equationname{Accumulated invariant change across the reactor train}
A(c_{n_r}-m)&=\sum_i\frac{V_i}{Q_0q_P}A\xi_i,\\
\equationname{Reactor-train invariant equation and external-addition total}
q_PAc_{n_r}-q_PAm&=D_{\rm add},\qquad
D_{\rm add}=\sum_i\frac{V_i}{Q_0}A\xi_i.
\end{align}
Next multiply the mixer balance by $A$ and substitute it for $q_PAm$,
\begin{align}
\equationname{Reactor-train invariant equation after mixer substitution}
q_PAc_{n_r}&=Ac_{\rm in}+r_IAc_{n_r}+r_RAc_U+D_{\rm add},\\
\equationname{Invariant equation after cancellation of mixed-liquor recycle}
(q_P-r_I)Ac_{n_r}&=Ac_{\rm in}+r_RAc_U+D_{\rm add}.
\end{align}
Because $q_P-r_I=q_C$, the left-hand side is the clarifier-feed inventory flow. The clarifier balance provides a second expression for the same quantity,
\begin{equation}
\equationname{Clarifier-feed invariant flow equals the two outlet invariant flows}
q_CAc_{n_r}=q_EAc_E+(r_R+w)Ac_U.
\end{equation}
Equating the two expressions and subtracting $r_RAc_U$ cancels the return-sludge flow. Mixed-liquor recycle has already canceled through $q_P-r_I$. The remaining external plant balance is
\begin{equation}
\equationname{External plant mass conservation after internal recycle cancellation}
Ac_{\rm in}+\sum_{i=1}^{n_r}\frac{V_i}{Q_0}A\xi_i
=q_EAc_E+wAc_U.
\label{eq:plant_external_invariant}
\end{equation}
Only fresh feed, stage additions, final effluent, and waste sludge remain. The recycle ratios continue to influence the internal state and treatment performance, but their cancellation from the external balance confirms that they are internal transfers rather than external sources or sinks.

For ASM2d-TSN, $F=20$, $P=28$, and $K=5$. The five inventories represent soluble inert material, total nitrogen including dissolved nitrogen gas, total phosphorus, an alkalinity-related combination, and a metal-related combination. The last two unnormalized combinations are
\begin{align}
\equationname{Alkalinity-related invariant combination}
I_{\rm alk}&=S_{ALK}-\frac{S_{NH4}}{14}
+\frac{S_{NO2}}{14}+\frac{S_{NO3}}{14}-\frac{S_{PO4}}{31},\\
\equationname{Metal-related invariant combination}
I_{\rm metal}&=3.45X_{MeP}+4.87X_{MeOH}.
\label{eq:plant_named_invariants}
\end{align}
The nitrogen and phosphorus rows use the constituent weights in Table~\ref{tab:component_basis_composites}, and the nitrogen inventory assigns unit weight to dissolved nitrogen gas. Under the stated parameterization, the stoichiometric rank is 15. These inventory identities are audited directly rather than inferred from reporting labels such as TN.

\subsection{Clarification and solids storage}

\subsubsection{A nonreactive separator with two outlet flows}

The clarifier is represented as a nonreactive settling unit. In the adopted formulation, the clarifier transports and separates components but does not perform biological or chemical conversion. Soluble components follow the liquid phase, whereas particulate components can settle and become concentrated in the underflow. \citet{Takacs1991} provide the settling and thickening formulation used for the aggregate solids profile, while \citet{JeppssonDiehl1996} show why propagation of component composition through a clarifier requires more information than aggregate suspended-solids concentration alone.

Selection means extracting designated entries of a vector without altering their values. For the illustrative component order $(S_A,S_{NH4},X_H,X_I)$, define
\begin{equation}
\equationname{Illustrative soluble and particulate component selectors}
E_S=\begin{bmatrix}1&0&0&0\\0&1&0&0\end{bmatrix},\qquad
E_X=\begin{bmatrix}0&0&1&0\\0&0&0&1\end{bmatrix}.
\end{equation}
Then $E_Sc=(S_A,S_{NH4})^\top$ and $E_Xc=(X_H,X_I)^\top$. Transposing a selector places the selected entries back into their original positions while filling all other positions with zeros. In this example,
\begin{equation}
\equationname{Illustrative coordinate projectors from the component selectors}
E_S^\top E_S=\operatorname{diag}(1,1,0,0),\qquad
E_X^\top E_X=\operatorname{diag}(0,0,1,1).
\end{equation}
Their sum is the identity because every component belongs to exactly one group, and their cross-product is zero because the groups do not overlap. For the complete ordered model, $E_S=[I_{10}\ 0]$ and $E_X=[0\ I_{10}]$. More generally, let $E_S\in\mathbb R^{F_S\times F}$ and $E_X\in\mathbb R^{F_X\times F}$ select soluble and particulate components. They satisfy
\begin{equation}
\equationname{Completeness and orthogonality identities of the component selectors}
E_S^\top E_S+E_X^\top E_X=I_F,\qquad E_SE_X^\top=0.
\end{equation}
The column $t_X$ maps component concentrations to TSS. Its entries are non-negative and its soluble entries are zero.

For the four-component order above, $t_X=(0,0,0.90,0.75)^\top$. If $X_H=100$ and $X_I=40$ g COD m$^{-3}$, the resulting TSS concentration is $0.90(100)+0.75(40)=120$ g TSS m$^{-3}$. The coefficients convert the native COD-based component concentrations to suspended-solids mass. Adding the numerical concentrations directly would therefore calculate a different quantity.

For one component, clarifier mass conservation gives
\begin{equation}
\equationname{Scalar mass conservation across the clarifier}
Q_Cc_{n_r,f}=Q_Ec_{E,f}+Q_Uc_{U,f}.
\end{equation}
The statistical response stores component flows divided by $Q_0$,
\begin{equation}
\equationname{Flow-weighted clarifier outlet component responses}
g_E=q_Ec_E,\qquad g_U=q_Uc_U.
\label{eq:plant_outlet_flows}
\end{equation}
Because $q_E$ and $q_U$ are dimensionless, these coordinates have the same units as concentration. Their corresponding physical discharge rates are $Q_0g_E$ and $Q_0g_U$. Representing the outlet responses in this way avoids division by a potentially small underflow ratio inside the equality system. The complete component balance and soluble pass-through relations are
\begin{equation}
\equationname{Joint clarifier component and soluble-transport equalities}
g_E+g_U=q_Cc_{n_r},\qquad
E_Sg_U=q_UE_Sc_{n_r}.
\label{eq:plant_clarifier_equalities}
\end{equation}
Together, these equations imply $E_Sg_E=q_EE_Sc_{n_r}$. Settling therefore cannot directly remove a soluble nutrient in the adopted clarifier model.

For soluble component $f$, pass-through gives $c_{U,f}=c_{n_r,f}$. Substituting $g_{U,f}=q_Uc_{n_r,f}$ into the total component balance yields
\begin{align}
\equationname{Flow-weighted soluble component in the overflow}
g_{E,f}&=(q_C-q_U)c_{n_r,f}=q_Ec_{n_r,f},\\
\equationname{Equality of overflow and feed soluble concentrations}
c_{E,f}&=g_{E,f}/q_E=c_{n_r,f}.
\end{align}
For a hypothetical feed concentration of 20 g m$^{-3}$ and flow ratios $q_C=1.5$, $q_U=0.52$, and $q_E=0.98$, the feed, underflow, and overflow component flows are 30, 10.4, and 19.6 g m$^{-3}$ on the fresh-flow basis. Dividing by the respective outlet flow ratios recovers 20 g m$^{-3}$ in both outlets.

Particulate components instead satisfy the densification requirement
\begin{equation}
\equationname{Particulate densification requirement for the underflow}
E_Xg_U\ge q_UE_Xc_{n_r}.
\label{eq:plant_densification}
\end{equation}
After division by $q_U$, every underflow particulate concentration is at least as large as its feed concentration. Together with non-negative overflow and the total component balance, this relation also prevents the underflow from recovering more particulate material than enters the clarifier.

For an illustrative particulate component, take $q_C=1.5$, $q_U=0.52$, and a feed concentration of 100 g m$^{-3}$. Densification requires $g_U\ge52$ g m$^{-3}$ on the fresh-flow basis. Mass conservation and non-negative overflow require $g_U\le150$ g m$^{-3}$. The fraction of incoming particulate flow recovered in the underflow therefore lies between $0.52/1.5$ and one. A soluble component with the same feed concentration instead has $g_U=52$ and equal feed and outlet concentrations. These two cases clarify the different roles of the soluble and particulate constraints.

Choose, for example, an admissible particulate flow $g_U=140$. Mass conservation then gives $g_E=150-140=10$. Dividing by the respective flow ratios gives $c_U=140/0.52=269.23$ and $c_E=10/0.98=10.20$ g m$^{-3}$. The underflow has therefore recovered $140/150=93.33\%$ of the incoming component flow. This recovery cannot be inferred from concentration alone because the outlet water flows differ.

For any particulate component with positive feed concentration,
\begin{equation}
\equationname{Bounds on the fraction of feed particulate material recovered in underflow}
\frac{q_U}{q_C}\le\frac{g_{U,f}}{q_Cc_{n_r,f}}\le1.
\end{equation}
Because the TSS weights are non-negative, multiplication by $t_X$ preserves these inequalities. Hence $X_U\ge X_N$. The solids balance $q_CX_N=q_EX_E+q_UX_U$ then gives
\begin{equation}
\equationname{Overflow solids bound implied by particulate densification}
q_EX_E\le(q_C-q_U)X_N=q_EX_N.
\end{equation}
Division by $q_E>0$ establishes the corresponding overflow bound.

Applying $t_X$ gives the feed and outlet suspended-solids concentrations,
\begin{equation}
\equationname{Feed, overflow, and underflow suspended-solids concentrations}
X_N=t_X^\top c_{n_r},\qquad
X_E=\frac{t_X^\top g_E}{q_E},\qquad
X_U=\frac{t_X^\top g_U}{q_U}.
\label{eq:plant_tss_endpoints}
\end{equation}
The declared constraints imply $X_E\le X_N\le X_U$. These relations bound the direction of separation, but they do not determine its magnitude. Additional information is therefore required from either the nonlinear settling calculation or an empirical endpoint model.

\subsubsection{The layer-resolved mechanistic model}

The mechanistic clarifier resolves the solids profile across $L$ layers, numbered from the surface to the hopper. Their TSS concentrations are $s_1,\ldots,s_L$, with $s_1=X_E$ and $s_L=X_U$. The exact settling velocity at concentration $s$ is
\begin{equation}
\equationname{Bounded double-exponential settling velocity}
v(s,X_N)=\max\left\{0,\min\left[v_{\max},
v_0\left(e^{-r_h(s-f_{\rm ns}X_N)}-e^{-r_p(s-f_{\rm ns}X_N)}\right)\right]\right\}.
\label{eq:plant_settling_velocity}
\end{equation}
The gravitational solids flux is $sv(s,X_N)$. The nonsettleable fraction $f_{\rm ns}$ makes the settling offset depend on feed solids concentration. At an interface, a receiver concentration above $X_t$ activates a limitation by the smaller adjacent gravitational flux. Otherwise, the upper-layer gravitational flux passes to the receiver \citep{Takacs1991}. Hydraulic transport uses the overflow velocity above the feed layer and the underflow velocity below it.

Each finite-volume layer equation is formed as incoming solids flux minus outgoing solids flux, multiplied by $A_{\rm cl}/V_{{\rm cl},\ell}$. The feed layer additionally receives $Q_0q_CX_N/V_{{\rm cl},\ell_F}$. Internal interface fluxes cancel when the layer equations are summed. At steady state, the resulting total solids balance reduces to Equation~\eqref{eq:plant_clarifier_equalities} after application of $t_X$.

A three-layer example makes this cancellation explicit. Place the feed in layer 2, and let $B_{12}$ denote the net solids rate from layer 2 toward layer 1 and $B_{23}$ the net rate from layer 2 toward layer 3. These signed rates combine hydraulic and gravitational transport in g d$^{-1}$. The layer balances are
\begin{align}
\equationname{Illustrative solids equation for the upper clarifier layer}
V_1\dot s_1&=B_{12}-Q_Es_1,\\
\equationname{Illustrative solids equation for the clarifier feed layer}
V_2\dot s_2&=Q_CX_N-B_{12}-B_{23},\\
\equationname{Illustrative solids equation for the lower clarifier layer}
V_3\dot s_3&=B_{23}-Q_Us_3.
\end{align}
Adding the three equations eliminates both interface rates. At steady state, the remaining relation is $Q_CX_N=Q_EX_E+Q_UX_U$. The whole-clarifier balance therefore holds even though the internal interface fluxes depend nonlinearly on neighboring concentrations. Mass conservation alone cannot determine the internal fluxes or the shape of the layer profile.

At a mechanistic steady state, particulate outlet components are reconstructed using the final-reactor particulate composition,
\begin{equation}
\equationname{Particulate outlet reconstruction from endpoint solids concentrations}
E_Xc_E=\frac{s_1}{X_N}E_Xc_{n_r},\qquad
E_Xc_U=\frac{s_L}{X_N}E_Xc_{n_r}.
\label{eq:plant_outlet_reconstruction}
\end{equation}
Soluble outlet components remain equal to their feed concentrations. The mapping is defined for $X_N>0$ and assumes that the particulate composition leaving through the overflow and underflow retains the final-reactor composition \citep{JeppssonDiehl1996}.

\subsubsection{An exact interval for aggregate inventory under the endpoint envelope}

The clarifier inventory is the concentration in each layer multiplied by its volume and summed over the vessel. For an illustrative four-layer clarifier,
\begin{equation}
\equationname{Illustrative four-layer clarifier solids inventory}
M=V_1X_E+V_2s_2+V_3s_3+V_4X_U.
\end{equation}
If each layer has volume 100 m$^3$ and the concentrations are $(10,200,600,1{,}000)$ g m$^{-3}$, the total stored solids are $100(10+200+600+1{,}000)=181{,}000$ g, or 181 kg. For the general clarifier,
\begin{equation}
\equationname{General clarifier solids inventory and total volume}
M_{\rm cl}=\sum_{\ell=1}^{L}V_{{\rm cl},\ell}s_\ell,\qquad
V_{\rm cl}=\sum_{\ell=1}^{L}V_{{\rm cl},\ell}.
\label{eq:plant_inventory_definition}
\end{equation}
The outlet concentrations alone cannot determine this inventory because two clarifiers can share the same endpoints while storing different amounts internally. Whole-plant solids retention time therefore requires both the internal inventory and the solids leaving through the outlet streams. The benchmark formulation of \citet{Alex2008} likewise includes reactor and settler inventories in its sludge-age calculation.

The mechanistic acceptance envelope requires
\begin{equation}
\equationname{Interior-layer solids envelope between the outlet concentrations}
X_E\le s_\ell\le X_U,\qquad \ell=2,\ldots,L-1.
\label{eq:plant_layer_envelope}
\end{equation}
This condition constrains each interior layer between the two endpoint concentrations, but it does not require monotonic variation from one layer to the next. With positive layer volumes, the minimum inventory is obtained by assigning every interior layer $X_E$, while the maximum is obtained by assigning every interior layer $X_U$.

For four layers, the two unknown interior contributions satisfy
\begin{align}
\equationname{Lower inventory bound for two illustrative interior layers}
V_2X_E+V_3X_E&\le V_2s_2+V_3s_3\\
\equationname{Upper inventory bound for two illustrative interior layers}
&\le V_2X_U+V_3X_U.
\end{align}
Adding the fixed endpoint contributions gives
\begin{align}
\equationname{Lower total-inventory bound for an illustrative four-layer clarifier}
(V_1+V_2+V_3)X_E+V_4X_U&\le M,\\
\equationname{Upper total-inventory bound for an illustrative four-layer clarifier}
M&\le V_1X_E+(V_2+V_3+V_4)X_U.
\end{align}
For the hypothetical equal-volume clarifier above, these limits are 103{,}000 and 301{,}000 g, and the 181{,}000-g inventory lies within them. If the two internal concentrations were both 400 g m$^{-3}$, the total inventory would again be 181{,}000 g. Matching the endpoints and total inventory therefore does not uniquely determine the internal profile.

Figure~\ref{fig:plant_inventory_envelope} shows the same four-layer arithmetic as concentration bars. The surface and hopper concentrations remain fixed while the two internal layers change, altering the stored mass. The left and right panels attain the lower and upper inventory limits implied by the envelope.

\begin{figure}[htbp]
\centering
\includegraphics[width=\linewidth]{ch6_concept_inventory_envelope.pdf}
\caption{Hypothetical four-layer profiles with identical overflow and underflow concentrations. Each layer has volume 100 m$^3$. The lower-bound, illustrative interior, and upper-bound profiles store 103, 181, and 301 kg of solids. Hatching identifies the fixed endpoints, while dotted fills identify internal layers. The diagram represents the adopted endpoint envelope.}
\label{fig:plant_inventory_envelope}
\end{figure}

The middle profile is only one of many profiles that can share the same endpoints and total inventory. The scalar inventory retains the amount of stored solids, whereas the nonlinear layer equations provide the additional information needed to determine a mechanistic distribution.

For $L$ layers, the same argument replaces $V_2+V_3$ with $\sum_{\ell=2}^{L-1}V_{{\rm cl},\ell}$. Because this internal-volume sum is $V_{\rm cl}-V_{{\rm cl},1}-V_{{\rm cl},L}$, collecting endpoint coefficients gives
\begin{align}
\equationname{Lower clarifier-inventory bound from endpoint concentrations}
M_{\rm cl}^{-}&=(V_{\rm cl}-V_{{\rm cl},L})X_E+V_{{\rm cl},L}X_U,\\
\equationname{Upper clarifier-inventory bound from endpoint concentrations}
M_{\rm cl}^{+}&=V_{{\rm cl},1}X_E+(V_{\rm cl}-V_{{\rm cl},1})X_U,\\
\equationname{Admissible interval for the clarifier solids inventory}
M_{\rm cl}^{-}&\le M_{\rm cl}\le M_{\rm cl}^{+}.
\label{eq:plant_inventory_bounds}
\end{align}

These bounds are both necessary and sufficient for the existence of an endpoint-envelope profile with the stated total inventory. To see sufficiency, assign every interior layer a common concentration $s_*$ between $X_E$ and $X_U$. The total inventory becomes
\begin{equation}
\equationname{Inventory represented by a common interior-layer concentration}
V_{{\rm cl},1}X_E+V_{{\rm cl},L}X_U
+\left(V_{\rm cl}-V_{{\rm cl},1}-V_{{\rm cl},L}\right)s_*.
\end{equation}
As $s_*$ varies continuously between $X_E$ and $X_U$, this expression spans the entire admissible interval. For equal-volume layers, it also gives a simple constructive profile for every inventory admitted by the bounds.

As an illustrative calculation, consider ten layers of 600 m$^3$, with $X_E=10$ g m$^{-3}$ and $X_U=5{,}000$ g m$^{-3}$. The corresponding inventory interval is 3{,}054{,}000--27{,}006{,}000 g. Every inventory in this interval is compatible with at least one profile satisfying the endpoint envelope. Joint projection enforces this aggregate interval, whereas layer-resolved mechanistic evaluation checks the internal settling equations.

For the ten equal layers, the arithmetic is
\begin{align}
\equationname{Solids inventory of the illustrative ten-layer clarifier}
M_{\rm cl}&=600(X_E+s_2+s_3+\cdots+s_9+X_U),\\
\equationname{Lower inventory bound for the illustrative ten-layer clarifier}
M_{\rm cl}^{-}&=600(9X_E+X_U),\\
\equationname{Upper inventory bound for the illustrative ten-layer clarifier}
M_{\rm cl}^{+}&=600(X_E+9X_U).
\end{align}
At the lower endpoint, nine layers are at $X_E$ and one is at $X_U$; at the upper endpoint, one layer is at $X_E$ and nine are at $X_U$. This makes the coefficients in the compact interval physically transparent rather than treating them as fitted constants.

\subsection{Statistical prediction of the connected response}

\subsubsection{A response that preserves network and storage information}

The connected statistical response must retain enough information to reconstruct the physical relationships used in treatment and operating analysis. For a teaching example with one reactor and two components, let $S$ denote a soluble component and $X$ a particulate component. The stored response is the explicit nine-entry vector
\begin{equation}
\equationname{Illustrative joint response of a one-stage, two-component plant}
\chi_{\rm toy}=
(m_S,m_X,c_{1,S},c_{1,X},g_{E,S},g_{E,X},g_{U,S},g_{U,X},M_{\rm cl})^\top.
\end{equation}
Adding a second reactor inserts $c_{2,S}$ and $c_{2,X}$ before the outlet coordinates, increasing the vector length to eleven. In general, every reactor contributes $F$ component concentrations. The mixer and two clarifier outlets each contribute another $F$, and the clarifier inventory contributes one scalar. The connected surrogate therefore predicts
\begin{equation}
\equationname{Complete joint plant-response vector and its dimension}
\chi=\operatorname{col}(m,c_1,\ldots,c_{n_r},g_E,g_U,M_{\rm cl})
\in\mathbb R^{n_\chi},\qquad n_\chi=(n_r+3)F+1.
\label{eq:plant_joint_response}
\end{equation}
Here $\operatorname{col}$ stacks its arguments into one column. The response retains concentrations needed for reactor inventories, outlet flows needed for material balances, and clarifier storage needed for solids retention time. It intentionally omits the internal clarifier-layer profile. The full mechanistic state is
\begin{equation}
\equationname{Mechanistic reactor-and-clarifier state vector and its dimension}
y=\operatorname{col}(c_1,\ldots,c_{n_r},s_1,\ldots,s_L)
\in\mathbb R_+^{n_y},\qquad n_y=n_rF+L.
\label{eq:plant_mechanistic_state}
\end{equation}
The mapping from $y$ to $\chi$ deterministically reconstructs the mixer, both outlet streams, and the clarifier inventory. The reduced statistical response and the full mechanistic state therefore have different dimensions while still supporting the same engineering calculations.

For the five-reactor case, the connected response has $20+5(20)+20+20+1=161$ coordinates, whereas the full mechanistic state has $5(20)+10=110$. Table~\ref{tab:plant_response_positions} records the response ordering explicitly. A position identifies one coordinate and does not imply that neighboring coordinates have the same physical unit.

\begin{table}[htbp]
\centering\small
\caption{Blocks in the 161-coordinate connected response.}
\label{tab:plant_response_positions}
\begin{tabular}{lll}
\toprule Block & Positions & Coordinates\\\midrule
Mixer & 1--20 & $m_1,\ldots,m_{20}$\\
Reactor 1 & 21--40 & $c_{1,1},\ldots,c_{1,20}$\\
Reactor 2 & 41--60 & $c_{2,1},\ldots,c_{2,20}$\\
Reactor 3 & 61--80 & $c_{3,1},\ldots,c_{3,20}$\\
Reactor 4 & 81--100 & $c_{4,1},\ldots,c_{4,20}$\\
Reactor 5 & 101--120 & $c_{5,1},\ldots,c_{5,20}$\\
Overflow flows & 121--140 & $g_{E,1},\ldots,g_{E,20}$\\
Underflow flows & 141--160 & $g_{U,1},\ldots,g_{U,20}$\\
Clarifier inventory & 161 & $M_{\rm cl}$\\
\bottomrule
\end{tabular}
\end{table}

\subsubsection{Complete second-order features and ridge estimation}

The statistical input contains every free operating control together with every fresh-influent component concentration. Standardization first subtracts a fitting-set mean and divides by a fitting-set standard deviation. For one operating coordinate and one influent coordinate,
\begin{equation}
\equationname{Illustrative standardization of a control and an influent concentration}
z_1=\frac{\vartheta_1-\bar\vartheta_1}{s_{\vartheta_1}},\qquad
z_2=\frac{c_{{\rm in},1}-\bar c_{{\rm in},1}}{s_{{\rm in},1}}.
\end{equation}
The resulting coordinates are dimensionless. For hypothetical aeration $0.8$, fitting mean $0.5$, and fitting scale $0.1$, the standardized value is 3. For hypothetical ammonium concentration 40 g m$^{-3}$ with mean 30 and scale 10 g m$^{-3}$, the standardized value is 1. Standardization therefore prevents the regression penalty from being determined directly by the native units of the inputs. With fitting-set centers and positive diagonal scales,
\begin{equation}
\equationname{Standardized plant input vector and its dimension}
z=\operatorname{col}\left[D_{\vartheta,\rm fit}^{-1}(\vartheta-\bar\vartheta),
D_{{\rm in},\rm fit}^{-1}(c_{\rm in}-\bar c_{\rm in})\right]
\in\mathbb R^{d_z},\qquad d_z=d_\vartheta+F.
\label{eq:plant_standardized_input}
\end{equation}

For two standardized inputs, their symmetric product matrix is
\begin{equation}
\equationname{Illustrative quadratic product matrix of two standardized inputs}
zz^\top=\begin{bmatrix}z_1^2&z_1z_2\\z_2z_1&z_2^2\end{bmatrix}.
\end{equation}
Because multiplication commutes, $z_1z_2$ and $z_2z_1$ represent the same monomial. Keeping one triangular copy therefore yields the unique quadratic terms $(z_1^2,z_1z_2,z_2^2)$. Together with the intercept and two linear terms, this gives
\begin{equation}
\equationname{Illustrative unique-monomial second-order feature vector}
(1,z_1,z_2,z_1^2,z_1z_2,z_2^2)^\top.
\end{equation}
With three inputs, the feature vector contains an intercept, three linear terms, three squares, and three distinct cross-products, for ten features. In general, there are $d_z$ squares and $d_z(d_z-1)/2$ unique cross-products, giving
\begin{align}
\equationname{Linear and unique-quadratic feature coordinates}
r(z)&=\operatorname{col}\left[z,\operatorname{vech}(zz^\top)\right],\\
\equationname{Centered and scaled plant regression feature vector}
\phi(\vartheta,c_{\rm in})&=\operatorname{col}\left[1,D_{r,\rm fit}^{-1}(r(z)-\bar r)\right],\\
\equationname{Dimension of the unique-monomial plant feature vector}
n_\phi&=1+d_z+\frac{d_z(d_z+1)}{2}.
\label{eq:plant_features}
\end{align}
The operator $\operatorname{vech}$ keeps one triangular copy of the symmetric matrix, so every square and pairwise product appears once. A product between ammonium loading and aeration allows the aeration response to depend on influent loading, while a squared recycle term permits local curvature \citep{Nazlabadi2021,Selerio2026}. These features approximate the nonlinear input-response relation.

Quadratic features are themselves centered and scaled after construction. For the hypothetical standardized inputs $(3,1)$, the unscaled nonconstant terms are $(3,1,9,3,1)$. If the first squared feature has fitting mean 1 and standard deviation 2, its final value is $(9-1)/2=4$. If the cross-product has fitting mean 0.2 and standard deviation 0.5, its final value is $(3-0.2)/0.5=5.6$. This second stage of standardization prevents a feature with naturally large variance from receiving an effectively weaker ridge penalty simply because of its numerical scale \citep{Hastie2009}.

A monomial is one polynomial term, such as $z_1^2$ or $z_1z_2$. A feature set containing each distinct monomial once spans the same quadratic function family as an ordered set that includes both $z_1z_2$ and $z_2z_1$. The two duplicate coefficients enter prediction only through their sum. If the effective cross-product coefficient is $b$, the ordered representation can assign $b/2$ to each duplicate, producing a squared penalty of $b^2/2$ rather than $b^2$. Removing duplicate monomials therefore preserves the polynomial family but changes the regularization convention unless the penalty is adjusted. The connected fit selects its penalty for the stated feature family containing each distinct term once. With 27 inputs, the resulting feature count is $1+27+378=406$.

For response coordinate $k$, let its standardized target be $\widetilde\chi_{ik}=(\chi_{ik}-\bar\chi_k)/s_{\chi,k}$. Its fitted value is $C_{k0}+\sum_{j=1}^{n_\phi-1}C_{kj}\phi_{ij}$. The multiresponse ridge objective is
\begin{equation}
\equationname{Scalar expansion of the multiresponse ridge objective}
\frac1{n n_\chi}\sum_{i=1}^{n}\sum_{k=1}^{n_\chi}
\left(\widetilde\chi_{ik}-\sum_{j=0}^{n_\phi-1}C_{kj}\phi_{ij}\right)^2
+\frac{\gamma}{n_\chi}\sum_{k=1}^{n_\chi}\sum_{j=1}^{n_\phi-1}C_{kj}^2.
\end{equation}
The intercept corresponds to $j=0$. It appears in prediction but is excluded from the penalty. Dividing a physical prediction error by $s_{\chi,k}$ is equivalent to weighting its squared error by $1/s_{\chi,k}^2$. Response scaling therefore gives coordinates of different physical magnitudes comparable influence in the joint statistical score.

For example, an illustrative concentration error of 10 g m$^{-3}$ with response scale 20 g m$^{-3}$ contributes $0.5^2=0.25$ before averaging, whereas an inventory error of 10{,}000 g with scale 200{,}000 g contributes $0.05^2=0.0025$. Comparing the unscaled errors alone would suggest the opposite ordering. The fitted scales therefore define what constitutes a large error in the joint regression.

Let $\Phi_{\mathcal I}$ contain the feature rows for fitting indices $\mathcal I$, and let $\widetilde Y_{\mathcal I}$ contain the standardized joint responses. With $R_\phi=\operatorname{diag}(0,1,\ldots,1)$, the coefficient matrix solves
\begin{equation}
\equationname{Matrix form of the multiresponse ridge fitting problem}
\widehat C_\gamma=\arg\min_C
\left\{\frac{\|\widetilde Y_{\mathcal I}-\Phi_{\mathcal I}C^\top\|_F^2}
{|\mathcal I|n_\chi}+\frac{\gamma}{n_\chi}\|CR_\phi\|_F^2\right\}.
\label{eq:plant_ridge_fit}
\end{equation}
The intercept remains unpenalized. Positive $\gamma$ stabilizes the fit in the presence of correlated quadratic features. \citet{Hastie2009} describe this trade-off between flexibility and coefficient shrinkage. The raw prediction on the physical response scale is
\begin{equation}
\equationname{Raw joint plant prediction on the physical response scale}
\chi_{\rm raw}=\bar\chi+D_\chi\widehat C_\gamma\phi(\vartheta,c_{\rm in}).
\label{eq:plant_raw_prediction}
\end{equation}
The fixed response scale $D_\chi$ places inventory, concentration, and outlet-flow coordinates on comparable statistical scales. The joint projection described below then adjusts this raw response to satisfy the declared physical requirements.

For one response coordinate with coefficient vector $\beta$, differentiation of the scalar ridge objective gives the normal equations
\begin{equation}
\equationname{Normal equations for one ridge response coordinate}
(\Phi^\top\Phi+n\gamma R_\phi)\widehat\beta=\Phi^\top\widetilde\chi.
\end{equation}
The term $n\gamma R_\phi$ adds positive curvature to the nonintercept coefficient directions. A simple one-feature example illustrates the resulting shrinkage. Let three feature and target values both be $(-\sqrt{3/2},0,\sqrt{3/2})$, giving zero means and unit population standard deviations. Symmetry gives a zero intercept. With slope $b$, the objective becomes $(1-b)^2+\gamma b^2$. Differentiation gives $2(b-1)+2\gamma b=0$, so $\widehat b=1/(1+\gamma)$. For $\gamma=1$, the fitted slope is 0.5 rather than the unpenalized value of one.

Finally, the physical prediction is recovered coordinate by coordinate as $\chi_{{\rm raw},k}=\bar\chi_k+s_{\chi,k}\sum_j\widehat C_{kj}\phi_j$. A hypothetical mean of 100 g m$^{-3}$, scale of 20 g m$^{-3}$, and standardized prediction of 0.5 therefore return 110 g m$^{-3}$. Conversion back to physical units is required before component balances can be imposed.

The quadratic feature family is retained from ICSOR \citep{Selerio2026}. Its coefficients describe the raw statistical mapping. When projection is active, however, a raw coefficient is not generally the derivative of the final projected response because the projected derivative also depends on the constraints, their dependence on the controls, and the active inequality set.

\subsubsection{A separate positive model for overflow solids}

Overflow TSS occupies a much smaller concentration range than reactor solids, so its prediction requires particular care. An absolute error of 1 g m$^{-3}$ is only 1\% of 100 g m$^{-3}$ but 100\% of 1 g m$^{-3}$. A standardized joint-response score may therefore describe reactor and outlet solids adequately in aggregate while still obscuring relative error at the discharge point. The relative-error argument of \citet{Tofallis2015} motivates a separate logarithmic model for overflow solids.

For positive mechanistic targets, define
\begin{equation}
\equationname{Dimensionless logarithmic overflow-solids target}
\ell_E=\log(X_E/X_\star),\qquad X_\star=1\ \mathrm{mg\,L^{-1}}.
\label{eq:plant_log_target}
\end{equation}
The ratio inside the logarithm is dimensionless. The scalar ridge fit uses the same complete quadratic feature map,
\begin{equation}
\equationname{Ridge fitting problem for the logarithmic overflow model}
\widehat\beta_{E,\gamma_E}=\arg\min_\beta
\left\{\frac{\|\ell_{E,\mathcal I}-\Phi_{\mathcal I}\beta\|_2^2}{|\mathcal I|}
+\gamma_E\|R_\phi\beta\|_2^2\right\}.
\end{equation}
Its back-transformed prediction is
\begin{equation}
\equationname{Positive overflow-solids prediction from logarithmic back-transformation}
\widehat X_{E,\log}=X_\star\exp\left(\widehat\beta_{E,\gamma_E}^\top\phi\right)>0.
\label{eq:plant_log_prediction}
\end{equation}
A twofold concentration error therefore produces the same logarithmic discrepancy at low and high concentrations. Exponentiation also guarantees a positive statistical endpoint before the joint mass-conservation projection is solved.

A small example clarifies the interpretation of the logarithmic fit. For hypothetical targets of 2 and 8 mg L$^{-1}$, the dimensionless logarithms are $\log2$ and $\log8$. If one fitted log value represents them equally under squared logarithmic error, the minimizer is
\begin{equation}
\equationname{Illustrative logarithmic average and geometric concentration mean}
\frac{\log2+\log8}{2}=\log4.
\end{equation}
Back-transformation gives 4 mg L$^{-1}$, the geometric rather than arithmetic mean. A fitted log value of $-2$ gives 0.1353 mg L$^{-1}$, which is still positive even though it lies below $X_\star$. The reference concentration supplies units; it does not impose a lower concentration bound.

Under squared logarithmic error, direct exponentiation represents a conditional geometric center rather than an unbiased concentration-scale mean \citep{Tofallis2015}. No concentration-scale bias adjustment is therefore applied. The separate model also avoids imposing a constant solids-removal fraction across all influents and controls, which would ignore variation in feed composition, hydraulics, recycle, and biological state. Instead, the fitted endpoint supplies the empirical closure introduced in Chapter~\ref{ch:literature}. The physical balances determine the total solids entering and leaving the clarifier, but they do not determine how those solids divide between overflow and underflow. The fitted overflow solids concentration supplies that missing relation. It remains an empirical estimate rather than a mass-conservation identity.

\subsection{Joint projection onto the declared plant constraints}

\subsubsection{Equalities and inequalities at fixed controls}

Joint projection adjusts all connected response blocks simultaneously. The structure of the constraint matrix is easiest to see in a hypothetical one-reactor, two-component plant using the nine-coordinate response $\chi_{\rm toy}$. Let $r_I=1$, $r_R=0.5$, and $w=0.1$, so $q_P=2.5$, $q_C=1.5$, $q_U=0.6$, and $q_E=0.9$. Let $S$ be soluble and $X$ particulate, with unit TSS weight, and suppose the sum of the two components is fixed by reactor mass conservation. If the fresh concentrations are $(30,20)$ g m$^{-3}$, the mixer, invariant, clarifier, soluble-pass-through, and empirical-endpoint equalities are
\begin{align}
\equationname{Illustrative mixer equality for the soluble component}
2.5m_S-c_{1,S}-\tfrac56g_{U,S}&=30,\\
\equationname{Illustrative mixer equality for the particulate component}
2.5m_X-c_{1,X}-\tfrac56g_{U,X}&=20,\\
\equationname{Illustrative reactor mass conservation equality}
c_{1,S}+c_{1,X}-m_S-m_X&=0,\\
\equationname{Illustrative clarifier equality for the soluble component}
g_{E,S}+g_{U,S}-1.5c_{1,S}&=0,\\
\equationname{Illustrative clarifier equality for the particulate component}
g_{E,X}+g_{U,X}-1.5c_{1,X}&=0,\\
\equationname{Illustrative soluble underflow transport equality}
g_{U,S}-0.6c_{1,S}&=0,\\
\equationname{Illustrative empirical overflow-solids equality}
g_{E,X}&=0.9(10)=9.
\end{align}
The final row corresponds to a fitted overflow concentration of 10 g m$^{-3}$. The invariant row is written without its unit-length normalization factor because multiplying a zero equality by a positive constant does not change its feasible set. In the chosen response order, the complete equality system is
\begin{equation}
\equationname{Matrix form of the illustrative plant projection equalities}
\left[\begin{array}{rrrrrrrrr}
2.5&0&-1&0&0&0&-5/6&0&0\\
0&2.5&0&-1&0&0&0&-5/6&0\\
-1&-1&1&1&0&0&0&0&0\\
0&0&-1.5&0&1&0&1&0&0\\
0&0&0&-1.5&0&1&0&1&0\\
0&0&-0.6&0&0&0&1&0&0\\
0&0&0&0&0&1&0&0&0
\end{array}\right]\chi_{\rm toy}
=\begin{bmatrix}30\\20\\0\\0\\0\\0\\9\end{bmatrix}.
\end{equation}
Each row corresponds to one previously derived physical or empirical relation. A zero entry simply means that the corresponding response coordinate does not appear in that equation. The inventory coordinate has a zero column because these equalities do not determine stored clarifier solids.

The illustrative response
\begin{equation}
\equationname{Illustrative joint response satisfying the plant projection equalities}
\chi_{\rm toy}=(30,50,30,50,27,9,18,66,18{,}000)^\top
\end{equation}
satisfies all seven rows. For example, the particulate mixer balance is $2.5(50)-50-(5/6)(66)=20$, while the clarifier particulate balance is $9+66=1.5(50)$. The underflow concentration is $66/0.6=110$ g m$^{-3}$. These calculations make the linear constraint assembly directly auditable.

For the full plant, the equality operator $\mathcal H(\vartheta)\chi=b(\vartheta,c_{\rm in})$ collects the same row types,
\begin{equation}
\equationname{Joint mixer, reactor-invariant, and clarifier projection equalities}
\operatorname{col}\left[
\begin{array}{c}
q_Pm-r_Ic_{n_r}-(r_R/q_U)g_U-c_{\rm in}\\
\{A(c_i-c_{i-1})-d_i^{-1}A\xi_i\}_{i=1}^{n_r}\\
g_E+g_U-q_Cc_{n_r}\\
E_Sg_U-q_UE_Sc_{n_r}
\end{array}\right]=0.
\label{eq:plant_projection_equalities}
\end{equation}
The empirical overflow relation is appended as
\begin{equation}
\equationname{Empirical overflow-solids closure equality}
t_X^\top g_E=q_E\widehat X_{E,\log}(\vartheta,c_{\rm in}).
\label{eq:plant_overflow_closure}
\end{equation}
The complete operator is denoted $\mathcal H_{\log}$ with right-hand side $b_{\log}$. It has $2F+n_rK+F_S+1$ rows. The final row fixes total overflow TSS but does not prescribe the particulate composition responsible for that total.

The remaining teaching-plant requirements appear as inequalities. With three 100-m$^3$ clarifier layers, $V_{\rm cl}=300$ m$^3$ and $V_1=V_3=100$ m$^3$. Particulate densification requires $0.6c_{1,X}-g_{U,X}\le0$. Substituting $X_E=g_{E,X}/0.9$ and $X_U=g_{U,X}/0.6$ into the inventory interval, then multiplying by the positive product $0.9(0.6)=0.54$, gives
\begin{align}
\equationname{Illustrative lower inventory bound in outlet-flow coordinates}
120g_{E,X}+90g_{U,X}-0.54M_{\rm cl}&\le0,\\
\equationname{Illustrative upper inventory bound in outlet-flow coordinates}
-60g_{E,X}-180g_{U,X}+0.54M_{\rm cl}&\le0.
\end{align}
The coefficient $120=0.6(300-100)$ and the coefficient $90=0.9(100)$ arise from hydraulic and volume factors, not from statistical fitting. At the illustrative response above, densification gives $30-66=-36$, and both inventory rows give $-2{,}700$. The implied inventory interval is 13{,}000--23{,}000 g, which contains the selected value of 18{,}000 g.

In the same response ordering, the three inequalities are
\begin{equation}
\equationname{Matrix form of illustrative densification and inventory inequalities}
\left[\begin{array}{rrrrrrrrr}
0&0&0&0.6&0&0&0&-1&0\\
0&0&0&0&0&120&0&90&-0.54\\
0&0&0&0&0&-60&0&-180&0.54
\end{array}\right]\chi_{\rm toy}\le0.
\end{equation}
Non-negativity would add one row $-\chi_{{\rm toy},k}\le0$ for each response coordinate. Because the equality and inequality matrices act on the same vector, their coordinate ordering must remain identical; otherwise, the wrong physical quantities would be constrained.

For a plant of arbitrary size, the same row patterns give
\begin{equation}
\equationname{General densification and clarifier-inventory projection inequalities}
\mathcal G(\vartheta)\chi=
\operatorname{col}\left[
\begin{array}{c}
q_UE_Xc_{n_r}-E_Xg_U\\
q_U(V_{\rm cl}-V_{{\rm cl},L})t_X^\top g_E
+q_EV_{{\rm cl},L}t_X^\top g_U-q_Eq_UM_{\rm cl}\\
q_Eq_UM_{\rm cl}-q_UV_{{\rm cl},1}t_X^\top g_E
-q_E(V_{\rm cl}-V_{{\rm cl},1})t_X^\top g_U
\end{array}\right]\le0.
\label{eq:plant_projection_inequalities}
\end{equation}
The first block enforces particulate densification. The final two rows represent the lower and upper clarifier-inventory bounds after multiplication by the positive product $q_Eq_U$, which removes divisions without changing feasibility. Non-negativity is imposed separately through $\chi\ge0$.

At fixed $(\vartheta,c_{\rm in})$, every constraint is affine in $\chi$. This remains true for the empirical overflow closure because its fitted endpoint becomes a known scalar once the inputs are fixed. The inventory inequalities are likewise affine in the response even though their coefficients vary with the controls. Distinguishing fixed-control linearity from control dependence is important for the numerical formulation and the sensitivity analysis that follows.

\subsubsection{The scaled nearest-response problem}

Projection is defined in standardized response coordinates. For response coordinate $k$, let $\chi_{{\rm raw},k}$ be the raw prediction and $s_k>0$ its fixed scale. The standardized displacement is
\begin{equation}
\equationname{Standardized projection displacement of one response coordinate}
u_k=\frac{\chi_k-\chi_{{\rm raw},k}}{s_k}.
\end{equation}
The projection objective is $\tfrac12\sum_{k=1}^{n_\chi}u_k^2$. A movement of two scale units in one coordinate therefore contributes $\tfrac12(2)^2=2$, while a movement of half a scale unit contributes $\tfrac12(0.5)^2=0.125$. With $D_\chi=\operatorname{diag}(s_1,\ldots,s_{n_\chi})$, the physical response is reconstructed as $\chi=\chi_{\rm raw}+D_\chi u$.

For a mixer component, substitution gives
\begin{align}
\equationname{Scaled mixer equation before the return-sludge term}
q_P(m_{{\rm raw},f}+s_{m,f}u_{m,f})
&-r_I(c_{{\rm raw},n_r,f}+s_{n_r,f}u_{n_r,f})\\
\equationname{Return-sludge term and target of the scaled mixer equation}
&-\frac{r_R}{q_U}(g_{{\rm raw},U,f}+s_{U,f}u_{U,f})
=c_{{\rm in},f}.
\end{align}
The raw terms are known at the current inputs, while the scale factors multiply the unknown standardized displacements. Moving the known quantities to the right leaves a linear equation in $u$. The remaining equalities and inequalities are transformed in the same way. Let $D_b$ and $D_g$ be positive diagonal row scales. The projected displacement solves
\begin{align}
\equationname{Minimum-distance objective for joint plant projection}
\widehat u=\arg\min_{u\in\mathbb R^{n_\chi}}\quad&\tfrac12u^\top u,\\
\text{subject to}\quad&
D_b^{-1}\{\mathcal H_{\log}(\vartheta)(\chi_{\rm raw}+D_\chi u)
-b_{\log}(\vartheta,c_{\rm in})\}=0,\notag\\
&\chi_{\rm raw}+D_\chi u\ge0,\notag\\
&D_g^{-1}\mathcal G(\vartheta)(\chi_{\rm raw}+D_\chi u)\le0,\notag\\
\equationname{Reconstruction of the projected physical plant response}
\widehat\chi&=\chi_{\rm raw}+D_\chi\widehat u.
\label{eq:plant_joint_projection}
\end{align}
The identity Hessian makes this a strictly convex quadratic program. In physical coordinates, the corresponding metric is $D_\chi^{-\top}D_\chi^{-1}$. Response scaling therefore determines which physical adjustments count as large, while positive row scaling improves numerical conditioning without altering the feasible set. \citet{SturmSilva2024} motivate species-sensitive weighting for projection, and \citet{Boyd2004} provide the convex optimization principles used for feasibility, uniqueness, and optimality.

The physically constrained plant surrogate extends the original ICSOR separation between statistical prediction and physical projection to a connected response \citep{Selerio2026}. The manifold projection of \citet{Valente2025} provides related output-projection context. Alternative enforcement locations include analytic neural-output constraints \citep{Beucler2021}, flux learning for mass conservation \citep{SturmWexler2020}, constrained training \citep{Ruckstuhl2021}, and governing-equation residual penalties \citep{Li2024}. These approaches introduce physical information during fitting or output construction. The present formulation instead imposes the declared plant balances and inventory limits directly on each admitted predicted response.

For a nonempty feasible set, closedness, coercivity, and strict convexity give a unique projection. Feasibility must nevertheless be checked jointly because the empirical endpoint can conflict with the physical constraints. For example, a fitted overflow-solids value exceeding every feed-solids value allowed by the other balances would contradict the implied relation $X_E\le X_N$. Such an endpoint would make an otherwise feasible physical polyhedron empty. Projection infeasibility is therefore recorded as a failure rather than hidden by further adjustment.

Figure~\ref{fig:plant_projection_architecture} separates the three inputs to the connected projection. The raw statistical response supplies the point being corrected. The derived plant relations supply the physical equalities and inequalities. The logarithmic model supplies a fitted endpoint equality. The distinction matters because the latter is empirical even though it enters the same mathematical equality system.

\begin{figure}[htbp]
\centering
\begin{tikzpicture}[>=Latex,
block/.style={draw=black!75,rounded corners,fill=cyan!6,align=center,font=\footnotesize,inner sep=4pt},
physical/.style={block,fill=black!5},
empirical/.style={block,dashed,fill=white},
flow/.style={->,thick,draw=black!80},
failflow/.style={->,dashed,draw=black!70},
lab/.style={font=\scriptsize,fill=white,inner sep=2pt}]
\node[block,text width=8.6cm] (inputs) at (0,0) {Declared controls and fresh-influent components};
\node[block,text width=4.0cm,minimum height=3.4cm] (raw) at (-4.65,-2.35)
{Raw connected response\\$\chi_{\rm raw}$\\$m,c_1,\ldots,c_5$\\$g_E,g_U,M_{\rm cl}$};
\node[physical,text width=4.0cm,minimum height=3.4cm] (phys) at (0,-2.35)
{Physical relations\\Mixer and inventories\\Clarifier balances\\Separation bounds\\Inventory bounds\\non-negativity};
\node[empirical,text width=4.0cm,minimum height=3.4cm] (closure) at (4.65,-2.35)
{Empirical closure\\Logarithmic model\\Positive overflow solids\\$t_X^\top g_E=q_E\widehat X_{E,\log}$};
\node[block,text width=8.6cm] (projection) at (0,-5.45)
{Joint nearest-response projection\\Minimum standardized projection\\Physical relations and fitted closure imposed together};
\node[physical,text width=8.6cm] (audit) at (0,-7.45)
{Independent projection audit\\Feasibility, non-negativity, and optimality residuals};
\node[block,text width=7.7cm] (output) at (-2.2,-9.45)
{Audited projected response $\widehat\chi$\\Used for the objective and trust screening};
\node[empirical,text width=3.1cm] (failure) at (4.5,-9.45)
{Projection fails\\or audit fails\\Retain status};
\draw[flow] (inputs.south)--(phys.north);
\draw[flow] (inputs.west)-|(raw.north);
\draw[flow] (inputs.east)-|(closure.north);
\draw[flow] (raw.south)--(projection.north west);
\draw[flow] (phys.south)--(projection.north);
\draw[flow] (closure.south)--(projection.north east);
\draw[flow] (projection)--(audit);
\draw[flow] (audit.south west)--(output.north) node[midway,left,lab] {Pass};
\draw[failflow] (audit.south east)--(failure.north) node[midway,right,lab] {Fail};
\end{tikzpicture}
\caption{Conceptual information flow in the connected projection. The dashed closure block denotes an empirically fitted relation rather than a physical identity. Projection corrects all response blocks jointly. The audit checks the imposed relations at their numerical tolerances.}
\label{fig:plant_projection_architecture}
\end{figure}

The returned response is therefore the nearest admissible state under the stated response scales and constraint set. Both the physical rows and the empirical endpoint determine that set.

\subsubsection{Geometry of simultaneous physical enforcement}

Figure~\ref{fig:plant_projection_geometry} gives a geometric interpretation of Equation~\eqref{eq:plant_joint_projection}. Define the scaled response coordinates $\eta=D_\chi^{-1}\chi$ and the corresponding raw point $\eta_{\rm raw}=D_\chi^{-1}\chi_{\rm raw}$. These coordinates describe the response itself, whereas $u=\eta-\eta_{\rm raw}$ is the standardized displacement used above. Consequently, minimizing $\tfrac12u^\top u$ selects the point in the complete feasible set nearest to the raw response in this scaled geometry. Positive diagonal scaling retains the meaning of non-negativity while expressing concentrations, component flows and inventory in comparable scale units.

The analytical illustration uses the equality $\eta_1+\eta_2=1$, non-negativity $\eta_1,\eta_2\ge0$, and an additional inequality $\eta_1\ge0.20$. The equality defines a line. Non-negativity reduces it to the segment between $(0,1)$ and $(1,0)$, and the additional inequality restricts it further to the green segment between $(0.20,0.80)$ and $(1,0)$. The last inequality stands in for a separator or inventory restriction. Its numerical value is chosen only to make the geometry visible; it is not a case-study operating limit or a new plant-model assumption.

\begin{figure}[htbp]
\centering
\begin{tikzpicture}[>=Latex,x=3.5cm,y=3.5cm,font=\small,
annotation/.style={fill=white,inner sep=2pt,align=center,font=\small}]
\fill[plantblue!4] (0,0) rectangle (1.30,1.85);
\begin{scope}
\clip (-.91,-.22) rectangle (1.31,1.91);
\draw[plantorange!45,densely dotted] (-.60,1.30) circle (.40);
\draw[plantorange!55,densely dotted] (-.60,1.30) circle (.70);
\draw[plantorange!75,dashed] (-.60,1.30) circle (.9433981);
\draw[plantblue!80!black,dashed,thick] (-.85,1.85)--(1.20,-.20);
\end{scope}
\draw[plantblue!55,line width=3pt] (0,1)--(1,0);
\draw[plantteal,line width=4pt] (.20,.80)--(1,0);
\draw[plantpurple,dash dot,thick] (.20,-.12)--(.20,1.65);
\draw[->,black!70] (-.87,0)--(1.37,0) node[right] {$\eta_1$};
\draw[->,black!70] (0,-.16)--(0,1.96) node[above] {$\eta_2$};
\foreach \tick in {0.2,1.0}{\draw (\tick,.025)--(\tick,-.025) node[below,font=\scriptsize] {$\tick$};}
\draw (.025,1)--(-.025,1);
\draw[plantorange,->,line width=1.2pt] (-.60,1.30)--(.20,.80);
\draw[plantpurple,->,densely dashed] (-.60,1.30)--(-.45,1.45);
\fill[plantorange] (-.60,1.30) circle (2.4pt);
\node[annotation,anchor=north east,text=plantorange!80!black] at (-.61,1.26) {Raw prediction\\$(-0.60,1.30)$};
\draw[red!70!black,thick] (-.477,1.423)--(-.423,1.477);
\draw[red!70!black,thick] (-.477,1.477)--(-.423,1.423);
\node[annotation,anchor=south,text=red!70!black] at (-.45,1.49) {Equalities only\\$(-0.45,1.45)$};
\draw[plantblue,fill=white,thick] (0,1) circle (2.0pt);
\fill[plantteal] (.20,.80) circle (2.6pt);
\node[annotation,anchor=north west,text=plantteal!75!black] at (.25,.77) {Joint projection\\$(0.20,0.80)$};
\fill[plantteal] (1,0) circle (1.8pt);
\node[annotation,anchor=west,text=plantblue!80!black] at (.49,1.48) {$\eta_1+\eta_2=1$};
\node[annotation,anchor=west,text=plantpurple] at (.49,1.27) {$\eta_1\ge0.20$};
\draw[plantteal,line width=4pt] (.53,1.05)--(.71,1.05);
\node[font=\small,anchor=west,text=plantteal!75!black] at (.75,1.05) {Joint feasible set};
\node[annotation,anchor=west,font=\scriptsize,align=left] at (.49,1.76) {Dotted/dashed circles:\\equal displacement from the raw state};
\end{tikzpicture}
\caption{Analytical illustration of joint plant-response projection in two scaled coordinates. The equality line, non-negativity and an additional inequality intersect to form the green feasible segment. The orange arrow gives the shortest admissible correction from the raw point, while the red cross shows an equality-only projection that remains negative. Circles denote equal standardized displacement from the raw point. The illustration explains the projection criterion and is not a simulated plant result.}
\label{fig:plant_projection_geometry}
\end{figure}

For the raw point $(-0.60,1.30)$, the nearest point on the equality line alone is $(-0.45,1.45)$. Its first coordinate is negative, so conservation alone does not make it admissible. The point $(0,1)$ satisfies the equality and non-negativity but still violates the additional inequality. The complete projection is instead $(0.20,0.80)$, where the smallest displacement circle that reaches the feasible segment meets its active boundary. Clipping the negative raw coordinate would give $(0,1.30)$ and violate the equality. A sequence of separate corrections therefore cannot be assumed to retain every previously satisfied relation.

In the connected plant, the same criterion acts on all mixer, reactor, outlet-flow and inventory coordinates together. A correction to one stream can require compensating changes at another location, and the fixed response scales determine the cost assigned to those movements. The diagram makes this simultaneous enforcement visible, but it does not prove agreement with the nonlinear kinetic or settling equations, nor does it establish predictive improvement relative to a reference outside the complete constraint set. The following weighted two-component example complements this geometric explanation by checking feasibility, multiplier signs, stationarity and complementarity explicitly.

\subsubsection{An independently auditable optimality certificate}

A simple two-component example illustrates how the projection can be checked independently. Let the raw prediction be $(-2,9)$, the coordinate scales be $(1,2)$, and the conserved total be 10. In physical coordinates,
\begin{equation}
\equationname{Illustrative weighted projection with mass conservation and non-negativity}
\min_{c_1,c_2}\ \frac12(c_1+2)^2+\frac18(c_2-9)^2,
\qquad c_1+c_2=10,\quad c_1\ge0,\quad c_2\ge0.
\end{equation}
The coefficient $1/8$ equals $1/(2\times2^2)$. A one-unit physical change in the second component is therefore less costly than the same change in the first because its scale is larger.

Ignoring non-negativity initially and introducing an equality multiplier $\lambda$ gives
\begin{align}
\equationname{Illustrative projection stationarity for the first component}
c_1+2+\lambda&=0,\\
\equationname{Illustrative projection stationarity for the second component}
\tfrac14(c_2-9)+\lambda&=0,\\
\equationname{Illustrative projection equality for the component total}
c_1+c_2&=10.
\end{align}
The first two equations imply $c_1=-2-\lambda$ and $c_2=9-4\lambda$. Substitution into the conserved-total equation gives $\lambda=-0.6$ and $(c_1,c_2)=(-1.4,11.4)$. Mass conservation alone therefore leaves the first component negative.

Imposing the active boundary $c_1=0$ forces $c_2=10$. Let $\mu_1,\mu_2\ge0$ be multipliers for $-c_1\le0$ and $-c_2\le0$. Because $c_2>0$, its inequality is inactive and $\mu_2=0$. Stationarity for component 2 gives $\lambda=-0.25$, and stationarity for component 1 gives $2-0.25-\mu_1=0$, so $\mu_1=1.75$. Both inequality multipliers are non-negative and both complementarity products vanish. These conditions certify the boundary solution instead of merely showing that it is feasible.

In standardized coordinates, $u=(2,0.5)^\top$, and the same problem becomes
\begin{equation}
\equationname{Standardized matrices of the illustrative two-component projection}
E=\begin{bmatrix}1&2\end{bmatrix},\quad e=3,\quad
G_u=\begin{bmatrix}-1&0\\0&-2\end{bmatrix},\quad
g=\begin{bmatrix}-2\\9\end{bmatrix}.
\end{equation}
The equality is $u_1+2u_2=3$, while the inequality rows are the physical non-negativity requirements after substitution. This form generalizes directly to the larger plant projection.

At fixed inputs, write the standardized problem as
\begin{equation}
\equationname{Standardized quadratic projection with equalities and inequalities}
\min_u\tfrac12u^\top u,\qquad Eu=e,\qquad G_uu\le g.
\label{eq:plant_standard_projection}
\end{equation}
Let $\lambda^=$ denote unrestricted equality multipliers and $\lambda^\le\ge0$ the inequality multipliers. Optimality requires primal feasibility, dual feasibility, stationarity, and complementarity,
\begin{align}
\equationname{Stationarity condition for the plant quadratic projection}
u+E^\top\lambda^=+G_u^\top\lambda^\le&=0,\\
\equationname{Complementarity condition for the projection inequality multipliers}
\lambda^\le_j(g_j-(G_uu)_j)&=0\quad\text{for every }j.
\label{eq:plant_projection_optimality}
\end{align}

The dual value follows by collecting the linear terms in the Lagrangian. Define $v=E^\top\lambda^=+G_u^\top\lambda^\le$. Completing the square gives
\begin{align}
\equationname{Completed-square expansion of the projection Lagrangian}
\tfrac12u^\top u+v^\top u-e^\top\lambda^=-g^\top\lambda^\le
&=\tfrac12\|u+v\|_2^2-\tfrac12\|v\|_2^2\\
\equationname{Target terms in the completed-square projection Lagrangian}
&\quad-e^\top\lambda^=-g^\top\lambda^\le.
\end{align}
The first term is minimized at $u=-v$, giving
\begin{equation}
\equationname{Dual objective of the standardized quadratic projection}
d(\lambda^=,\lambda^\le)=
-\tfrac12\|E^\top\lambda^=+G_u^\top\lambda^\le\|_2^2
-e^\top\lambda^=-g^\top\lambda^\le.
\end{equation}
For primal- and dual-feasible quantities, the resulting gap is
\begin{equation}
\equationname{Primal-dual gap and its non-negative decomposition}
\Gamma=\tfrac12\|u+E^\top\lambda^=+G_u^\top\lambda^\le\|_2^2
+(\lambda^\le)^\top(g-G_uu)\ge0.
\label{eq:plant_projection_gap}
\end{equation}
The gap vanishes exactly when stationarity and complementarity hold. This provides a mathematical certificate of the unique projection \citep{Boyd2004}. Numerical acceptance still checks each residual separately because a small aggregate gap does not replace explicit equality or non-negativity checks.

The two-component example permits direct verification of every certificate quantity. With $\lambda^{=}=-0.25$ and $\lambda^\le=(1.75,0)^\top$,
\begin{align}
\equationname{Illustrative cancellation of the stationarity residual}
E^\top\lambda^=+G_u^\top\lambda^\le&=(-2,-0.5)^\top=-u,\\
\equationname{Illustrative inequality slack vector}
g-G_uu&=(0,10)^\top,\\
\equationname{Illustrative value of the primal projection objective}
\tfrac12u^\top u&=\tfrac12(4+0.25)=2.125,\\
\equationname{Illustrative value of the dual projection objective}
d&=-2.125-3(-0.25)-[-2(1.75)]=2.125.
\end{align}
The primal and dual values agree. The stationarity residual is zero coordinate by coordinate, and the complementarity sum is $1.75(0)+0(10)=0$. A multiplier with the wrong sign would fail dual feasibility even if another residual were small.

Every distinct trial therefore receives an independently initialized projection solve. An operator-splitting quadratic solver of the type studied by \citet{Stellato2020} supplies the numerical candidate, while the audit reconstructs multipliers independently through bounded-variable least squares. Numerical retries may alter solver settings, but they do not alter the mathematical projection problem.

\subsubsection{Error bounds in the projection metric}

Let $\mathcal C$ denote the fixed-input feasible set and define the weighted norm
$\|v\|_{D_\chi^{-1}}=\|D_\chi^{-1}v\|_2$. Its scalar form is
\begin{equation}
\equationname{Weighted response norm from coordinate scales}
\|v\|_{D_\chi^{-1}}^2=\sum_{k=1}^{n_\chi}\frac{v_k^2}{s_k^2}.
\end{equation}
The geometric effect of projection can be derived most clearly in standardized coordinates. Write $a_k=\chi_{{\rm raw},k}/s_k$, $q_k=\widehat\chi_k/s_k$, and $t_k=\chi_{{\rm ref},k}/s_k$. The transformed feasible set remains convex because the scaling is an invertible linear transformation.

For any feasible comparison point $z$, the segment $q+\tau(z-q)$ remains feasible for $0\le\tau\le1$. Because $q$ minimizes distance from $a$, the projection objective cannot decrease along that segment. Its directional derivative at zero therefore satisfies
\begin{equation}
\equationname{Variational inequality characterizing a Euclidean projection}
\sum_k(q_k-a_k)(z_k-q_k)\ge0.
\end{equation}
If the reference itself is feasible, choose $z=t$. Reversing both signs gives $\sum_k(a_k-q_k)(q_k-t_k)\ge0$. For each coordinate,
\begin{align}
\equationname{Coordinatewise squared-error decomposition}
(a_k-t_k)^2&=(a_k-q_k)^2+(q_k-t_k)^2\\
\equationname{Cross term in the coordinatewise squared-error decomposition}
&\quad+2(a_k-q_k)(q_k-t_k).
\end{align}
Summing these identities gives
\begin{equation}
\equationname{Squared-error bound for a feasible reference}
\sum_k(q_k-t_k)^2\le\sum_k(a_k-t_k)^2-\sum_k(a_k-q_k)^2.
\end{equation}
The final sum is exactly the squared standardized projection distance. Returning to physical coordinates, if $\chi_{\rm ref}\in\mathcal C$,
\begin{equation}
\equationname{Weighted error reduction bound when the reference is projection-feasible}
\|\widehat\chi-\chi_{\rm ref}\|_{D_\chi^{-1}}^2
\le\|\chi_{\rm raw}-\chi_{\rm ref}\|_{D_\chi^{-1}}^2
-\|\widehat u\|_2^2.
\label{eq:plant_projection_error_exact}
\end{equation}
Projection therefore cannot increase aggregate squared error in this particular weighted metric to an exactly feasible reference. This result does not imply improvement in every component, every physical-unit score, or every operating objective, which are evaluated separately.

The earlier two-component problem illustrates this distinction. Let the feasible reference be $(3,7)$. Raw squared weighted error is $(-2-3)^2+(9-7)^2/4=26$, projected squared error is $(0-3)^2+(10-7)^2/4=11.25$, and projection-size squared is $2^2+0.5^2=4.25$. The bound gives $11.25\le26-4.25=21.75$. Yet the second component's absolute error increases from 2 to 3. Aggregate weighted improvement therefore does not require every component error to decrease.

The fitted overflow equality generally differs from the mechanistic endpoint, so the mechanistic reference need not lie in $\mathcal C$. Define its distance from the complete projection set by
\begin{equation}
\equationname{Distance from the reference state to the projection set}
\delta_{\rm feas}=\|P_{\mathcal C}(\chi_{\rm ref})-\chi_{\rm ref}\|_{D_\chi^{-1}}.
\end{equation}
Let $p$ denote the closest feasible point to $t$ in standardized coordinates. Applying the previous segment argument to $z=p$ gives
\begin{align}
\equationname{Cross-term decomposition using a projected reference}
\sum_k(a_k-q_k)(q_k-t_k)
&=\sum_k(a_k-q_k)(q_k-p_k)\\
\equationname{Reference-infeasibility term in the cross-term decomposition}
&\quad+\sum_k(a_k-q_k)(p_k-t_k).
\end{align}
The first term is non-negative. By Cauchy--Schwarz, the magnitude of the second is bounded by
\begin{equation}
\equationname{Product of projection displacement and reference infeasibility}
\left[\sum_k(a_k-q_k)^2\right]^{1/2}
\left[\sum_k(p_k-t_k)^2\right]^{1/2}
=r_{\rm corr}\delta_{\rm feas}.
\end{equation}
Here $r_{\rm corr}=\|\chi_{\rm raw}-P_{\mathcal C}(\chi_{\rm raw})\|_{D_\chi^{-1}}$. Expanding the squared error and applying this bound yields
\begin{align}
\equationname{Squared-error bound including reference infeasibility}
\|P_{\mathcal C}(\chi_{\rm raw})-\chi_{\rm ref}\|_{D_\chi^{-1}}^2
&\le\|\chi_{\rm raw}-\chi_{\rm ref}\|_{D_\chi^{-1}}^2
-r_{\rm corr}^2+2r_{\rm corr}\delta_{\rm feas}\\
\equationname{Upper error bound using squared reference infeasibility}
&\le\|\chi_{\rm raw}-\chi_{\rm ref}\|_{D_\chi^{-1}}^2+\delta_{\rm feas}^2.
\end{align}
Taking square roots gives the more interpretable bound
\begin{equation}
\equationname{Distance bound when the reference is not projection-feasible}
\|P_{\mathcal C}(\chi_{\rm raw})-\chi_{\rm ref}\|_{D_\chi^{-1}}
\le\|\chi_{\rm raw}-\chi_{\rm ref}\|_{D_\chi^{-1}}+\delta_{\rm feas}.
\label{eq:plant_projection_error_general}
\end{equation}
The distance term includes any mismatch between the mechanistic reference and the full projection set, including the empirical overflow equality. The weighted-projection discussion of \citet{SturmSilva2024} similarly separates mass conservation from predictive improvement. Independent original-model evaluation then checks the biological kinetics, nonlinear settling behavior, and local dynamic stability that projection does not enforce.

A hypothetical example shows why the empirical closure matters. Let the conserved total remain 10, but additionally require the first component to equal a fitted value of 4. The feasible set then contains only $(4,6)$. Suppose both the mechanistic reference and raw prediction are $(2,8)$ with scales $(1,2)$. The raw error is zero, but projection introduces weighted error $\sqrt{(4-2)^2+(6-8)^2/4}=\sqrt5$. The reference distance to the projection set is also $\sqrt5$, so the general bound is attained exactly. Treating only the physical identities as the feasible set would have incorrectly suggested the stronger exact-feasibility guarantee.

\subsubsection{Control-dependent sensitivities}

If the projection matrices were fixed and only the right-hand side depended affinely on parameters, the solution of a parametric quadratic program could be piecewise affine \citep{Tondel2003}. In the present formulation, however, both matrix coefficients and targets vary with the controls. Hydraulic ratios, inverse flow factors, the statistical feature vector, and the exponential overflow endpoint all depend on $\vartheta$. The projected plant response is therefore generally piecewise smooth where regularity conditions hold, with active-set changes defining boundaries between smooth regions.

A scalar example illustrates why derivatives of the constraint coefficients must be retained. Consider
\begin{equation}
\equationname{Scalar equality-constrained projection with a varying parameter}
\min_u\tfrac12u^2,\qquad b(\theta)u=e(\theta),
\end{equation}
with $b(\theta)\ne0$. Stationarity and feasibility are $u+b\lambda=0$ and $bu=e$, or
\begin{equation}
\equationname{Optimality system for the scalar parametric projection}
\begin{bmatrix}1&b\\b&0\end{bmatrix}
\begin{bmatrix}u\\\lambda\end{bmatrix}
=\begin{bmatrix}0\\e\end{bmatrix}.
\end{equation}
Differentiating both equations gives $u'+b'\lambda+b\lambda'=0$ and $b'u+bu'=e'$. Rearrangement yields
\begin{equation}
\equationname{Derivative system for the scalar parametric projection}
\begin{bmatrix}1&b\\b&0\end{bmatrix}
\begin{bmatrix}u'\\\lambda'\end{bmatrix}
=\begin{bmatrix}-b'\lambda\\e'-b'u\end{bmatrix}.
\end{equation}
The optimality matrix remains unchanged, while the right-hand side contains derivatives of both the coefficient and the target.

For $b=1+\theta$ and $e=\theta$, the solution is $u=\theta/(1+\theta)$ and $\lambda=-\theta/(1+\theta)^2$. At $\theta=1$, $u=0.5$ and $\lambda=-0.25$, giving
\begin{equation}
\equationname{Numerical example of the parametric projection derivative}
\begin{bmatrix}1&2\\2&0\end{bmatrix}
\begin{bmatrix}u'\\\lambda'\end{bmatrix}
=\begin{bmatrix}0.25\\0.5\end{bmatrix}.
\end{equation}
The second row gives $u'=0.25$, and the first gives $\lambda'=0$. Direct differentiation of $u=\theta/(1+\theta)$ gives the same result. Omitting $b'$ would instead differentiate a different optimization problem in which the constraint coefficient was incorrectly treated as constant.

For a stable active set, stack all equality rows and active inequality rows into $B(\vartheta)$, with corresponding right-hand side $b_a(\vartheta)$. The stationary system is
\begin{equation}
\equationname{Linear system for a fixed active set of projection constraints}
\begin{bmatrix}I&B^\top\\B&0\end{bmatrix}
\begin{bmatrix}u\\\lambda_a\end{bmatrix}
=\begin{bmatrix}0\\b_a\end{bmatrix}.
\label{eq:plant_active_system}
\end{equation}
Applying the product rule to $u+B^\top\lambda_a=0$ and $Bu=b_a$ gives
\begin{align}
\equationname{Differentiated stationarity for a control-dependent active set}
\partial_j u+B^\top\partial_j\lambda_a&=-(\partial_jB)^\top\lambda_a,\\
\equationname{Differentiated feasibility for a control-dependent active set}
B\partial_j u&=\partial_jb_a-(\partial_jB)u.
\end{align}
These equations combine into
\begin{equation}
\equationname{Joint sensitivity system for a control-dependent active set}
\begin{bmatrix}I&B^\top\\B&0\end{bmatrix}
\begin{bmatrix}\partial_j u\\\partial_j\lambda_a\end{bmatrix}
=\begin{bmatrix}-(\partial_jB)^\top\lambda_a\\
\partial_jb_a-(\partial_jB)u\end{bmatrix}.
\label{eq:plant_active_sensitivity}
\end{equation}
Because $D_\chi$ is fixed after fitting, the physical-response derivative is $\partial_j\widehat\chi=\partial_j\chi_{\rm raw}+D_\chi\partial_j\widehat u$.

The fitted overflow endpoint contributes an additional control-dependent target. Differentiating the scalar log prediction gives $\partial_j(\widehat\beta^\top\phi)=\widehat\beta^\top\partial_j\phi$. Exponentiation then gives $\partial_j\widehat X_{E,\log}=\widehat X_{E,\log}\widehat\beta^\top\partial_j\phi$, and the flow-weighted endpoint has derivative
\begin{equation}
\equationname{Derivative of the control-dependent overflow-solids target}
\partial_j(q_E\widehat X_{E,\log})
=(\partial_jq_E)\widehat X_{E,\log}
+q_E\widehat X_{E,\log}\widehat\beta_{E,\gamma_E}^\top\partial_j\phi.
\end{equation}
Ignoring this term would again differentiate a different optimization problem.

These derivatives are used only when the active rows are independent, the stationary matrix is well conditioned, active inequality multipliers are strictly positive, and independent perturbation checks pass. Near degeneracy, different active-set descriptions can correspond to the same unique projected state. Uniqueness of the projection therefore does not guarantee a reliable classical derivative. A value-only outer method remains available when the derivative chain cannot be validated.

\subsection{Treatment priorities, resource measures, and operating safeguards}

\subsubsection{A common engineering objective}

The effluent-quality map in Table~\ref{tab:component_basis_composites} gives
\begin{equation}
\equationname{Effluent component-to-water-quality map}
z_E=I_{\rm comp}c_E=\operatorname{col}(\mathrm{COD},\mathrm{TN},\mathrm{TP},\mathrm{TSS}).
\end{equation}
Let the diagonal entries of $D_T$ be $T_{COD},T_{TN},T_{TP},T_{TSS}$. Dividing each effluent quantity by its own scale makes the four quality terms dimensionless. With non-negative quality weights $\varpi_Q$ satisfying $\mathbf1^\top\varpi_Q=1$, the weighted quality score is
\begin{equation}
\equationname{Scalar expansion of the weighted effluent-quality objective}
\begin{aligned}
j_Q={}&\varpi_{Q,COD}\frac{z_{E,COD}}{T_{COD}}
+\varpi_{Q,TN}\frac{z_{E,TN}}{T_{TN}}\\
&+\varpi_{Q,TP}\frac{z_{E,TP}}{T_{TP}}
+\varpi_{Q,TSS}\frac{z_{E,TSS}}{T_{TSS}}.
\end{aligned}
\end{equation}
Equivalently,
\begin{equation}
\equationname{Vector form of the weighted effluent-quality objective}
j_Q=\varpi_Q^\top D_T^{-1}z_E.
\label{eq:plant_quality_objective}
\end{equation}
The four reported totals summarize different treatment concerns, but they do not identify the individual biological transformations that produce them. TN may remain nearly constant while ammonium is converted to nitrite or nitrate, and TP may decrease across the clarifier because particulate phosphorus is separated rather than because a new clarifier reaction occurs. The nutrient-control guidance of \citet{USEPA2009} and the component-resolved framework of \citet{Zhang2026} therefore support examining species and solids flows alongside aggregate treatment metrics.

Resource terms are normalized according to their physical interpretation. A bounded setting is placed between zero and one by subtracting its lower bound and dividing by its allowed range. Thus $H=18$ h in the interval 6--36 h gives $(18-6)/(36-6)=0.4$. Aeration is treated differently because its resource proxy is the total transfer capacity $\sum_iV_i k_La_i$, which has units of m$^3$ d$^{-1}$. Wasted solids are proportional to $Q_0wX_U$, and normalization by $Q_0w^UX_U^{\rm ref}$ cancels $Q_0$. The resulting terms are
\begin{align}
\equationname{Normalized retention-time and aeration objective terms}
j_H&=\frac{H-H^L}{H^U-H^L},\qquad
j_A=\frac{\sum_{i\in\mathcal A}V_i k_La_i}{E_A^{\rm ref}},\\
\equationname{Normalized mixed-liquor and return-sludge recycle objective terms}
j_I&=\frac{r_I-r_I^L}{r_I^U-r_I^L},\qquad
j_R=\frac{r_R-r_R^L}{r_R^U-r_R^L},\\
\equationname{Normalized waste-solids objective term}
j_W&=\frac{wX_U}{w^UX_U^{\rm ref}}.
\label{eq:plant_resource_objective}
\end{align}
The set $\mathcal A$ contains the aerated reactors. The reference $E_A^{\rm ref}>0$ is fixed before optimization as the maximum transfer-capacity numerator over the declared operating box. Fixing this denominator prevents a larger selected volume from being canceled by a simultaneously changing reference.

The aeration term represents transfer capacity rather than actual oxygen transfer, which also depends on the dissolved-oxygen deficit in Equation~\eqref{eq:plant_reactor}. The energy-management model of \citet{Miederer2025} separates aeration from other energy demands, while \citet{Zhang2026} explicitly retains the oxygen-transfer driving force. The present objective terms should therefore be read as dimensionless engineering measures of capacity and resource use.

Let the six overall coefficients correspond to the quality, HRT, aeration, internal recycle, return sludge, and wasting terms. The scalar objective is
\begin{equation}
\equationname{Scalar expansion of the six-term engineering objective}
J=\varpi_1j_Q+\varpi_2j_H+\varpi_3j_A+\varpi_4j_I+\varpi_5j_R+\varpi_6j_W.
\end{equation}
Its compact form is
\begin{equation}
\equationname{Vector engineering objective and normalized priority weights}
J(\vartheta,\chi,\varpi)=\varpi^\top
\operatorname{col}(j_Q,j_H,j_A,j_I,j_R,j_W),\qquad
\varpi\ge0,\quad\mathbf1^\top\varpi=1.
\label{eq:plant_engineering_objective}
\end{equation}
The vector $\varpi$ determines the trade-off between quality, capacity, and resource use, while $\varpi_Q$ determines the relative importance of COD, TN, TP, and TSS within the quality term. These coefficients encode preferences; they are not percentages of the realized objective.

For equal-volume reactors,
\begin{equation}
\equationname{Aeration objective for the illustrative equal-volume reactor train}
j_A=\frac{H}{36\ \mathrm{h}}\frac{a_3+a_4+a_5}{3}.
\label{eq:plant_aeration_proxy_example}
\end{equation}
This expression follows from $V_i=Q_0H/(24n_r)$ and $k_La_i=47a_i$. If all three aeration controls are at one, reducing $H$ from 36 h to 12 h reduces $j_A$ from one to one-third. The aeration settings have not changed, but the total installed transfer capacity has changed with reactor volume. A comparison based only on the dimensionless aeration controls would miss this capacity effect.

The baseline priorities are
\begin{equation}
\equationname{Declared engineering and effluent-quality priority vectors}
\varpi=(0.50,0.15,0.20,0.05,0.05,0.05)^\top,\qquad
\varpi_Q=\tfrac14\mathbf1_4.
\label{eq:plant_objective_priorities}
\end{equation}
Quality receives the largest overall weight, while aeration receives the largest resource weight. The importance of oxygen supply and solids age to treatment and energy use is discussed by \citet{Leu2012}. The exact values remain case-study choices. A different set of weights would define a different decision problem and would require its own reported optimization results.

A hypothetical calculation illustrates how the six terms combine. Let the effluent values be $(40,5,1,4)$ mg L$^{-1}$ and the fixed quality scales be $(100,10,2,20)$ mg L$^{-1}$. With equal quality priorities, $j_Q=(0.4+0.5+0.5+0.2)/4=0.4$. Let $H=18$ h, $(a_3,a_4,a_5)=(0.8,0.6,0.4)$, $r_I=1$, $r_R=0.5$, $w=0.02$, and $X_U=6{,}000$ g m$^{-3}$.

The total biological volume is $Q_0H/24=7{,}500$ m$^3$, so each of the five equal tanks has volume 1{,}500 m$^3$. Aeration capacity is $1{,}500(47)(0.8+0.6+0.4)=126{,}900$ m$^3$ d$^{-1}$. Table~\ref{tab:plant_hypothetical_objective} lists the normalized values and weighted contributions. Their sum is 0.353.

\begin{table}[htbp]
\centering\small
\caption{Hypothetical calculation of dimensionless objective contributions.}
\label{tab:plant_hypothetical_objective}
\begin{tabular}{llll}
\toprule Term & Calculation & Value & Weighted value\\\midrule
$j_Q$ & $(0.4+0.5+0.5+0.2)/4$ & 0.4 & 0.200\\
$j_H$ & $(18-6)/(36-6)$ & 0.4 & 0.060\\
$j_A$ & $126{,}900/423{,}000$ & 0.3 & 0.060\\
$j_I$ & $(1-0)/(4-0)$ & 0.25 & 0.0125\\
$j_R$ & $(0.5-0.25)/(1.25-0.25)$ & 0.25 & 0.0125\\
$j_W$ & $(0.02)(6{,}000)/[(0.05)(15{,}000)]$ & 0.16 & 0.008\\
\bottomrule
\end{tabular}
\end{table}

The overall quality weight of 0.50 multiplies $j_Q$ only after its four internal quality weights have been applied. An internal quality weight of one-quarter is therefore not an overall objective weight of one-quarter. Likewise, an objective reference value is not necessarily a feasibility limit. In particular, $j_W$ may exceed one if its numerator exceeds the declared reference. Normalization changes scale; it does not impose a hard bound.

\subsubsection{Solids retention time and clarifier loading}

Solids retention time combines stored solids mass with the rate at which solids leave the external plant boundary. For two reactors, stored solids would be $V_1X_1+V_2X_2+M_{\rm cl}$ in grams. The corresponding loss rate would be $Q_EX_E+Q_0wX_U$ in grams per day. Their ratio therefore has units of days. The same accounting applies to the five-reactor case.

External solids leave only through the clarifier overflow and the waste-sludge stream,
\begin{equation}
\equationname{External solids loss through effluent and waste sludge}
\dot M_X=Q_0(q_EX_E+wX_U).
\label{eq:plant_external_solids_loss}
\end{equation}
Return sludge is internal and does not contribute to external solids loss. With $X_i=t_X^\top c_i$, the plant solids retention time is
\begin{equation}
\equationname{Plant solids retention time from inventory and external loss}
\theta_X=\frac{\sum_{i=1}^{n_r}V_iX_i+M_{\rm cl}}{\dot M_X},
\qquad\dot M_X\ge\epsilon_X>0.
\label{eq:plant_solids_age}
\end{equation}
The clarifier inventory is part of the numerator. Omitting it would underestimate the solids age for a given external loss rate. Similarly, the wasting ratio alone cannot determine $\theta_X$ because both the product $wX_U$ and the internal solids inventory matter \citep{Alex2008}.

For illustration, let $Q_0=10{,}000$ m$^3$ d$^{-1}$, $w=0.02$, $X_E=10$, and $X_U=6{,}000$ g m$^{-3}$. The two external loss terms are
\begin{align}
\equationname{Illustrative effluent solids-loss calculation}
Q_0q_EX_E&=10{,}000(0.98)(10)=98{,}000\ \mathrm{g\,d^{-1}},\\
\equationname{Illustrative waste-sludge solids-loss calculation}
Q_0wX_U&=10{,}000(0.02)(6{,}000)=1{,}200{,}000\ \mathrm{g\,d^{-1}}.
\end{align}
Their total is 1{,}298{,}000 g d$^{-1}$. If the biological reactors contain 22{,}500{,}000 g of solids and the clarifier stores 7{,}500{,}000 g, total inventory is 30{,}000{,}000 g. The corresponding illustrative solids retention time is $30{,}000{,}000/1{,}298{,}000=23.11$ d.

Clarifier surface loading and solids loading describe different physical burdens on the separator. Surface loading divides overflow volume rate by clarifier area, while solids loading divides the feed solids mass rate by area \citep{Henze2008}. Because solids are represented in grams, a factor of $10^{-3}$ converts to kilograms:
\begin{equation}
\equationname{Clarifier surface and solids loading measures}
\ell_{\rm surface}=\frac{Q_0q_E}{A_{\rm cl}},\qquad
\ell_{\rm solids}=\frac{10^{-3}Q_0q_CX_N}{A_{\rm cl}}.
\label{eq:plant_loading_measures}
\end{equation}
The resulting units are m d$^{-1}$ and kg m$^{-2}$ d$^{-1}$, respectively. These quantities, together with $\theta_X$, are descriptive in the baseline comparison rather than optimization constraints.

For the same hypothetical fresh flow, take $r_R=0.5$, so $q_C=1.5$, with $A_{\rm cl}=1{,}500$ m$^2$ and $X_N=3{,}000$ g m$^{-3}$. The surface loading is $9{,}800/1{,}500=6.53$ m d$^{-1}$. The clarifier-feed solids rate is $10{,}000(1.5)(3{,}000)=45{,}000{,}000$ g d$^{-1}$, or 45{,}000 kg d$^{-1}$, giving a solids loading of 30 kg m$^{-2}$ d$^{-1}$. Using underflow solids in place of clarifier-feed solids would define a different loading measure.

The common engineering safeguards are
\begin{equation}
\equationname{Underflow, feed-solids, and external-loss safeguards}
X_U\le X_U^{\max},\qquad X_N\ge X_F^{\min}>0,\qquad
\dot M_X\ge\epsilon_X>0.
\label{eq:plant_engineering_safeguards}
\end{equation}
The feed-solids floor keeps particulate outlet reconstruction defined, while the external-loss floor keeps solids retention time defined. The underflow limit represents the admitted thickening or solids-handling range. Their numerical values are plant-specific inputs fixed before search. An underflow limit should reflect accepted sludge-handling conditions \citep{Henze2008}, whereas the feed and loss floors should prevent undefined calculations without excluding physically relevant low-solids operation. These safeguards are separate from the fixed $X_U^{\rm ref}=15{,}000$ g m$^{-3}$ used only to normalize the wasting objective.

\subsubsection{Trust diagnostics for operating search}

The trust diagnostics restrict the operating search to specified levels of data support and model consistency. They do not replace mechanistic verification. Instead, they test whether a candidate lies within an empirically supported part of the surrogate response and whether the projected state remains close to the raw prediction and compatible with selected mechanistic relations. The fixed thresholds therefore define the region in which the surrogate is permitted to support candidate selection. This role is consistent with approximation-management approaches such as \citet{Alexandrov1998} and \citet{Hameed2026}.

The projection diagnostic measures the root-mean-square standardized correction. For the hypothetical two-component projection $u=(2,0.5)$, it is $\sqrt{(4+0.25)/2}=1.458$. Dividing by the square root of the response dimension makes the quantity interpretable as an average standardized movement rather than a total movement that grows automatically with the number of response coordinates.

Feature leverage measures input support in the fitted quadratic feature space. In a hypothetical intercept-plus-one-feature example with development features $(-\sqrt{3/2},0,\sqrt{3/2})$ and $\gamma=1$,
\begin{equation}
\equationname{Illustrative regularized feature Gram matrix}
\Phi^\top\Phi+3\gamma R_\phi=
\begin{bmatrix}3&0\\0&6\end{bmatrix}.
\end{equation}
Its inverse is $\operatorname{diag}(1/3,1/6)$. At query feature $(1,z)$, leverage is $1/3+z^2/6$: it equals $1/3$ at the center, $1/2$ at $z=1$, and $11/6$ at $z=3$. The maximum development leverage in this example is $7/12$. The calculation shows how joint feature support weakens as a query moves away from the fitted design, even when individual controls remain inside their coordinate bounds.

The particulate split diagnostic asks whether all particulate components follow the common recovery relation used by Equation~\eqref{eq:plant_outlet_reconstruction}. Suppose two particulate components have unit TSS weights and feed concentrations $(40,60)$, so $X_N=100$. An underflow component-flow vector $(20,30)$ has total flow 50 and yields numerators $100(20)-50(40)=0$ and $100(30)-50(60)=0$. Reallocating the same total underflow solids to $(25,25)$ preserves aggregate solids but gives numerators 500 and $-500$. With fixed split scales of 500, the standardized residuals are 1 and $-1$, producing an RMS diagnostic of one. Aggregate solids consistency therefore does not establish composition consistency.

The reactor diagnostic measures the residual of the full smooth kinetic equations at the projected state. In the hypothetical two-component reactor discussed earlier, the projected concentrations $(9,6)$ preserve the upstream total $(12,3)$. With $d=2$ and reaction rate 8, the component residuals are $2(12-9)-8=-2$ and $2(3-6)+8=2$. If each balance scale is 10, the standardized residuals are $(-0.2,0.2)$ and the RMS diagnostic is 0.2. The invariant is satisfied even though the kinetic equations are not. This example makes the distinction between physical conservation and mechanistic-equation agreement explicit.

The four diagnostics are
\begin{align}
\equationname{Root-mean-square standardized projection displacement diagnostic}
\kappa_{\rm corr}&=\frac{\|D_\chi^{-1}(\widehat\chi-\chi_{\rm raw})\|_2}{\sqrt{n_\chi}},\\
\equationname{Ridge feature-leverage diagnostic}
\kappa_{\rm lev}&=\phi^\top(\Phi^\top\Phi+n_{\rm dev}\gamma R_\phi)^{-1}\phi,\\
\equationname{Particulate split-consistency diagnostic}
\kappa_{\rm split}&=\left[\frac1{F_X}\sum_{f\in X}
\left\{\frac{X_Ng_{U,f}-(t_X^\top g_U)c_{n_r,f}}{d_{{\rm split},f}}\right\}^2\right]^{1/2},\\
\equationname{Scaled reactor-equation residual diagnostic}
\kappa_{\rm reac}&=\frac{\|D_{\rm bal}^{-1}f_{{\rm sm},\rm reactor}(\vartheta,\widehat\chi)\|_2}{\sqrt{n_rF}}.
\label{eq:plant_trust_diagnostics}
\end{align}
The first measures the size of the physical correction applied to the raw surrogate response. The second measures data support in the fitted feature space. The third tests compatibility with the common particulate-recovery assumption. The fourth checks the smooth reactor equations at the projected concentrations. Clarifier-layer equations are assessed separately through the layer-resolved mechanistic model.

Residual scaling is based on the sizes of the physical terms before cancellation. Suppose a balance row contains terms $(100,-100)$ in one development state and $(200,-200)$ in another. Its net residual is zero in both states, but the individual terms are large. Their mean squared magnitude is $(100^2+100^2+200^2+200^2)/4=25{,}000$, giving an RMS scale of 158.11. Scaling only the net balance would discard this information. For a row with $p$ physical terms $b_{ij}$ over $n_{\rm dev}$ development states,
\begin{equation}
\equationname{Term-based scale for an independently checked residual}
d_b=\max\left\{1,\left[\frac1{n_{\rm dev}p}
\sum_{i=1}^{n_{\rm dev}}\sum_{j=1}^{p}b_{ij}^2\right]^{1/2}\right\}.
\label{eq:plant_term_scales}
\end{equation}
The same rule defines inequality and split scales. The overflow-closure scale uses the mechanistic overflow flow and the out-of-fold fitted endpoint flow as its two terms. The floor of one is expressed in the numerical units of the row and is not a concentration floor.

Projection, split, and reactor thresholds are set at the nearest-rank 95th percentile of grouped out-of-fold development diagnostics. If the sorted values are $v_{(1)},\ldots,v_{(n_{\rm dev})}$, the threshold is $v_{(\lceil0.95n_{\rm dev}\rceil)}$. The leverage threshold is the maximum development leverage under the final fitted feature map. These values are fixed before optimization. A diagnostic whose empirical threshold is zero remains constrained to zero; the value one is used only as a numerical row-scale floor.

The surrogate decision problem is therefore
\begin{equation}
\equationname{Surrogate-assisted plant operating optimization problem}
\begin{aligned}
\min_{\vartheta}\quad&J(\vartheta,\widehat\chi(\vartheta,c_{\rm in}),\varpi),\\
\text{subject to}\quad&\vartheta\in\mathcal V,\\
&g_{\rm eng}(\vartheta,\widehat\chi)\le0,\\
&g_{\rm trust}(\vartheta,\chi_{\rm raw},\widehat\chi)\le0,\\
&\widehat\chi\text{ satisfies Equation~\eqref{eq:plant_joint_projection}}.
\end{aligned}
\label{eq:plant_surrogate_optimization}
\end{equation}
The lower projection problem is convex, but the upper operating problem remains generally nonconvex because the statistical response, the control-dependent constraints, and the engineering objective are nonlinear in the controls. This nested structure is consistent with the bilevel perspective described by \citet{Colson2007}. The trust-region analysis of \citet{Alexandrov1998} further motivates supplementing average prediction accuracy with local admission checks near candidate operating points.

\subsection{Bounded products and convex relaxation}

The plant formulation also contains products of bounded variables, which helps explain why process optimization remains nonlinear even when the statistical coefficients themselves appear linearly. Let $x_L\leq x\leq x_U$, $y_L\leq y\leq y_U$, and let $z$ represent the product $xy$. The four quantities
\[
(x-x_L)(y-y_L),\quad (x_U-x)(y_U-y),\quad
(x_U-x)(y-y_L),\quad (x-x_L)(y_U-y)
\]
are non-negative. Expanding them and replacing $xy$ by $z$ gives the McCormick inequalities \citep{McCormick1976}
\begin{equation}
\equationname{McCormick inequalities for a bounded bilinear product}
\begin{aligned}
z&\geq x_Ly+y_Lx-x_Ly_L,\\
z&\geq x_Uy+y_Ux-x_Uy_U,\\
z&\leq x_Uy+y_Lx-x_Uy_L,\\
z&\leq x_Ly+y_Ux-x_Ly_U.
\end{aligned}
\label{eq:teaching_mccormick_relaxation}
\end{equation}
These inequalities are linear in $x$, $y$, and $z$, and every exact product satisfying the stated bounds obeys them. Their intersection is nevertheless larger than the exact nonlinear surface $z=xy$.

For example, if $0\leq x,y\leq1$, the inequalities reduce to $z\geq0$, $z\geq x+y-1$, $z\leq x$, and $z\leq y$. At $x=y=0.5$, they allow any $0\leq z\leq0.5$, while the exact product is only $z=0.25$. An objective value obtained at another permitted value of $z$ may therefore be a relaxation artifact rather than an attainable physical response. Tighter variable bounds can tighten the relaxation, but they do not generally recover the exact product relation. Process-network optimization uses this distinction when constructing bounds and interpreting solutions \citep{RubioCastro2013,YangNetwork2014}.

This example clarifies how convex relaxation differs from the projection used in this chapter. A relaxation enlarges a nonlinear feasible set to obtain a tractable optimization bound. The physical projection instead takes a statistical prediction and finds a nearby response satisfying the declared mass-conservation, non-negativity, separator, and inventory relations. The fixed plant topology is retained throughout both the surrogate-assisted and direct mechanistic searches.

\subsection{A direct mechanistic comparator}

\subsubsection{The full-state constrained optimization problem}

The direct route retains both the operating controls and the full reduced mechanistic state. It reconstructs $\chi(\vartheta,c_{\rm in},y)$ from that state and evaluates the same engineering objective used by the surrogate route. The finite-volume clarifier equations contain a positive feed-solids denominator at solver iterates. To keep this calculation defined during infeasible trial steps, an auxiliary variable $\widetilde X_F\ge X_F^{\min}$ supplies the denominator while feasibility requires $\widetilde X_F=t_X^\top c_{n_r}$.

A one-reactor, two-component, three-layer example illustrates the distinction between the direct and reduced formulations. The full state is $(c_S,c_X,s_1,s_2,s_3)$, which has five coordinates rather than the nine coordinates of $\chi_{\rm toy}$. With a simple conversion rate $kc_S$, the reactor equations are
\begin{align}
\equationname{Illustrative direct reactor equation for the consumed component}
0&=d(m_S-c_S)-kc_S,\\
\equationname{Illustrative direct reactor equation for the produced component}
0&=d(m_X-c_X)+kc_S.
\end{align}
The three clarifier equations are the layer balances defined earlier. The direct problem therefore imposes all five mechanistic equations, whereas the reduced physical projection retains only the aggregate reactor invariant and clarifier relations. Its three-layer envelope has two scalar inequalities, $s_1-s_2\le0$ and $s_2-s_3\le0$. Four layers would require $(s_1-s_2,s_1-s_3,s_2-s_4,s_3-s_4)\le0$. This pattern extends to the ten-layer case.

The auxiliary denominator has a numerical rather than physical role. During an infeasible trial, feed solids can approach zero while nonzero layer concentrations are still being adjusted. Division by the actual feed concentration would then be undefined. The bounded auxiliary $\widetilde X_F$ keeps the residual calculation finite. Because feasibility also imposes $\widetilde X_F=X_N$, it cannot be used to accept a zero-feed state as physically valid.

Let $f_{\rm sm}$ collect the $n_rF$ smooth reactor equations and the $L$ smooth clarifier-layer equations. Define
\begin{equation}
\equationname{Interior-layer envelope residuals for direct optimization}
g_{\rm env}(y)=\operatorname{col}
(s_1\mathbf1_{L-2}-s_{2:L-1},\ s_{2:L-1}-s_L\mathbf1_{L-2}).
\end{equation}
The direct nonlinear program is
\begin{equation}
\equationname{Direct mechanistic plant operating optimization problem}
\begin{aligned}
\min_{\vartheta,y,\widetilde X_F}\quad&J(\vartheta,\chi(\vartheta,c_{\rm in},y,\widetilde X_F),\varpi),\\
\text{subject to}\quad&f_{\rm sm}(\vartheta,c_{\rm in},y,\widetilde X_F)=0,\\
&\widetilde X_F=t_X^\top c_{n_r},\quad\widetilde X_F\ge X_F^{\min},\\
&\vartheta\in\mathcal V,\quad y\ge0,\quad g_{\rm env}(y)\le0,\\
&g_{{\rm eng},\rm red}(\vartheta,\chi)\le0.
\end{aligned}
\label{eq:plant_direct_optimization}
\end{equation}
The reduced engineering vector omits the feed-solids row because the auxiliary equality and lower bound already impose it. The auxiliary therefore changes only the numerical path taken through infeasible iterates, not the physical state admitted at feasibility.

\subsubsection{Smooth switches and continuation}

The exact kinetic and settling equations contain minima, maxima, positive parts, possible zero denominators, and a receiver switch. These operations are physically meaningful but introduce nondifferentiabilities or singular behavior that can hinder derivative-based optimization. The chemical-process treatment of \citet{BuzziFerraris2013} provides the broader context for derivative quality and numerical conditioning.

The exact identities
\begin{equation}
\equationname{Exact maximum and minimum expressed through absolute values}
\max(a,b)=\frac{a+b+|a-b|}{2},\qquad
\min(a,b)=\frac{a+b-|a-b|}{2}
\end{equation}
show that the nondifferentiability enters through $|a-b|$. Replacing $|r|$ by $\sqrt{r^2+\delta^2}$ produces a smooth approximation whenever $\delta>0$. The positive part uses the same substitution in $(v+|v|)/2$, whereas a potentially singular quotient uses a different regularization. For width $\delta=\epsilon S$,
\begin{align}
\equationname{Smooth approximation of a maximum}
\max_\delta(a,b)&=\tfrac12\left[a+b+\sqrt{(a-b)^2+\delta^2}\right],\\
\equationname{Smooth approximation of a minimum}
\min_\delta(a,b)&=\tfrac12\left[a+b-\sqrt{(a-b)^2+\delta^2}\right],\\
\equationname{Smooth approximation of the positive-part function}
[v]_{+,\delta}&=\tfrac12\left[v+\sqrt{v^2+\delta^2}\right],\\
\equationname{Regularized smooth approximation of a quotient}
\left(\frac ab\right)_\delta&=\frac{ab}{b^2+\delta^2}.
\label{eq:plant_smooth_operations}
\end{align}
For negative $v$, the positive-part approximation is evaluated as $\delta^2/[2(\sqrt{v^2+\delta^2}-v)]$ to avoid catastrophic cancellation. These approximations alter the equations near their switches, so the final selected controls are always reevaluated with the original nonsmoothed model.

For $a=2$, $b=1$, and $\delta=0.1$, the smooth maximum is 2.00249 and the smooth minimum is 0.99751, compared with exact values 2 and 1. For $v=-0.2$, the smooth positive part is 0.01180 rather than zero. At $b=0$, the regularized quotient remains defined and equals zero. At $a=1$, $b=0.1$, and $\delta=0.1$, it equals 5 rather than the exact quotient 10. These differences explain why the continuation gradually reduces the smoothing width and why an exact-model evaluation remains necessary.

The stable positive-part expression follows by rationalization,
\begin{align}
\equationname{Rationalization of the smoothed positive-part expression}
\frac{v+\sqrt{v^2+\delta^2}}2
&=\frac{(v+\sqrt{v^2+\delta^2})(\sqrt{v^2+\delta^2}-v)}
{2(\sqrt{v^2+\delta^2}-v)}\\
\equationname{Cancellation-resistant form of the smoothed positive part}
&=\frac{\delta^2}{2(\sqrt{v^2+\delta^2}-v)}.
\end{align}
The numerator reduces to $\delta^2$, eliminating subtraction of nearly equal terms.

The receiver switch is smoothed over a finite band around $X_t$. Subtracting the lower band edge $X_t-h$ and dividing by the full width $2h$ maps the transition interval to $[0,1]$. For receiver concentration $s$, the switch is zero below $X_t-h$, one above $X_t+h$, and
\begin{equation}
\equationname{Quintic receiver switch and its transition coordinate}
6\zeta^5-15\zeta^4+10\zeta^3,\qquad
\zeta=\frac{s-X_t+h}{2h}
\label{eq:plant_receiver_switch}
\end{equation}
inside the band. The quintic matches values and first two derivatives at both endpoints, producing a smooth transition between gravitational-flux rules.

For $X_t=3{,}000$ and $h=10$ g m$^{-3}$, concentrations 2{,}990, 3{,}000, and 3{,}010 correspond to $\zeta=0,0.5,1$ and switch values 0, 0.5, and 1. At 2{,}995, the switch is 0.10352; at 3{,}005, it is 0.89648. If $J_1$ denotes unrestricted gravitational flux and $J_2$ the receiver-limited flux, the smoothed interface uses $(1-\sigma)J_1+\sigma J_2$. The transition therefore modifies the flux law itself rather than merely smoothing a reported concentration.

Each smoothing scale $S$ is the development-set root mean square of the corresponding argument with a floor of one. Reactor scales are pooled across stages and include total oxidized nitrogen, fermentable substrate plus acetate, the hydrolysis denominator, polyphosphate and stored-carbon denominators, and remaining storage capacity. Clarifier scales include $s-f_{\rm ns}X_N$ and gravitational solids flux. The velocity scale is 250 m d$^{-1}$. These scales remain fixed during continuation.

The direct route uses three continuation stages,
\begin{equation}
\equationname{Smoothing and receiver-width continuation sequence}
(\epsilon,h)=(10^{-6},10),\ (10^{-7},3),\ (10^{-8},1),
\label{eq:plant_continuation_sequence}
\end{equation}
with $h$ in g TSS m$^{-3}$. Each stage initializes from the preceding solution. The sequence progressively sharpens the approximations while preserving a usable starting point. The final smooth candidate is still treated as a search output rather than as the final physical reference.

\subsubsection{Jacobian structure and globalization}

The direct mechanistic problem also benefits from its structured Jacobian. Entry $J_{jk}$ is the derivative of residual $f_j$ with respect to state coordinate $y_k$ and therefore describes how one balance responds to a local change in one state variable. For a hypothetical two-reactor soluble conversion with rate $kc_{i,S}$, soluble pass-through in the recycle gives
\begin{equation}
\equationname{Illustrative recycle-dependent soluble mixer concentration}
m_S=\frac{c_{{\rm in},S}+(r_I+r_R)c_{2,S}}{q_P}.
\end{equation}
The soluble residuals are $f_{1,S}=d_1(m_S-c_{1,S})-kc_{1,S}$ and $f_{2,S}=d_2(c_{1,S}-c_{2,S})-kc_{2,S}$. Their Jacobian block is
\begin{equation}
\equationname{Illustrative two-stage reactor Jacobian with recycle coupling}
\frac{\partial(f_{1,S},f_{2,S})}{\partial(c_{1,S},c_{2,S})}
=\begin{bmatrix}
-(d_1+k)&d_1(r_I+r_R)/q_P\\
d_2&-(d_2+k)
\end{bmatrix}.
\end{equation}
The upper-right entry arises entirely from recycle: changing the final-reactor concentration changes the mixer and therefore changes the first-reactor residual. Ignoring this term would incorrectly treat the mixer as independent of the downstream state. With $d_1=d_2=2$, $k=1$, $r_I=1$, $r_R=0.5$, and $q_P=2.5$, the block is $\begin{bmatrix}-3&1.2\\2&-3\end{bmatrix}$.

Finite differences provide an independent check on such derivatives. Increasing $c_{1,S}$ by 0.01 g m$^{-3}$ changes $f_{1,S}$ by $-0.03$ g m$^{-3}$ d$^{-1}$, giving derivative $-3$. Increasing $c_{2,S}$ by the same amount changes $m_S$ by 0.006 and $f_{1,S}$ by 0.012, giving 1.2. For a nonlinear residual, the corresponding central difference is
\begin{equation}
\equationname{Central finite-difference approximation of a Jacobian entry}
J_{jk}\approx\frac{f_j(y+he_k)-f_j(y-he_k)}{2h}.
\end{equation}
Near a non-negative boundary, a symmetric perturbation may leave the admissible state region. A second-order forward difference is then used,
\begin{equation}
\equationname{Second-order forward finite-difference approximation of a Jacobian entry}
J_{jk}\approx\frac{-3f_j(y)+4f_j(y+he_k)-f_j(y+2he_k)}{2h}.
\end{equation}
The perturbation $h$ retains the physical units of $y_k$. Scaling coordinates keeps its relative size comparable across different component bases.

For three clarifier layers at fixed feed conditions, each interface involves only neighboring layers. The layer Jacobian therefore has the tridiagonal pattern
\begin{equation}
\equationname{Tridiagonal sparsity pattern of the clarifier-layer Jacobian}
J_{\rm layer}=\begin{bmatrix}\star&\star&0\\\star&\star&\star\\0&\star&\star\end{bmatrix},
\end{equation}
where $\star$ indicates a potentially nonzero entry. Feed solids and outlet reconstruction add couplings to reactor coordinates. The zeros are consequences of the declared dependency structure rather than approximations introduced for computational convenience.

The complete direct residual remains sparse. A reactor balance depends primarily on its own component vector and the preceding stage, with additional recycle dependence on final-reactor and underflow quantities. A clarifier layer depends on neighboring layer concentrations, feed solids, and hydraulic controls. Inventory and reporting maps are linear aggregations. These local dependencies produce a structured Jacobian rather than an arbitrary dense response map.

Scaling is applied before Jacobian conditioning or residual magnitudes are assessed. The component basis includes oxygen, COD, nitrogen, phosphorus, solids, and charge equivalents, so an unscaled norm could be dominated by a large solids balance while concealing a nutrient-balance defect. Analytic derivatives or differentiated smooth equations supply the optimization Jacobian, while independent finite differences remain available for derivative verification. Sparse linear algebra exploits the local reactor and layer structure without dropping recycle couplings.

The direct search uses the interior-point filter method studied by \citet{WachterBiegler2006}. Its line search balances progress in objective value and constraint violation. This globalization mechanism governs acceptance of local steps during the search. As with the surrogate route, numerical search progress remains distinct from the final physical and stationarity audits.

\section{Methodology}

\subsection{Case-study domain and fixed parameters}

\subsubsection{Topology and free controls}

The case study contains five equal-volume biological reactors and ten clarifier layers. Reactors 1 and 2 are unaerated, while reactors 3--5 have independent aeration controls. The free-control vector is
\begin{equation}
\equationname{Seven freely varied controls of the connected-plant case}
\vartheta=(H,a_3,a_4,a_5,r_I,r_R,w)^\top\in\mathbb R^7.
\label{eq:plant_case_controls}
\end{equation}
For all five stages, $h_i=H/5$. The unaerated controls are fixed at $a_1=a_2=0$, the aerated-stage transfer coefficients are $k_La_i=47a_i$ d$^{-1}$, and $\xi_i=0$. These choices define the present case study; they are not structural requirements of the connected formulation.

\begin{table}[htbp]
\centering\small
\caption{Plant geometry and operating domain.}
\label{tab:plant_operating_domain}
\begin{tabular}{p{0.47\linewidth}p{0.40\linewidth}}
\toprule
Quantity & Value or interval\\
\midrule
Fresh-influent flow $Q_0$ & 10{,}000 m$^3$ d$^{-1}$\\
Reactors $n_r$, layers $L$, feed layer $\ell_F$ & 5, 10, 5\\
Reactor volume fraction $f_i^V$ & $1/5$ for every stage\\
Clarifier area and total depth & 1{,}500 m$^2$ and 4.0 m\\
Clarifier total and individual layer volumes & 6{,}000 m$^3$ and 600 m$^3$\\
Nominal total HRT $H$ & 6--36 h\\
Each of $a_3,a_4,a_5$ & 0--1 independently\\
Mixed-liquor recycle $r_I$ & 0--4\\
Return-sludge ratio $r_R$ & 0.25--1.25\\
Waste-sludge ratio $w$ & 0.001--0.05\\
Aeration reference $E_A^{\rm ref}$ & 423{,}000 m$^3$ d$^{-1}$\\
Wasted-solids reference $X_U^{\rm ref}$ & 15{,}000 g TSS m$^{-3}$\\
\bottomrule
\end{tabular}
\end{table}

The equal-volume reactors provide a compact unaerated-to-aerated treatment train. Each clarifier layer has depth 0.4 m. The operating bounds preserve positive overflow and underflow flow ratios throughout the domain. Together with the 20 fresh-influent components, the seven free controls produce 27 statistical inputs, 406 quadratic features, and 161 reduced-response coordinates. The mechanistic state contains 110 coordinates. The connected projection contains 76 equality rows and 12 physical inequality rows; after response non-negativity is included, the total number of inequality rows is 173.

Changing the number of reactors or free controls changes the statistical input and response dimensions. Changing the clarifier-layer count changes the mechanistic state and inventory mapping but does not add coordinates to the aggregate surrogate response. In either case, the fitted models, physical constraints, and predictive validation would need to be re-established. The layer discretization affects the mechanistic targets and the admissible inventory bounds even though the surrogate stores only aggregate clarifier inventory.

\subsubsection{Fresh-influent intervals}

Table~\ref{tab:plant_influent_intervals} defines the independent fresh-influent intervals used in the connected-plant experiment. The nominal influent is the componentwise midpoint of these intervals. Concentration bases and the composite reporting weights follow Table~\ref{tab:component_basis_composites}. The plant alkalinity interval is 1.6--5.2 mol charge equivalents m$^{-3}$.

\begin{table}[htbp]
\centering\small
\caption{Independent fresh-influent intervals for the connected plant.}\label{tab:plant_influent_intervals}
\begin{tabular}{lrrl}
\toprule
Component & Lower & Upper & Constituent basis\\
\midrule
$S_O$ & 0 & 0.5 & oxygen\\
$S_F$ & 20 & 180 & COD\\
$S_A$ & 5 & 80 & COD\\
$S_{NH4}$ & 12 & 55 & nitrogen\\
$S_{NO2}$ & 0 & 3 & nitrogen\\
$S_{NO3}$ & 0 & 8 & nitrogen\\
$S_{N2}$ & 0 & 2 & nitrogen\\
$S_{PO4}$ & 2 & 18 & phosphorus\\
$S_I$ & 10 & 90 & COD\\
$S_{ALK}$ & 1.6 & 5.2 & charge equivalents\\
$X_I$ & 20 & 120 & COD\\
$X_S$ & 60 & 280 & COD\\
$X_H$ & 15 & 100 & COD\\
$X_{PAO}$ & 5 & 60 & COD\\
$X_{PP}$ & 2 & 20 & phosphorus\\
$X_{PHA}$ & 1 & 30 & COD\\
$X_{AOB}$ & 0.5 & 8 & COD\\
$X_{NOB}$ & 0.5 & 8 & COD\\
$X_{MeP}$ & 0 & 12 & solids\\
$X_{MeOH}$ & 0 & 12 & solids\\
\bottomrule
\end{tabular}
\end{table}

All concentration values are expressed as grams of the stated constituent per m$^3$, except alkalinity. Numerically, g m$^{-3}$ equals mg L$^{-1}$. Independent interval sampling distributes candidates throughout the declared experimental domain without imposing correlations among the input coordinates.

\subsubsection{Stoichiometric, kinetic, and settling parameters}

Tables~\ref{tab:plant_stoichiometric_constants}--\ref{tab:plant_settling_constants} list the fixed parameters used in the connected-plant case. These quantities define the mechanistic reference model and are not fitted from the surrogate data. Each concentration-based parameter follows the constituent basis of its associated component. Symbols grouped on one row share a common numerical value. The two nitrifier populations remain distinct through their different growth rates and oxygen half-saturation constants.

\begin{table}[htbp]
\centering\small
\caption{Fixed yields, constituent contents, and conversion factors.}\label{tab:plant_stoichiometric_constants}
\begin{tabular}{p{0.38\linewidth}p{0.13\linewidth}p{0.38\linewidth}}
\toprule Parameter & Value & Unit or interpretation\\\midrule
$Y_H,Y_{PAO}$ & 0.625 & g biomass COD per g substrate COD\\
$Y_{PHA}$ & 0.20 & g stored COD per g phosphorus\\
$Y_{PO4}$ & 0.40 & g phosphorus per g stored COD\\
$Y_{AOB}$ & 0.18 & g biomass COD per g nitrogen\\
$Y_{NOB}$ & 0.08 & g biomass COD per g nitrogen\\
$f_{SI}$ & 0 & soluble inert fraction\\
$f_{XI}$ & 0.10 & particulate inert fraction\\
$i_{N,SI}$ & 0.01 & g nitrogen per g COD\\
$i_{N,SF},i_{N,XS}$ & 0 & g nitrogen per g COD\\
$i_{N,XI}$ & 0.02 & g nitrogen per g COD\\
$i_{N,BM}$ & 0.07 & g nitrogen per g biomass COD\\
$i_{P,SI},i_{P,SF},i_{P,XS}$ & 0 & g phosphorus per g COD\\
$i_{P,XI}$ & 0.01 & g phosphorus per g COD\\
$i_{P,BM}$ & 0.02 & g phosphorus per g biomass COD\\
$i_{P,MeP}$ & $1/4.87$ & g phosphorus per g metal-phosphate solids\\
$\gamma_{OH}$ & 3.45 & g metal-hydroxide solids per g phosphorus\\
\bottomrule
\end{tabular}
\end{table}

\begin{table}[htbp]
\centering\small
\caption{Fixed hydrolysis and heterotrophic rate and saturation constants.}
\label{tab:plant_kinetic_constants}
\begin{tabular}{p{0.48\linewidth}p{0.10\linewidth}p{0.31\linewidth}}
\toprule Parameter & Value & Unit\\\midrule
$K_H$ & 3.0 & d$^{-1}$\\
$\eta_{{\rm hyd},NO3},\eta_{{\rm hyd},NO2}$ & 0.6 & dimensionless\\
$\eta_{{\rm hyd},fe}$ & 0.4 & dimensionless\\
$K_{O,\rm hyd},K_{O,H},K_{O,PAO}$ & 0.20 & g oxygen m$^{-3}$\\
$K_{NO3,\rm hyd},K_{NO2,\rm hyd},K_{NOx,\rm hyd}$ & 0.5 & g nitrogen m$^{-3}$\\
$K_X$ & 0.10 & g substrate COD per g biomass COD\\
$\mu_H$ & 6.0 & d$^{-1}$\\
$q_{fe}$ & 3.0 & d$^{-1}$\\
$b_H$ & 0.40 & d$^{-1}$\\
$\eta_{H,NO3},\eta_{H,NO2}$ & 0.9 & dimensionless\\
$K_F,K_{fe},K_A$ & 4.0 & g COD m$^{-3}$\\
$K_{NO3,H},K_{NO2,H},K_{NOx,H}$ & 0.5 & g nitrogen m$^{-3}$\\
$K_{NH4,H},K_{NH4,PAO},K_{NH4,NOB}$ & 0.05 & g nitrogen m$^{-3}$\\
$K_{PO4,H},K_{PO4,PAO},K_{PO4,\rm nit}$ & 0.01 & g phosphorus m$^{-3}$\\
$K_{ALK,H},K_{ALK,PAO}$ & 0.10 & mol equivalents m$^{-3}$\\
\bottomrule
\end{tabular}
\end{table}

\begin{table}[htbp]
\centering\small
\caption{Fixed storage, nitrifier, and chemical rate and saturation constants.}
\label{tab:plant_storage_nitrifier_constants}
\begin{tabular}{p{0.48\linewidth}p{0.10\linewidth}p{0.31\linewidth}}
\toprule Parameter & Value & Unit\\\midrule
$q_{PHA}$ & 5.0 & g stored COD per g biomass COD per d\\
$q_{PP}$ & 0.60 & g phosphorus per g biomass COD per d\\
$\mu_{PAO}$ & 0.56 & d$^{-1}$\\
$\eta_{PAO,NO3}$ & 0.07 & dimensionless\\
$\eta_{PAO,NO2}$ & 0.90 & dimensionless\\
$b_{PAO},b_{PP},b_{PHA},b_{AOB}$ & 0.20 & d$^{-1}$\\
$K_{NO3,PAO},K_{NO2,PAO},K_{NOx,PAO}$ & 0.5 & g nitrogen m$^{-3}$\\
$K_{PS}$ & 0.20 & g phosphorus m$^{-3}$\\
$K_{PP}$ & 0.01 & g phosphorus per g biomass COD\\
$K_{\max}$ & 0.34 & g phosphorus per g biomass COD\\
$K_{IPP}$ & 0.02 & g phosphorus per g biomass COD\\
$K_{PHA}$ & 0.01 & g stored COD per g biomass COD\\
$\mu_{AOB}$ & 1.81 & d$^{-1}$\\
$\mu_{NOB}$ & 1.52 & d$^{-1}$\\
$b_{NOB}$ & 0.17 & d$^{-1}$\\
$K_{O,AOB}$ & 0.74 & g oxygen m$^{-3}$\\
$K_{O,NOB}$ & 1.75 & g oxygen m$^{-3}$\\
$K_{NH4,AOB},K_{NO2,NOB}$ & 0.5 & g nitrogen m$^{-3}$\\
$K_{ALK,\rm nit},K_{ALK,PRE},K_{ALK,\rm chem}$ & 0.50 & mol equivalents m$^{-3}$\\
$k_{PRE}$ & 1.0 & m$^3$ per g solids per d\\
$k_{RED}$ & 0.60 & d$^{-1}$\\
\bottomrule
\end{tabular}
\end{table}

\begin{table}[htbp]
\centering\small
\caption{Fixed settling and oxygen-transfer parameters.}
\label{tab:plant_settling_constants}
\begin{tabular}{p{0.46\linewidth}p{0.39\linewidth}}
\toprule Quantity & Value\\\midrule
Maximum settling velocity $v_{\max}$ & 250 m d$^{-1}$\\
Theoretical settling velocity $v_0$ & 474 m d$^{-1}$\\
Hindered-settling coefficient $r_h$ & 0.000576 m$^3$ per g TSS\\
Low-concentration coefficient $r_p$ & 0.00286 m$^3$ per g TSS\\
Nonsettleable fraction $f_{\rm ns}$ & 0.00228\\
Receiver threshold $X_t$ & 3{,}000 g TSS m$^{-3}$\\
Oxygen saturation $S_O^*$ & 8.5 g oxygen m$^{-3}$\\
Maximum aerated-stage $k_La$ & 47 d$^{-1}$\\
\bottomrule
\end{tabular}
\end{table}

\subsection{Mechanistic sampling, fitting, and prediction assessment}

\subsubsection{Simulation groups and accepted cases}

The connected input domain contains seven operating controls and twenty fresh-influent components. Strength-one Latin-hypercube sampling is used to distribute candidates across this 27-dimensional space. Each coordinate receives one point in every marginal stratum. Independent random permutations prevent the same ordering from being shared across coordinates, while open-unit jitter places each sampled value strictly inside its selected stratum. The coverage rationale follows the sampling design of \citet{McKay1979}.

The study attempts 8{,}000 development candidates and 2{,}000 holdout candidates using separate seeded random streams. Rejected candidates remain part of the accounting and are not replaced. A two-start physical acceptance procedure determines which attempted cases become statistical states. Acceptance and rejection counts are reported separately for the development and holdout designs.

\begin{table}[htbp]
\centering\small
\caption{Randomization and fitting settings specified before assessment.}
\label{tab:plant_design_settings}
\begin{tabular}{p{0.54\linewidth}p{0.32\linewidth}}
\toprule Setting & Value\\\midrule
Development Latin-hypercube seed & 100042\\
Holdout Latin-hypercube seed & 100043\\
Influent-scenario Latin-hypercube seed & 314159\\
Grouped-fold permutation seed & 271828\\
Number of fitting folds & 5\\
Raw-response and logarithmic ridge grids & $\{10^{-8},10^{-7},\ldots,10^2\}$\\
Standard-deviation divisor & fitting-set size\\
Near-constant coordinate cutoff & $10^{-12}\max(1,\max_i|v_i|)$\\
Low-overflow group & at or below development 25th percentile\\
High-overflow group & at or above development 90th percentile\\
\bottomrule
\end{tabular}
\end{table}

A four-stratum example illustrates the Latin-hypercube mapping. Each unit-interval stratum has width $1/4$. Selecting stratum two with within-stratum fraction 0.6 gives $(2+0.6)/4=0.65$. If the physical coordinate spans 10 to 20, this maps to $10+0.65(20-10)=16.5$. Independent permutations of the remaining coordinates retain marginal coverage while avoiding a shared coordinate ordering.

The exact jitter rule uses the same idea at much finer resolution. Dividing an unsigned 64-bit random integer $U$ by $2^{11}$ and taking the floor gives an integer from zero to $2^{53}-1$. Adding 0.5 selects the midpoint of that fine stratum, and division by $2^{53}$ maps it to the open unit interval. Thus
\begin{equation}
\equationname{Conversion of a uniform integer to open-unit jitter}
\zeta_U=\frac{\lfloor U/2^{11}\rfloor+0.5}{2^{53}}\in(0,1).
\end{equation}
In exact arithmetic, the extreme values are $2^{-54}$ and $1-2^{-54}$. The procedure therefore excludes exact stratum boundaries, although endpoint checks remain necessary after finite-precision mapping to physical variables.

Each coordinate uses a descending Fisher--Yates permutation of its strata. Accepted mechanistic states retain their original candidate order before the grouped-fold permutation is applied. The fitting and prediction assessments therefore use only states that have already passed the mechanistic integration, balance, branch, and stability checks, while rejection counts preserve the outcomes of the attempted design.

\subsubsection{Two initial states and the nonsmoothed steady solve}

Mechanistic generation uses two deliberately different initial conditions. The first assigns $c_{\rm in}$ to every reactor and fresh-influent TSS to every clarifier layer. The second retains the fresh soluble concentrations but multiplies the particulate concentrations in reactors 1--5 by $(1.5,2,2.5,3,3.5)$. Its final-reactor TSS then sets the initial clarifier profile through
\begin{equation}
\equationname{Ten-layer initial solids-profile multipliers}
(0.002,0.005,0.01,0.03,0.10,0.50,1.25,2.25,3.25,4.00)
\label{eq:plant_initial_layer_multipliers}
\end{equation}
from the surface to the hopper. The two starts probe sensitivity to distinct biomass and storage states. Agreement between their converged physical states is part of the acceptance procedure.

A logarithmic transformation keeps integrated coordinates positive. For a positive native-coordinate value $y_f$, define $\eta_f=\log y_f$. Because $y_f=\exp(\eta_f)$,
\[
\begin{aligned}
y_f&=\exp(\eta_f),\\
\frac{\mathrm d y_f}{\mathrm dt}
&=\exp(\eta_f)\frac{\mathrm d\eta_f}{\mathrm dt}
=y_f\frac{\mathrm d\eta_f}{\mathrm dt}.
\end{aligned}
\]
The original balance requires $\mathrm d y_f/\mathrm dt=f_f(y)$, so
$\mathrm d\eta_f/\mathrm dt=f_f(y)/y_f$. The transformation therefore changes only the numerical coordinates, not the underlying physical balance.

For example, if a concentration is 2 native units and its balance rate is $-0.4$ native units per day, the transformed rate is $-0.2$ d$^{-1}$. A half-day illustrative Euler step in $\eta$ gives
\[
\eta_{new}=\log2+0.5(-0.2),\qquad
y_{new}=\exp(\eta_{new})=2\exp(-0.1)\approx1.8097.
\]
The actual integration uses the backward differentiation method with the numerical tolerances specified below; this arithmetic is only intended to show the effect of the positive-state transformation.

Before the logarithmic transformation, any initial coordinate below $10^{-8}$ is raised to $10^{-8}$ in its native physical unit. This safeguard applies only to initialization and is distinct from clipping a final solution or prediction.

The integration horizon depends on wasting because low wasting can lead to long solids residence times. It is defined as
\begin{equation}
\equationname{Waste-dependent horizon for dynamic relaxation}
T=\max\{400,50/\max(w,0.001)\}\ \mathrm{d}.
\label{eq:plant_integration_horizon}
\end{equation}
At $w=0.01$, this gives 5{,}000 d; at $w=0.20$, it gives the 400-d minimum; and at $w=0$, the internal floor gives 50{,}000 d. The values zero and 0.20 are used only to illustrate the branches of the rule and lie outside the admitted case-study wasting interval.

Dynamic relaxation uses the backward differentiation method for stiff systems \citep{CurtissHirschfelder1952}. The integration may stop before $T$ when the scaled residual falls below $10^{-9}$. Its purpose is to produce a physically informed initial state for the algebraic steady calculation, consistent with the continuation rationale of \citet{KelleyKeyes1998} and the activated-sludge procedure of \citet{Selerio2026}. A bounded least-squares polishing calculation is attempted only if the integrated endpoint fails the physical audit. Solver convergence and scientific acceptance are therefore kept separate.

\begin{table}[htbp]
\centering\small
\caption{Steady-state generation and physical acceptance settings.}\label{tab:plant_generation_settings}
\begin{tabular}{p{0.61\linewidth}p{0.25\linewidth}}
\toprule Setting & Value\\\midrule
Integration relative tolerance & $10^{-7}$\\
Transformed integration absolute tolerance & $10^{-9}$\\
Maximum integration step & $T/100$\\
Polishing evaluation limit & 5{,}000\\
Polishing step, function, and gradient tolerances & $10^{-9}$ each\\
Term-scaled balance tolerance & $10^{-6}$\\
State non-negativity tolerance & $10^{-10}$\\
Reaction-rate non-negativity tolerance & $10^{-12}$\\
Two-start scaled state agreement & $10^{-6}$\\
Rightmost reduced-state Jacobian eigenvalue & at most $-10^{-8}$ d$^{-1}$\\
Eigenvalue agreement between two step sizes & $10^{-6}$ d$^{-1}$\\
Matrix mass conservation and rank audit tolerance & $10^{-10}$\\
\bottomrule
\end{tabular}
\end{table}

The balance audit scales each residual relative to the physical terms that form it. Suppose incoming, consumption, and outgoing contributions are $100$, $-60$, and $-39.99999$ in one consistent rate unit. Their net residual is
\[
100-60-39.99999=0.00001.
\]
The largest term magnitude is 100, so the scaled residual is $10^{-7}$. This comparison evaluates the size of the cancellation error relative to the terms being canceled. A floor of one prevents division by an arbitrarily small scale when all terms are near zero. For terms $b_j$,
\begin{equation}
\equationname{Scaled component-equation acceptance residual}
r_{\rm bal}=\frac{|\sum_jb_j|}{\max(1,\max_j|b_j|)}.
\label{eq:plant_acceptance_balance}
\end{equation}
The residual is evaluated separately for each balance because the component equations use different constituent bases and cannot be meaningfully summed into one unscaled physical quantity.

Reactor generation scales repeat $\max(1,c_{{\rm in},f}^U)$ for each component, while layer scales use $\max(1,X_N)$. Agreement between the two starts is tested in these scaled coordinates.

Local stability is assessed from the dynamic Jacobian at the candidate equilibrium. For a small perturbation, $\dot{\delta y}=J_{dyn}\delta y$. In a simple independent two-coordinate example,
\[
\begin{bmatrix}\dot{\delta y}_1\\\dot{\delta y}_2\end{bmatrix}
=\begin{bmatrix}-0.2&0\\0&0.05\end{bmatrix}
\begin{bmatrix}\delta y_1\\\delta y_2\end{bmatrix}.
\]
The corresponding perturbations evolve as $\delta y_1(0)\exp(-0.2t)$ and $\delta y_2(0)\exp(0.05t)$. One decays while the other grows. Because the rightmost eigenvalue is positive, the equilibrium fails the local stability criterion.

For the coupled plant, each eigenvalue describes the local growth or decay of one state-space mode. The acceptance test uses the largest real part across all modes. The Jacobian is estimated with normalized finite-difference steps $\epsilon_{\rm mach}^{1/3}$ and $2\epsilon_{\rm mach}^{1/3}$. Central differences are used away from the non-negative boundary, and second-order forward differences are used near it. Both rightmost-eigenvalue estimates must satisfy the upper bound and agree within the stated tolerance.

Both initial states must produce finite solutions, satisfy concentration and rate non-negativity, pass the component and invariant mass-conservation audits, remain inside the operating domain and clarifier envelope, and satisfy the layer-flux residual tolerance. The two starts must also return identical branch-classification tuples, which record the active kinetic and settling branches. A state that satisfies one residual test but fails another scientific or physical condition is not admitted as a training target.

\subsubsection{Fold-local fitting and two separately selected penalties}

Accepted states are partitioned by complete plant case: every component and location belonging to one simulated state remains in the same fold. This grouping prevents any part of one mechanistic plant response from appearing simultaneously in fitting and fold assessment. Fold sizes differ by at most one. Input centers, input scales, quadratic-feature scales, and response scales are recomputed within each fitting fold. The holdout set contributes to none of these quantities. This separation follows the leakage concerns described by \citet{Kaufman2012}.

The one-standard-error rule can be illustrated with five hypothetical fold scores $0.40$, $0.42$, $0.38$, $0.41$, and $0.39$. Their mean is 0.40. The sum of squared deviations is
\[
0^2+0.02^2+(-0.02)^2+0.01^2+(-0.01)^2=0.001.
\]
The sample standard deviation is $\sqrt{0.001/4}\approx0.01581$, and the standard error is $0.01581/\sqrt5\approx0.00707$. The admissible mean-score threshold is therefore $0.40707$.

If a stronger penalty has mean score 0.406, it lies within this threshold and remains eligible. A still stronger penalty with mean score 0.430 does not. The selected penalty is the largest candidate within one standard error of the best mean score. The numerical threshold itself is not a penalty value.

For the raw connected response, each candidate penalty is scored by standardized root mean square error across the complete response vector. The largest penalty within one standard error of the minimum mean score is retained, favoring stronger regularization when the assessment does not clearly distinguish it from the empirical minimum \citep{Hastie2009}. Projected errors are not used to choose the raw-response penalty. The logarithmic overflow model uses the same grouped folds, the same candidate grid, and the same selection rule, but it is scored on the log-transformed endpoint target. Its penalty is therefore selected independently.

The augmented ridge equations are solved with column-pivoted orthogonal-triangular factorization and audited with singular-value decomposition. A coordinate whose standard deviation falls below the near-constant cutoff in Table~\ref{tab:plant_design_settings} is rejected rather than assigned an arbitrary scale. Positive ridge penalties and the unpenalized intercept give a unique coefficient solution for the stated fitting design. The selected penalties and fitted coefficient matrices are estimation outputs, not theoretical constants.

After selection, both models are refitted once using all accepted development states. Grouped out-of-fold predictions are retained for trust calibration and projection assessment. For each out-of-fold state, the overflow equality uses the logarithmic model trained without that state rather than the final model that has seen it. The final full-development fits are used only for holdout prediction and operating optimization.

The logarithmic target must be strictly positive. Predicted logarithms and their exponentiated values must also remain finite. The low- and high-overflow strata are defined by the development percentiles in Table~\ref{tab:plant_design_settings}. Separate assessment of these strata can expose response-dependent error that would be hidden by a central-range average. The imbalanced-regression analysis of \citet{Shahbazi2026} motivates explicit attention to sparse or extreme response regions.

\subsubsection{Scientific admission and projection audits}

Scientific admission determines whether the fitted connected surrogate is acceptable for operating assessment under the specified criteria. Grouped out-of-fold predictions provide the relevant evidence because each development state is predicted by a model that excluded it from fitting.

Admission requires finite and independently audited out-of-fold projections for every accepted development state. The complete raw-response normalized root mean square error and the raw clarifier-inventory normalized error must each be below one. Because both metrics use development response scales, the threshold requires aggregate prediction error below one fitted scale unit. This condition is specific to the case-study admission rule.

Every projected state is checked for equality feasibility, physical inequalities, non-negativity, dual feasibility, stationarity, and complementarity. The overflow-closure residual is reported separately. Failure of an admission criterion remains a failure; it does not trigger an unreported change to the fit, the selected penalty, or the projection formulation.

\begin{table}[htbp]
\centering\small
\caption{Projection solve and independent audit settings.}
\label{tab:plant_projection_settings}
\begin{tabular}{p{0.62\linewidth}p{0.24\linewidth}}
\toprule Setting & Value\\\midrule
Solver absolute and relative tolerances & $10^{-12}$ each\\
Initial iteration limit & 100{,}000\\
Polishing & enabled\\
Retry values of algorithm parameter $\rho_{\rm alg}$ & 0.01 and 10\\
Retry iteration limit & 200{,}000 each\\
Adaptive $\rho_{\rm alg}$ during retries & disabled\\
Scaled equality and inequality tolerances & $10^{-8}$ each\\
Dual, stationarity, and complementarity tolerances & $10^{-8}$ each\\
Physical non-negativity tolerance & $10^{-10}$\\
Active-set identification tolerance & $10^{-7}$\\
\bottomrule
\end{tabular}
\end{table}

\subsubsection{Prediction metrics and their denominators}

Prediction metrics are reported with their denominators made explicit because different denominators answer different questions. Consider two reference values 2 and 6 and predictions 3 and 4 for one concentration coordinate. The residuals are 1 and $-2$, giving
\[
\begin{aligned}
\text{mean squared error}&=\frac{1^2+(-2)^2}{2}=2.5,\\
\text{root mean square error}&=\sqrt{2.5}\approx1.5811,\\
\text{mean absolute error}&=\frac{|1|+|-2|}{2}=1.5,\\
\text{bias}&=\frac{1+(-2)}{2}=-0.5.
\end{aligned}
\]
The mean squared error has squared concentration units, whereas the other three retain the native concentration unit. Bias preserves the sign of the error and can therefore remain small when positive and negative errors cancel.

For the same values, the reference mean is 4. The prediction squared error is 5, whereas always predicting the reference mean gives squared error 8. The coefficient of determination is therefore $1-5/8=0.375$. Its denominator measures variation in the assessment reference and is unrelated to the training standardization or any engineering concentration limit.

For response coordinate $f$ and state $i$, define $a_{if}=p_{if}-t_{if}$, where $p$ is predicted and $t$ is mechanistic. Physical-unit root mean square error, mean absolute error, and bias are
\begin{equation}
\equationname{Coordinatewise root error, absolute error, and bias}
\left(\operatorname{mean}_i a_{if}^2\right)^{1/2},\qquad
\operatorname{mean}_i|a_{if}|,\qquad
\operatorname{mean}_i a_{if}.
\end{equation}
Normalized joint-response metrics replace $a_{if}$ by $a_{if}/(D_\chi)_{ff}$ and explicitly state whether errors are pooled across coordinates. Coordinatewise coefficients of determination are
\begin{equation}
\equationname{Coefficient of determination for a plant-response coordinate}
R_f^2=1-\frac{\sum_i(p_{if}-t_{if})^2}{\sum_i(t_{if}-\bar t_f)^2}.
\end{equation}
If the assessment target is constant, the denominator is zero and the score is recorded as undefined. Averaging coordinatewise $R^2$ values is not equivalent to computing one $R^2$ after pooling coordinates.

Normalization also changes the numerical magnitude of an error summary. In the preceding example, if the training response scale is two native units, the residuals become 0.5 and $-1$, giving standardized RMSE $\sqrt{0.625}\approx0.7906$. If a descriptive holdout report instead divides by the observed range $6-2=4$, the residuals become 0.25 and $-0.5$, giving RMSE approximately 0.3953. The predictions have not changed; only the denominator has. Training-scale admission scores and holdout-range reporting scores must therefore be identified separately.

Location-quality summaries retain the mixer, five reactors, overflow, and underflow because joint constraints couple their errors across the treatment train. Outlet-flow coordinates are first divided by their positive flow ratios to recover concentrations, after which the quality map is applied. For location $\ell$ and composite $j$, define
\begin{equation}
\equationname{Holdout-range normalization of a location-quality error}
r_{\ell j}=\max_it_{i\ell j}-\min_it_{i\ell j},\qquad
e_{i\ell j}=\frac{p_{i\ell j}-t_{i\ell j}}{r_{\ell j}}.
\end{equation}
For positive ranges, the aggregate error pools the $8\times4$ location-composite coordinates. A location score pools the four composites, while a composite score pools the eight locations. Displayed $R^2$ summaries average the corresponding coordinatewise values. These holdout ranges are used only for reporting; they do not enter fitting, admission, objective normalization, or trust calibration.

The separate overflow model is assessed both on the concentration and logarithmic scales. Concentration-scale reporting includes root mean square error, mean absolute error, bias, normalized error, and $R^2$, while logarithmic reporting includes root mean square error and bias. Positivity from Equation~\eqref{eq:plant_log_prediction} is a structural property of the model, not a substitute for accuracy assessment.

Exploratory simulation informed the eventual response definition, empirical closure, and acceptance rules. Development data then supplied all fitted transformations, selected penalties, overflow fitting, and trust thresholds. The independent holdout stream evaluates the resulting frozen procedure. Operating assessment finally uses new mechanistic calculations at the controls selected by the two search routes. Fitted coefficients, preprocessing, objective priorities, trust limits, safeguards, and numerical settings remain fixed throughout these assessments \citep{Kaufman2012}.

\subsection{Operating searches and convergence evidence}

\subsubsection{A common baseline and an explicit basin-sensitivity experiment}

The operating controls are normalized to the unit interval before search so that variables with different physical units can use comparable numerical step sizes. For a control spanning 10 to 30, a value of 15 maps to $(15-10)/(30-10)=0.25$, the midpoint 20 maps to 0.5, and the upper bound maps to one. The inverse transformation is
\[
\vartheta_j=\vartheta_j^L+(\vartheta_j^U-\vartheta_j^L)\eta_j.
\]
This normalization changes only the numerical search coordinates; it does not give all controls equal importance in the engineering objective.

The baseline search for both routes begins at the center of the normalized operating box,
\begin{equation}
\equationname{Operating-control normalization to the unit interval}
\eta_j=\frac{\vartheta_j-\vartheta_j^L}{\vartheta_j^U-\vartheta_j^L}\in[0,1].
\label{eq:plant_normalized_controls}
\end{equation}
The nominal influent is the midpoint of the fresh-influent intervals. Because the baseline uses only the center start, its outcome represents one local search basin.

A separate planned multistart experiment evaluates basin sensitivity. For seven controls, the center is the vector with all entries equal to 0.5. Moving the first coordinate up or down by one-quarter of its range gives
\[
\begin{aligned}
\eta_{\mathrm{first},+}&=(0.75,0.5,0.5,0.5,0.5,0.5,0.5)^\top,\\
\eta_{\mathrm{first},-}&=(0.25,0.5,0.5,0.5,0.5,0.5,0.5)^\top,\\
e_1&=(1,0,0,0,0,0,0)^\top.
\end{aligned}
\]
Repeating this displacement for each coordinate and both signs gives 14 displaced starts. Together with the center, the design contains 15 common starting vectors,
\begin{equation}
\equationname{Common center and displaced starts for the multistart comparison}
\eta^{(0)}=\tfrac12\mathbf1_7,\qquad
\eta^{(j,\pm)}=\tfrac12\mathbf1_7\pm\tfrac14e_j,
\qquad j=1,\ldots,7.
\label{eq:plant_multistart_design}
\end{equation}
Each start receives the same route-specific budget and audit procedure. A direct start obtains its mechanistic initial state from an independently accepted solve at the fixed starting controls. Failed initial states remain failures. The best independently validated feasible candidate is retained, while all starting-point outcomes and total search work remain part of the record. The single-center result remains the reported baseline; the 15-start design is a distinct sensitivity analysis.

This distinction matters when interpreting computation time. A method that explores fewer basins may finish sooner while missing a better candidate. The expensive-function benchmarking study of \citet{Bliek2023} therefore motivates reporting objective quality, evaluation work, and elapsed time as separate dimensions.

\subsubsection{Surrogate search and derivative fallback}

The initial surrogate outer search uses sequential least squares programming, following the constrained optimization framework of \citet{Kraft1988}. Every trial point receives a fully audited joint projection. The active-set derivative chain is used only when the projection active set satisfies the prescribed rank, conditioning, multiplier, and independent derivative checks. If the active set changes, the projection and its derivatives are reconstructed for the new trial.

If the derivative chain cannot be validated, a deterministic value-only quadratic-approximation method continues the search within the same normalized control bounds \citep{Ragonneau2022}. The fallback still uses the exact audited projection at every trial; it does not replace the lower problem with an approximate projection. Finite differences are reserved for independent derivative audits rather than used automatically as unqualified gradient substitutes. Before a candidate is retained, its best feasible visited point is projected again from an independent initialization.

\begin{table}[htbp]
\centering\small
\caption{Surrogate search, sensitivity, and finite-poll settings.}\label{tab:plant_surrogate_search_settings}
\begin{tabular}{p{0.64\linewidth}p{0.23\linewidth}}
\toprule Setting & Value\\\midrule
Initial outer iteration limit & 250\\
Initial outer function tolerance & $10^{-10}$\\
Active slack tolerance & $10^{-7}$\\
Required active inequality multiplier & greater than $10^{-8}$\\
Domain-aware active-set perturbation & $10^{-6}$\\
Stationary-system condition number times machine precision & at most $10^{-8}$\\
Sensitivity residual and scaled state reproduction & $10^{-8}$ each\\
Independent derivative steps & $10^{-6}$ and $5\times10^{-7}$\\
Derivative relative-error tolerance & $10^{-5}$\\
Upper feasibility and stationarity tolerance & $10^{-6}$ each\\
Value-only fallback iteration and evaluation limits & 250 each\\
Initial and final fallback trust radii & 0.25 and $10^{-6}$\\
Fallback feasibility tolerance & $10^{-6}$\\
Coarse and fine poll radii & $10^{-3}$ and $10^{-4}$\\
Coarse and fine direction counts & 14 and 106\\
Absolute and relative sufficient-decrease tolerances & $10^{-8}$ each\\
Direction-rank tolerance & $10^{-10}$\\
Certification evaluation limit & 10{,}000\\
Ray growth factor and maximum probes & 2 and 16\\
Objective relative tie tolerance & $10^{-10}$\\
\bottomrule
\end{tabular}
\end{table}

\subsubsection{First-order evidence and finite-resolution polling}

A retained surrogate candidate receives first-order stationarity qualification only if the derivative chain, lower projection certificate, and upper constrained-stationarity audit all pass. This qualification applies specifically to the declared surrogate optimization problem at that candidate.

When derivative evidence is unavailable, finite polling provides a separate, resolution-limited check. The need for coupled directions can be seen in a two-control example with constraint $x+y=1$, non-negative variables, and objective $J=x$. At $(0.5,0.5)$, any nonzero move along either coordinate axis violates the equality. A diagonal step along $(-1,1)^\top/\sqrt2$ instead preserves feasibility while reducing $x$. With step radius 0.1,
\[
\begin{aligned}
x_{trial}&=0.5-0.1/\sqrt2\approx0.4293,\\
y_{trial}&=0.5+0.1/\sqrt2\approx0.5707,\\
x_{trial}+y_{trial}&=1.
\end{aligned}
\]
The objective decreases by approximately 0.0707. An axis-only poll would therefore miss a feasible improving direction.

In seven controls, the signed coordinate directions contribute 14 vectors. The $7(6)/2=21$ unordered coordinate pairs each admit four sign combinations, giving 84 signed pairwise directions. Eight regular-simplex directions supplement them. The fine poll therefore contains $14+84+8=106$ directions. The simplex directions distribute around the origin rather than following only axes or coordinate pairs.

If first-order qualification is unavailable, complete feasible-direction polls are performed at normalized radii $10^{-3}$ and $10^{-4}$. The coarse poll contains the 14 signed coordinate directions. The fine poll adds the 84 pairwise diagonals and eight Helmert-based simplex directions. Coupled directions are included because active constraints can make improving feasible directions invisible to coordinate-only exploration.

Each complete poll must contain at least one feasible nonzero displacement. A rank calculation is retained as a diagnostic because active constraints can reduce the dimension of the feasible cone. If a sufficient improvement is found, a new complete poll is performed from the improved point. A fine-radius improvement also triggers revalidation at the coarse radius. Up to 16 search-only probes may then extend along the successful ray by factors of two. These probes accelerate exploration but provide no convergence certificate.

Failure to find sufficient decrease in both complete revalidated polls provides evidence only for the tested directions and radii. An incomplete poll, an evaluation failure, or exhaustion of the certification budget leaves stationarity unresolved. This interpretation follows the role of feasible polling in direct-search analysis \citep{AudetDennis2006}.

\subsubsection{Direct continuation, recovery, and endpoint checks}

The direct route uses the interior-point filter method of \citet{WachterBiegler2006} through the three smoothing stages defined earlier. Its endpoint audit independently checks the smooth reactor and clarifier equations, inequalities, non-negativity, multipliers, and stationarity. Solver termination status is therefore not treated as sufficient evidence of a valid direct candidate.

If the initial direct search fails, one recovery attempt may begin from the certified surrogate controls for the same influent and repeat the complete smoothing sequence. Recovery is not attempted after a successful initial direct search or from an uncertified surrogate candidate. Because the recovery start depends on the surrogate decision, it is a dependent rescue calculation rather than an independent multistart search. Its status and computational cost are reported separately.

\begin{table}[htbp]
\centering\small
\caption{Direct continuation and exact-evaluation qualification settings.}\label{tab:plant_direct_settings}
\begin{tabular}{p{0.64\linewidth}p{0.23\linewidth}}
\toprule Setting & Value\\\midrule
Interior-point iteration limit per smoothing stage & 2{,}500\\
Convergence and constraint tolerances & $10^{-8}$ each\\
Dual and complementarity tolerances & $10^{-6}$ each\\
Endpoint equality acceptance & $10^{-8}$\\
Endpoint inequality and stationarity acceptance & $10^{-6}$ each\\
Endpoint active-set tolerance & $10^{-7}$\\
Receiver-switch ambiguity distance from $X_t$ & 1 g TSS m$^{-3}$\\
Other branch ambiguity distances & $10^{-6}S$\\
Two-start agreement under both prescribed state scales & $10^{-6}$\\
Engineering feasibility tolerance after exact evaluation & $10^{-6}$\\
\bottomrule
\end{tabular}
\end{table}

A validated feasible candidate may still be retained when its search budget is exhausted or its stationarity qualification remains unresolved. In that case, the unresolved status remains visible. A case with no validated feasible candidate remains a search failure. A favorable solver flag alone is insufficient for either optimization route.

\subsection{Independent mechanistic verification and fair comparison}

The surrogate and direct searches use different internal response calculations, but their selected decisions must ultimately be compared on one common physical basis. Figure~\ref{fig:plant_route_verification} separates four stages: route-specific search, native candidate audit, independent original-model evaluation, and final comparison eligibility. Each stage retains its own status rather than being collapsed into one success indicator.

\begin{figure}[htbp]
\centering
\begin{tikzpicture}[>=Latex,
block/.style={draw=black!75,rounded corners,fill=cyan!6,align=center,font=\footnotesize,inner sep=4pt},
audit/.style={block,fill=black!5},
failure/.style={block,dashed,fill=white},
flow/.style={->,thick,draw=black!80},
failflow/.style={->,dashed,draw=black!70},
lab/.style={font=\scriptsize,fill=white,inner sep=2pt}]
\node[block,text width=12.9cm] (common) at (0,0)
{Common influent and control domain\\Fixed priorities and engineering safeguards};
\node[block,text width=6.0cm] (surrogate) at (-3.7,-1.85)
{Surrogate search\\Prediction and joint projection\\Trust screening};
\node[block,text width=6.0cm] (direct) at (3.7,-1.85)
{Direct mechanistic search\\Full mechanistic states\\Smoothing continuation};
\node[audit,text width=6.0cm] (snative) at (-3.7,-3.85)
{Surrogate candidate audit\\Native feasibility\\Convergence status retained};
\node[audit,text width=6.0cm] (mnative) at (3.7,-3.85)
{Direct candidate audit\\Native feasibility\\Convergence status retained};
\node[failure,text width=4.0cm] (sfail) at (-5.0,-6.3)
{No feasible candidate\\Retain route failure};
\node[failure,text width=4.0cm] (mfail) at (5.0,-6.3)
{No feasible candidate\\Retain route failure};
\node[block,text width=4.5cm] (evaluate) at (0,-6.3)
{Independent evaluation\\of each retained decision\\Fixed selected controls\\Two initial states\\Original nonsmooth model};
\node[audit,text width=6.6cm] (valid) at (0,-9.05)
{Original-model validity\\Balances and non-negativity\\Layer equations and branches\\Reduced-state stability};
\node[failure,text width=3.0cm] (invalid) at (5.7,-9.05)
{Invalid state\\Retain failure};
\node[audit,text width=6.6cm] (engineering) at (0,-11.4)
{Engineering eligibility\\Underflow solids limit\\Feed-solids and external-loss floors};
\node[failure,text width=3.0cm] (engfail) at (5.7,-11.4)
{Valid state\\Safeguard fails\\Retain diagnostic};
\node[block,text width=9.3cm] (comparison) at (0,-13.5)
{Common objective comparison\\Both decisions valid and engineering eligible};
\draw[flow] (common.south west)--(surrogate.north);
\draw[flow] (common.south east)--(direct.north);
\draw[flow] (surrogate)--(snative);
\draw[flow] (direct)--(mnative);
\draw[flow] (snative.south east)--(evaluate.north west) node[midway,right,lab] {Retain};
\draw[flow] (mnative.south west)--(evaluate.north east) node[midway,left,lab] {Retain};
\draw[failflow] (snative.south west)--(sfail.north) node[midway,left,lab] {Fail};
\draw[failflow] (mnative.south east)--(mfail.north) node[midway,right,lab] {Fail};
\draw[flow] (evaluate)--(valid);
\draw[flow] (valid)--(engineering) node[midway,left,lab] {Valid};
\draw[failflow] (valid)--(invalid) node[midway,above,lab] {Fail};
\draw[flow] (engineering)--(comparison) node[midway,left,lab] {Eligible};
\draw[failflow] (engineering)--(engfail) node[midway,above,lab] {Fail};
\end{tikzpicture}
\caption{Conceptual paired-search and verification workflow. Each retained control vector receives its own independent original-model evaluation. The surrogate and smooth-direct native responses do not supply the reference treatment outcome. Original-model validity and engineering eligibility are distinct checks. Only eligible pairs enter the common-objective comparison, while failed or qualified cases remain in the accounting.}
\label{fig:plant_route_verification}
\end{figure}

The independent evaluation stage supplies treatment outcomes at the selected controls without feeding information back into either fitted search model. A mechanistically valid response may still violate a retained engineering safeguard and remain useful as a diagnostic. Separating these outcomes prevents numerical convergence, physical validity, and decision eligibility from being treated as the same property.

\subsubsection{Evaluation at fixed selected controls}

Every available selected decision is reevaluated with its controls and influent held fixed. Both prescribed initial states are solved using the original nonsmoothed reactor and ten-layer clarifier equations. The independently converged state supplies the reference treatment outcome for both routes. This role of an original-model check is consistent with the validation requirements discussed by \citet{BhosekarIerapetritou2018}.

Exact evaluation uses the physical checks in Table~\ref{tab:plant_generation_settings}. Agreement between the two starts is tested under both the generation scales and fixed development-derived state scales. The fixed scales are coordinatewise development-set median absolute deviations with a floor of one. A valid evaluation must pass mass conservation, non-negativity, layer residuals, the operating domain, the endpoint envelope, and the reduced-state local stability test.

Branch tuples must agree unless one evaluation start lies inside a prescribed ambiguity band. These bands cover the clarifier receiver threshold, zero storage capacity, zero settling offset, zero or maximum settling velocity, and equality of adjacent gravitational fluxes. An ambiguity flag qualifies the state-level branch comparison; it does not automatically prove that every differing branch is physically equivalent. State agreement and all other audits must still pass. Smooth stationarity or smooth-reference equivalence near such a branch boundary therefore remains explicitly qualified.

This verification is particularly important when a selected prediction lies close to a treatment requirement. The aeration-scheduling framework of \citet{Ren2026} similarly emphasizes influent disturbances and effluent ammonium safeguards. The relevant engineering question is the value of the particular effluent quantity at the selected controls, not only the model's average error elsewhere.

\subsubsection{Native and common-reference objectives}

Let route $S$ denote the connected surrogate and route $M$ the direct smooth mechanistic search. Their native responses are $\chi_{{\rm nat},S}=\widehat\chi$ and $\chi_{{\rm nat},M}=\chi_{{\rm sm},M}$. Their independently evaluated responses are $\chi_{{\rm ref},S}$ and $\chi_{{\rm ref},M}$. Raw and projected surrogate predictions are also evaluated, where defined, at both retained control vectors.

Normalization can alter an objective even when the physical response is identical. Consider an illustrative two-quality response whose development standard deviations are 0.5 and 2 in their native units. A scale with a floor of one gives denominators 1 and 2, whereas an unfloored scale gives 0.5 and 2. For response values 0.8 and 4,
\[
\begin{aligned}
\text{floored scale}&=(0.8/1,\,4/2)=(0.8,2),\\
\text{unfloored scale}&=(0.8/0.5,\,4/2)=(1.6,2).
\end{aligned}
\]
With equal internal quality weights, these give means 1.4 and 1.8. The difference arises entirely from the normalization. In matrix form,
\[
D_{T,S}^{\rm toy}=\begin{bmatrix}1&0\\0&2\end{bmatrix},\qquad
D_{T,M}^{\rm toy}=\begin{bmatrix}0.5&0\\0&2\end{bmatrix}.
\]
The paired comparison therefore evaluates both decisions on a common reference normalization so that a route-specific scaling convention is not mistaken for a treatment improvement.

For positive development-set quality standard deviations $\sigma_j$, the route-native normalizations are
\begin{equation}
\equationname{Route-specific water-quality objective scales}
D_{T,S}=\operatorname{diag}(\max(1,\sigma_j)),\qquad
D_{T,M}=\operatorname{diag}(\sigma_j).
\label{eq:plant_route_normalization}
\end{equation}
The unit floor is expressed in the physical unit of each quality quantity. The common reference objective uses $D_{T,M}$ for both selected decisions. The fitted standard deviations remain estimation outputs rather than theoretical constants.

A hypothetical comparison clarifies the two distinct objective differences. Suppose two native objectives are 0.90 and 0.95, so the native calculations favor the first decision by 0.05. Independent mechanistic evaluation instead gives 1.10 and 0.98, reversing the ordering. For the first route, the native-minus-reference difference is $0.90-1.10=-0.20$, indicating that the native calculation was optimistic relative to the independent evaluation. If the second route has native objective 0.95 and reference objective 0.98, its native-minus-reference difference is $-0.03$. Comparing the two independently evaluated decisions gives $1.10-0.98=0.12$, meaning that the first selected decision has the larger common-reference objective.

These two comparisons are
\begin{align}
\equationname{Difference between route-native and reference-model objectives}
\Delta J_{{\rm nat-ref},k}&=
J_{{\rm nat},k}(\vartheta_k,\chi_{{\rm nat},k})
-J_{\rm ref}(\vartheta_k,\chi_{{\rm ref},k}),\quad k\in\{S,M\},\\
\equationname{Reference-model objective difference between optimization routes}
\Delta J_{S-M}&=J_{\rm ref}(\vartheta_S,\chi_{{\rm ref},S})
-J_{\rm ref}(\vartheta_M,\chi_{{\rm ref},M}).
\label{eq:plant_objective_comparison}
\end{align}
The first quantity can reflect both response approximation and route-specific normalization. The second compares the retained decisions on one shared physical and numerical basis \citep{BhosekarIerapetritou2018,Pedrozo2025}.

Relative objective differences also depend on their denominator. If the common-reference difference is 0.12 and the direct objective is 0.98, the direct-relative percentage is approximately 12.24\%. A symmetric scaled difference using $\max(1,1.10,0.98)=1.10$ gives approximately 0.1091. The first uses one route as a reference; the second uses a symmetric magnitude scale.

Control differences require a different normalization. On an illustrative interval from 10 to 30, two selected controls 25 and 20 differ by five physical units, corresponding to normalized difference 0.25. If this is the only differing coordinate among seven controls, the RMS normalized difference is $\sqrt{0.25^2/7}\approx0.0945$, while the maximum difference remains 0.25. The RMS summarizes overall separation; the maximum keeps the largest coordinate change visible.

The normalized control differences are
\begin{equation}
\equationname{Normalized operating-control differences and their summaries}
d_j=\frac{\vartheta_{S,j}-\vartheta_{M,j}}{\vartheta_j^U-\vartheta_j^L},\qquad
d_{\rm rms}=\left(\frac17\sum_{j=1}^7d_j^2\right)^{1/2},\qquad
d_{\max}=\max_j|d_j|.
\end{equation}
These metrics distinguish similarity in the objective from similarity in the actual operating decisions.

Removal is reported on the concentration basis of the corresponding composite. If influent concentration is 200 mg L$^{-1}$ and effluent concentration is 2 mg L$^{-1}$, removal is
\[
100\frac{200-2}{200}=99\%.
\]
In general,
\begin{equation}
\equationname{Percentage removal of a reported water-quality quantity}
\mathcal R_j=100\left(1-\frac{z_{E,j}}{z_{{\rm in},j}}\right),
\end{equation}
provided the influent denominator is positive. Removal therefore compares concentrations on one reporting basis, whereas a full component inventory balance also requires flow rates and all represented outlet streams.

Predicted-to-reference removal errors are expressed in percentage points. Concentration percentage errors instead divide by the mechanistic effluent concentration. These two denominators can produce very different numerical impressions. For example, with influent 200 mg L$^{-1}$ and effluent values 2 and 4 mg L$^{-1}$, the concentration prediction error is 100\% relative to 2 mg L$^{-1}$, while the removal difference is only one percentage point. Reporting both keeps these interpretations distinct.

\subsubsection{Validity, engineering eligibility, and failures}

A paired objective comparison is reported only when both routes provide a natively feasible candidate, each candidate produces a valid exact mechanistic evaluation, and both exact responses satisfy the common retained engineering safeguards. A mechanistically valid response that violates an underflow, feed-solids, or external-loss safeguard remains diagnostically useful but does not enter the eligible paired comparison. Solids retention time and loading measures are reported for every finite mechanistic evaluation and do not create additional eligibility thresholds.

\begin{table}[htbp]
\centering\small
\caption{Distinct questions retained in failure and qualification accounting.}
\label{tab:plant_failure_classes}
\begin{tabular}{p{0.31\linewidth}p{0.55\linewidth}}
\toprule Assessment & Question answered\\\midrule
Generation acceptance & Did both original-model starts pass the required physical and numerical checks?\\
Projection feasibility & Did the fitted closure and physical constraints admit an audited joint response?\\
Native engineering feasibility & Did the search candidate satisfy its declared optimization constraints?\\
Convergence qualification & Did derivative audits or complete finite polls supply their stated evidence?\\
Exact-model validity & Did independent original-model evaluation converge to an accepted state?\\
Retained engineering feasibility & Did the exact response pass the common safeguards?\\
Paired eligibility & Were both decisions available and admissible on the common reference basis?\\
\bottomrule
\end{tabular}
\end{table}

These status fields are intentionally not collapsed into one success label. Missing or nonfinite quantities, undefined feature scales, failed factorizations, and inconsistent audits remain computational failures. A finite result that fails a scientific threshold remains available as a diagnostic state with the failed criterion identified. Generation rejection, projection infeasibility, optimization failure, and branch qualification therefore retain distinct meanings and denominators.

\subsubsection{Influent scenarios and time accounting}

Ten additional influent vectors from a separate twenty-dimensional Latin hypercube define the scenario experiment. They use the seed in Table~\ref{tab:plant_design_settings} and are excluded from fitting, trust calibration, and model tuning. Both optimization routes are attempted for the nominal influent and all ten scenarios. Every available selected decision receives independent exact evaluation, even when the other route fails. The scenario experiment therefore examines how the fixed fitted procedure behaves under specified influent variation.

Optimization time is defined as the duration of the route's initial search. For Extended ICSOR, this interval includes regression evaluation, joint projection, and any value-only fallback used within the initial search. For Smooth NLP, it includes the initial three-stage smoothing sequence. Development sampling, model fitting, convergence certification, conditional recovery, and independent exact evaluation are excluded. Failed initial attempts remain in the timing account. Scenario summaries use S1--S10 and keep the nominal case separate. Both routes use the same processor allocation and thread count.

Timing is reported together with search workload. Recorded quantities include evaluated points, iterations, projection retries, fallback use, completed smoothing stages, recovery status, and convergence qualification. The benchmarking concerns of \citet{Bliek2023} motivate retaining these workload indicators alongside elapsed time and objective quality.

A complete workflow comparison would additionally include development, certification, recovery, and exact verification. These terms matter when repeated use is considered because development effort can be amortized only if later per-decision savings are large enough. \citet{BhosekarIerapetritou2018} and \citet{Bliek2023} motivate this complete-cost perspective.

The arithmetic can be illustrated hypothetically. Suppose surrogate development requires 100 s, and each later surrogate-assisted decision requires 3 s in total for search, qualification, and verification. Suppose the comparable direct decision requires 6 s. For 20 decisions, total times are $100+20(3)=160$ s and $20(6)=120$ s. For 50 decisions, they are 250 and 300 s. The development effort is therefore recovered only after sufficiently many repeated decisions.

More generally, let $n_{dec}$ be the number of later decisions, $T_{dev}$ the development cost, and $T_S$ and $T_M$ the complete per-decision costs. Surrogate use is faster in total when
\[
\begin{aligned}
T_{dev}+n_{dec}T_S&<n_{dec}T_M,\\
T_{dev}&<n_{dec}(T_M-T_S),\\
n_{dec}&>\frac{T_{dev}}{T_M-T_S}\qquad\text{when }T_M>T_S.
\end{aligned}
\]
For the illustrative values above, the threshold is $100/(6-3)=33\tfrac13$, so at least 34 complete decisions are required. If the surrogate does not reduce complete per-decision cost, positive development effort cannot be recovered through timing alone. The comparison also assumes that the resulting decisions satisfy the required quality and feasibility criteria.

The complete assessment therefore retains four distinct forms of evidence. Statistical assessment concerns approximation of accepted mechanistic states. Projection audits concern the imposed network and inventory relations. Search audits concern local numerical evidence for the declared optimization problems. Independent original-model evaluation concerns treatment performance and feasibility at the selected controls. Together, these four layers determine whether the connected surrogate is useful as an operating decision aid \citep{BhosekarIerapetritou2018,Pedrozo2025}.

\section{Results}
\label{sec:plant_results}

The operating assessment compares Extended ICSOR and Smooth NLP through independent evaluation with the original mechanistic equations. The nominal influent is denoted N, and the ten additional influent cases are denoted S1--S10. Prediction accuracy, physical compliance, decision quality, and optimization time are reported separately so that improvement in one dimension is not treated as evidence of improvement in another.

\subsection{Accepted development and holdout states}

The development design retained 7{,}386 of 8{,}000 attempted mechanistic states, while the holdout design retained 1{,}845 of 2{,}000. Their acceptance fractions were 92.33\% and 92.25\%, respectively. Table~\ref{tab:plant_sample_results} gives the corresponding rejection counts. Rejected candidates were not replaced, so the retained samples reflect the realized acceptance behavior of the predefined designs. These accepted states supplied the fitting and holdout assessments.

\begin{table}[htbp]
\centering\small
\caption{Realized mechanistic samples for the connected-plant assessment. The separate development and holdout designs retained accepted states without replacing rejected candidates.}
\label{tab:plant_sample_results}
\begin{tabular}{lrrrr}
\toprule
Design & Attempted & Accepted & Rejected & Accepted (\%)\\
\midrule
Development & 8000 & 7386 & 614 & 92.33\\
Holdout & 2000 & 1845 & 155 & 92.25\\
\bottomrule
\end{tabular}
\end{table}

The raw-response model selected a ridge penalty of 0.1. Its mean grouped cross-validation normalized root error was 0.659 on the complete joint response using fitting-defined response scales. This development score and the holdout composite scores reported below use different denominators and should therefore be interpreted within their own assessment roles.

\subsection{Aggregate and local holdout accuracy}

Joint projection improved the aggregate holdout prediction. The pooled nRMSE decreased from 0.04690 for the raw response to 0.04306 after projection, corresponding to an 8.19\% reduction before rounding. The nMAE decreased from 0.03044 to 0.02568, a 15.65\% reduction, while mean location-level $R^2$ increased from 0.925 to 0.938. Figure~\ref{fig:plant_holdout_overview_results} shows that this improvement was not confined to the aggregate score: all eight treatment-location summaries and all four composite summaries had lower nRMSE after projection.

\begin{figure}[htbp]
\centering
\includegraphics[width=\textwidth,height=0.72\textheight,keepaspectratio]{compile/optimization/results/figures/q01_holdout_accuracy_overview.png}
\caption{Raw and projected composite accuracy on 1845 holdout plant states. Rows show aggregate, location, and composite summaries. Columns show nRMSE, nMAE, and mean location-level $R^2$. Errors are normalized by the mechanistic holdout range of each location and composite. The $R^2$ summaries average coordinate-specific scores rather than pooling locations into one score.}
\label{fig:plant_holdout_overview_results}
\end{figure}

The composite-level changes were also consistent with this general improvement. COD nRMSE decreased from 0.04741 to 0.04586, TN from 0.03724 to 0.03454, TP from 0.04170 to 0.03349, and TSS from 0.05852 to 0.05476. TP showed the largest relative improvement among the four composites, whereas TSS retained the largest projected error. Because these values use location-specific holdout ranges, they describe relative agreement within each reported treatment quantity rather than one common physical-unit error scale.

The location-level assessment shows where the projection helped and where it did not. Of the 32 location--composite combinations, projection reduced nRMSE in 30. Figure~\ref{fig:plant_holdout_locations_results} identifies the two exceptions: COD error increased by 0.47\% in the final reactor and by 2.36\% in the overflow. By contrast, overflow TP nRMSE decreased from 0.02948 to 0.01795, a reduction of 39.10\%. The connected projection therefore produced a broad statistical benefit, but the adjustment was not uniformly beneficial for every reported quantity at every treatment location.

\begin{figure}[htbp]
\centering
\includegraphics[width=\textwidth]{compile/optimization/results/figures/q02_holdout_accuracy_by_location.png}
\caption{Holdout nRMSE by treatment location and composite. Raw and projected heatmaps share their color scale. The third panel shows percentage change relative to raw error. Negative values indicate improvement. The mixer, five reactors, overflow, and underflow are assessed separately.}
\label{fig:plant_holdout_locations_results}
\end{figure}

\subsection{Parity across the connected treatment system}

Figure~\ref{fig:plant_holdout_parity_results} provides a direct view of projected agreement across the connected plant. For each composite, the figure pools 14{,}760 location-observations while retaining the median of each treatment location. Much of the point density lies near the one-to-one line. The mean location-level $R^2$ values were 0.944 for COD, 0.963 for TN, 0.945 for TP, and 0.899 for TSS. These values summarize the eight location-specific coefficients for each composite rather than one pooled coefficient across all observations.

\begin{figure}[htbp]
\centering
\includegraphics[width=\textwidth,height=0.72\textheight,keepaspectratio]{compile/optimization/results/figures/q03_holdout_parity_all_locations.png}
\caption{Projected composite parity across eight locations and 1845 holdout states. Both concentration axes are logarithmic. Hexagon color gives the number of location-observations, and labelled markers identify each location's median. The one-to-one line indicates exact agreement. Metric annotations retain location-specific normalization and average location-level $R^2$.}
\label{fig:plant_holdout_parity_results}
\end{figure}

The overflow remained more difficult to reproduce than the reactor locations. Its projected location-average $R^2$ was 0.876, compared with reactor averages of approximately 0.947. The projected overflow TSS nRMSE was 0.07392, and TSS was underpredicted in 1{,}008 of the 1{,}845 holdout states, or 54.6\%. The remaining cases had either exact agreement or overprediction. The connected response therefore captured the broad treatment pattern well while retaining a recognizable error concentration at the discharge solids endpoint.

\subsection{Physical audits and search qualifications}

The physical audit showed a sharp distinction between the raw and projected responses. Every raw holdout response violated at least one imposed mass-conservation relation, and 1{,}765 of the 1{,}845 raw responses contained at least one negative entry. After joint projection, no mass-conservation, non-negativity, or network-inequality violations were recorded under the stated audit definitions. The projection therefore closed the declared physical-consistency gap throughout the holdout set.

The optimization searches produced a different type of evidence. Extended ICSOR retained an initially feasible candidate for the nominal influent and all ten scenarios. Smooth NLP retained candidates for the nominal influent and nine of the ten scenarios. The initial S1 direct search did not produce a validated feasible candidate, and the permitted recovery did not resolve that outcome. A separately evaluated valid direct decision for S1 is nevertheless included in the treatment and objective figures so that the paired engineering comparison remains available. It is not counted as a successful initial direct search.

Table~\ref{tab:plant_search_qualification_results} separates candidate availability from convergence qualification. Extended ICSOR received positive first-order qualification in N, S6, and S10. Seven additional surrogate decisions were qualified by the finite-poll procedure, while S4 remained unresolved. None of the retained initial Smooth NLP candidates received a positive first-order certificate. These qualification results concern numerical evidence for the corresponding search problems; the independently solved mechanistic responses provide the reference treatment outcomes used in the engineering comparison.

\begin{table}[htbp]
\centering\small
\caption{Qualification of the initial center-start searches. Each route was attempted in eleven cases. Finite-poll qualification uses the specified directions and radii. The separately evaluated S1 direct decision appears in the paired figures but does not alter these initial-search counts.}
\label{tab:plant_search_qualification_results}
\begin{tabular}{p{7.0cm}rr}
\toprule
Outcome & Extended ICSOR & Smooth NLP\\
\midrule
Initial searches & 11 & 11\\
Retained feasible initial candidates & 11 & 10\\
First-order qualification & 3 & 0\\
Additional finite-poll qualification & 7 & 0\\
Retained candidates with unresolved qualification & 1 & 10\\
Cases without a retained initial candidate & 0 & 1\\
\bottomrule
\end{tabular}
\end{table}

\subsection{Prediction fidelity at selected controls}

Average holdout performance does not determine the error at the specific controls selected by an optimizer. Figure~\ref{fig:plant_selected_parity_results} therefore compares Extended ICSOR predictions directly with independent mechanistic evaluations at the surrogate-selected controls. Across the 44 case--composite comparisons, the median absolute concentration error was 12.75\% of the mechanistic effluent concentration and the mean was 18.43\%. Median and mean absolute removal errors were 3.26 and 5.23 percentage points. These summaries use a different input distribution and different denominators from the holdout metrics.

Several cases illustrate the practical meaning of the discrepancies. In S6, Extended ICSOR predicted COD of 151.10 mg L$^{-1}$, whereas the mechanistic evaluation gave 99.21 mg L$^{-1}$. In S10, the corresponding values were 232.23 and 167.78 mg L$^{-1}$. The surrogate therefore underestimated COD removal by 8.51 and 10.44 percentage points in these two cases. These errors were conservative with respect to treatment performance because the surrogate predicted worse COD removal than the mechanistic model produced.

The S4 nutrient predictions had the opposite implication. Extended ICSOR predicted TN of 2.67 mg L$^{-1}$, while the independent mechanistic value was 7.47 mg L$^{-1}$. It predicted TP of 2.08 mg L$^{-1}$ against a mechanistic value of 4.68 mg L$^{-1}$. TN and TP removals were consequently overstated by 16.97 and 27.33 percentage points. These discrepancies remained after physical projection and trust screening, showing that a prediction can satisfy the declared network constraints while still being inaccurate in an engineering quantity important to the decision.

Overflow TSS showed a different pattern. The surrogate overpredicted TSS in ten of the eleven selected decisions. Predicted values ranged from 5.53 to 8.05 mg L$^{-1}$, while mechanistic values ranged from 5.28 to 6.41 mg L$^{-1}$. In S1, the predicted and reference values were 7.78 and 5.77 mg L$^{-1}$, yet the corresponding removal values differed by only 0.61 percentage points. A modest removal error can therefore coexist with a substantial relative error in the low effluent concentration itself.

\begin{figure}[htbp]
\centering
\includegraphics[width=\textwidth]{compile/optimization/results/figures/q04_effluent_and_removal_parity.png}
\caption{Extended ICSOR predictions at selected controls compared with independent original-model evaluation. The upper row shows effluent concentration parity in mg L$^{-1}$. The lower row shows removal parity in percent. Columns correspond to COD, TN, TP, and TSS. Teal identifies concentrations and orange identifies removals. Labels identify N and S1--S10. Dashed lines indicate exact agreement.}
\label{fig:plant_selected_parity_results}
\end{figure}

\subsection{Treatment outcomes and operating choices}

Figure~\ref{fig:plant_controls_results} compares the two optimization routes after both selected decisions are evaluated with the original mechanistic model. Smooth NLP produced lower COD, TN, and TP in the nominal case and in nine of the ten scenarios. For the nominal influent, COD was 151.26 mg L$^{-1}$ at the Extended ICSOR decision and 57.60 mg L$^{-1}$ at the Smooth NLP decision. In S9, the corresponding values were 203.24 and 19.08 mg L$^{-1}$. TSS did not follow the same ordering consistently: Smooth NLP produced lower TSS in six of the eleven paired comparisons.

The selected HRT values show that the two routes also differed in their capacity choices. Extended ICSOR selected the 6-h lower bound in N, S6, S7, S9, and S10. Smooth NLP selected the 36-h upper bound in S5 and S10 within numerical tolerance. Because fresh flow was fixed, these HRT choices changed total biological reactor volume. The optimization therefore compared different combinations of installed biological capacity and operating controls.

Aeration and recycle patterns reveal additional structure. Both routes selected the upper aeration bound in reactor 3 for every case within numerical tolerance, and aeration in reactors 4 and 5 frequently reached the upper bound as well. Both routes selected the return-sludge lower bound of 0.25. Smooth NLP selected internal recycle $r_I=4.00$ in S2 and S10 and 3.44 in S9; in its other cases, the selected ratio was effectively zero. Extended ICSOR selected internal recycle values between zero and 0.804.

S2 illustrates why nominal HRT and residence time per passage must be distinguished. Smooth NLP selected $H=31.40$ h with $r_I=4.00$ and $r_R=0.25$. The resulting reactor-train flow ratio was 5.25, so the actual residence time per passage through the train was only about 5.98 h. High recycle therefore reduced the time associated with each passage despite the large nominal fresh-flow HRT.

Wasting frequently reached its upper bound, but the corresponding resource contribution still depended on the selected underflow solids concentration through $wX_U$. Neither the wasting ratio nor an aeration setting alone is sufficient to describe the resulting material removal or resource proxy.

\begin{figure}[htbp]
\centering
\includegraphics[width=\textwidth,height=0.76\textheight,keepaspectratio]{compile/optimization/results/figures/q05_effluent_and_operating_values.png}
\caption{Mechanistic effluent outcomes and selected controls for both optimization routes. The top row shows COD, TN, TP, and TSS in mg L$^{-1}$. The remaining panels show HRT, aeration in reactors 3--5, internal recycle, return sludge, and wasting. The unused panel is omitted. N and S1--S10 appear in every comparison. The constant aeration and return-sludge panels use plain zero-to-one axes.}
\label{fig:plant_controls_results}
\end{figure}

\subsection{Mechanistic treatment-train profiles}

Figure~\ref{fig:plant_profile_results} traces the original-model concentrations through the plant at each route's selected controls. All four reported composites increased from fresh influent to the mixer in all eleven cases for both routes. This increase reflects the material returned by the recycle streams rather than an external source of pollution \citep{Alex2008}.

The subsequent changes through the biological reactors and the clarifier have different physical meanings. Changes across the reactors reflect the combined effects of biological conversion, hydraulic transport, and solids inventory. The concentration decrease across the clarifier reflects separation of particulate material under the nonreactive clarifier assumption. A reported total may also remain nearly constant while its internal species are redistributed. For example, one nitrogen form may be converted to another without producing a large change in TN \citep{Henze2006,JeppssonDiehl1996}.

The pronounced TSS decrease from reactor 5 to the final effluent shows the role of secondary clarification in treated-water quality. Interpreting the plant response therefore benefits from retaining the individual species, outlet solids flows, and reactor and clarifier inventories in addition to the four reported composites. These quantities connect the observed profiles to nitrification, denitrification, phosphorus uptake, separation, and solids age \citep{Henze2008,Wentzel1992}.

\begin{figure}[htbp]
\centering
\includegraphics[width=\textwidth,height=0.77\textheight,keepaspectratio]{compile/optimization/results/figures/q06_treatment_train_profiles.png}
\caption{Mechanistic concentration profiles at controls selected by Extended ICSOR and Smooth NLP. Rows show COD, TN, TP, and TSS from fresh influent through the mixer and reactors 1--5 to clarifier effluent. Each composite has identical logarithmic concentration limits for both routes. Colors identify N and S1--S10. Underflow is not placed in the main liquid-treatment path.}
\label{fig:plant_profile_results}
\end{figure}

\subsection{Quality and resource trade-offs}

Smooth NLP produced the lower common-reference objective in the nominal case and in all ten scenarios. Figure~\ref{fig:plant_objective_results} separates the total objective into normalized treatment quality and weighted resource contributions so that the reason for each preference remains visible. For the nominal influent, the objectives were 0.8647 for Extended ICSOR and 0.7862 for Smooth NLP, making the surrogate-selected value 9.99\% higher relative to the direct value. S5 had the largest absolute difference, with values 1.0101 and 0.7022. S8 had the smallest, with values 0.9279 and 0.9234.

The lower Smooth NLP objectives did not always arise from the same trade-off. Smooth NLP had lower normalized quality cost in ten of the eleven cases, but its weighted resource contribution was higher in every case except S8. In S10, normalized quality cost decreased from 1.4748 to 0.3699, while weighted resource contribution increased from 0.0557 to 0.4191. The resulting total objective decreased from 0.7931 to 0.6040. The direct decision therefore accepted substantially greater capacity and resource contribution to obtain much better weighted treatment performance.

S8 showed the opposite pattern. Its normalized quality cost increased from 1.5856 to 1.6110, and all four evaluated effluent composites were slightly higher under the Smooth NLP decision. The weighted resource contribution, however, decreased from 0.1351 to 0.1179. That resource reduction was sufficient to lower the total objective from 0.9279 to 0.9234. The objective therefore preferred the direct decision even though its effluent concentrations were slightly worse.

The development quality scales were 106.61, 13.07, 7.11, and 17.38 mg L$^{-1}$ for COD, TN, TP, and TSS. Because all four exceeded one, the route-specific normalization matrices in Equation~\eqref{eq:plant_route_normalization} were identical in this experiment. The objective differences were therefore not caused by different route-specific quality scaling.

The dependence of the aeration proxy on reactor volume also influenced the trade-offs. When $H=6$ h and all three aeration settings are at one, Equation~\eqref{eq:plant_aeration_proxy_example} gives $j_A=1/6$. Several surrogate-selected decisions combined maximum aeration settings with minimum HRT and therefore incurred relatively small normalized aeration-capacity contributions. The objective measures the declared capacity and resource proxies rather than the dimensionless aeration controls alone.

\begin{figure}[htbp]
\centering
\includegraphics[width=\textwidth]{compile/optimization/results/figures/q07_objective_quality_economic.png}
\caption{Common-reference total objective, normalized quality cost, and weighted resource contributions. Both routes are evaluated with the original mechanistic model and common quality scales. Resource stacks show HRT, aeration, internal recycle, return sludge, and wasting. Solid bars identify Extended ICSOR, and hatched bars identify Smooth NLP. The quality panel shows $j_Q$, whose contribution to the total objective is $0.50j_Q$.}
\label{fig:plant_objective_results}
\end{figure}

\subsection{Optimization time}

Extended ICSOR completed the recorded initial search more quickly for the nominal influent and for every scenario. Figure~\ref{fig:plant_time_results} shows all eleven comparisons. Across S1--S10, Extended ICSOR times ranged from 0.758 to 2.702 s, whereas Smooth NLP times ranged from 3.932 to 8.253 s. Their scenario means were 1.726 and 5.236 s, respectively, giving a ratio of approximately 3.03. The nominal times were 0.985 and 5.062 s.

The direct-to-surrogate time ratio varied by scenario from 1.90 to 7.99, with the largest ratio in S10. The timing definition includes failed initial search attempts, including the failed S1 direct attempt. Whether a feasible candidate was retained is therefore reported separately in Table~\ref{tab:plant_search_qualification_results} rather than inferred from timing.

These times cover only the initial optimization intervals defined in the methodology. They exclude development sampling, fitting, convergence certification, conditional recovery, and independent mechanistic evaluation. Extended ICSOR therefore demonstrated a shorter initial search interval, but the paired engineering assessment simultaneously showed that its selected decisions had higher mechanistic objectives.

\begin{figure}[htbp]
\centering
\includegraphics[width=\textwidth]{compile/optimization/results/figures/q08_optimization_time.png}
\caption{Recorded optimization time for N and S1--S10 on a logarithmic seconds axis. The interval includes the initial surrogate search with projection and fallback or the initial direct smoothing sequence. It excludes generation, fitting, certification, conditional recovery, and independent mechanistic evaluation. Initial failed attempts remain in the timing account. Scenario mean times use S1--S10 only.}
\label{fig:plant_time_results}
\end{figure}

\section{Discussion}
\label{sec:plant_discussion}

\subsection{Connected enforcement and its statistical benefit}

The connected surrogate was designed to predict the plant as one material system rather than as a collection of independent unit outputs. Its response retained the mixer, five biological reactors, clarifier overflow and underflow, and stored solids. Joint projection then linked these quantities through the declared mixer, reactor-inventory, clarifier, separation, and storage relations. The holdout results show that this physical reconciliation was accompanied by a broad statistical benefit. Aggregate nRMSE and nMAE both decreased, all eight location summaries improved, all four composite summaries improved, and 30 of the 32 individual location--composite errors decreased.

This result is important because projection was not introduced merely as a numerical feasibility repair. The connected constraints couple locations that are physically related, so correcting one part of the response redistributes information across the treatment system. In this case, the redistribution generally moved the predictions closer to the mechanistic holdout states. The especially large reduction in overflow TP error illustrates how a network-level correction can improve a reported quantity at a downstream location even though the projection acts on the complete component response.

The improvement was nevertheless not uniform. COD error increased slightly in the final reactor and overflow, and the overflow remained the least accurate treatment location overall. The separate logarithmic model supplied a positive overflow-solids endpoint, while the joint projection reconciled that endpoint with component balances, soluble continuity, particulate densification, and inventory limits. The resulting response therefore reflected both physical equalities and an empirical closure. Equations~\eqref{eq:plant_projection_error_exact} and \eqref{eq:plant_projection_error_general} explain why error reduction depends on the projection metric and on whether the mechanistic reference lies in the complete projection set. The two local COD increases are consistent with that limitation: enforcing network consistency can improve the response in aggregate while worsening an individual reported coordinate.

The practical implication is that physical consistency and statistical accuracy should be assessed together but not conflated. The projection succeeded completely in its declared physical role, while its statistical effect was broad but not universal. This distinction is consistent with the projection analyses of \citet{SturmSilva2024} and \citet{KircherVotsmeier2025}. For connected treatment models, the relevant question is therefore not simply whether projection enforces the constraints, but also how the resulting adjustment redistributes predictive error across the locations and quantities that matter for engineering use.

\subsection{Holdout accuracy and selected-decision fidelity}

The holdout assessment showed that Extended ICSOR approximated the connected plant well over the sampled domain, but the selected-control assessment exposed a second and more demanding question. Optimization does not sample the domain uniformly. It concentrates evaluations around controls that appear favorable under the surrogate objective. Prediction error at those selected controls can therefore matter more to the final decision than average error over the broader holdout design \citep{BhosekarIerapetritou2018,Pedrozo2025}.

S4 illustrates this distinction clearly. The projected surrogate predicted TP of 2.08 mg L$^{-1}$, whereas independent mechanistic evaluation at the same controls gave 4.68 mg L$^{-1}$. TN was similarly underestimated, at 2.67 versus 7.47 mg L$^{-1}$. If, for illustration, a treatment requirement were 3 mg L$^{-1}$ for TP, the surrogate and mechanistic responses would place the same operating decision on opposite sides of the requirement. The engineering consequence is therefore not captured by the relatively favorable average holdout statistics.

The selected-control discrepancies also show why physical projection and trust screening cannot be treated as substitutes for mechanistic verification. Projection ensured consistency with the declared plant balances and bounds, while the trust diagnostics restricted search to regions of specified data support, correction size, split consistency, and reactor residual. Neither step guarantees that the surrogate reproduces the nonlinear biological and settling equations closely enough at every optimizer-selected point. A physically admissible, trusted response can still be quantitatively wrong in a nutrient concentration that governs the operating decision.

Independent evaluation closes this gap by calculating treatment outcomes from the original biological kinetics and nonlinear clarifier model at the selected controls. In this workflow, the surrogate is used to generate and screen candidate decisions, while the mechanistic model remains the reference used to judge them. The approximation-management work of \citet{Hameed2026} and the treatment-safeguard setting of \citet{Ren2026} support this separation between rapid candidate selection and reference evaluation. The results therefore support a verified decision workflow rather than unrestricted substitution of the mechanistic plant model.

\subsection{Decision quality under declared priorities}

The common-reference objective comparison gives the clearest assessment of the operating decisions themselves. Smooth NLP produced a lower mechanistic objective in the nominal case and in all ten influent scenarios. Extended ICSOR, by contrast, consistently completed the recorded initial search more quickly. The two routes therefore exhibited a trade-off between initial search computation and the quality of the retained local decisions under the declared priorities.

The objective decomposition shows that this difference cannot be summarized as ``better treatment'' versus ``lower resource use'' in one universal direction. In S10, Smooth NLP accepted a much larger capacity and resource contribution to obtain a much lower quality cost, and the treatment gain dominated the total objective. In S8, the direct route instead accepted slightly higher effluent concentrations but reduced the weighted resource contribution enough to achieve a lower total objective. Both outcomes are internally consistent with the same declared weights. The objective prefers different treatment-resource combinations because the relative contributions differ by scenario \citep{PadronPaez2020,Aboagye2021}.

This decomposition is especially important because HRT is a free decision in the present case. At fixed influent flow, increasing HRT increases biological reactor volume. Aeration cost also depends on that selected volume through the transfer-capacity proxy. The direct and surrogate routes therefore differed not only in aeration, recycle, and wasting, but also in the treatment capacity they selected. Frequent 6-h surrogate solutions favored smaller reactor volume and correspondingly smaller aeration-capacity contributions, whereas several direct solutions accepted much larger HRT to improve effluent quality.

The results should therefore be interpreted as local decisions under the specific capacity and resource priorities defined for the case study. An installed treatment plant would impose fixed reactor volumes and equipment capacities, changing the decision problem substantially. The component-resolved analysis of \citet{Zhang2026} and the energy treatment of \citet{Miederer2025} support explicit accounting of hydraulic capacity, oxygen-transfer capacity, and resource demand when interpreting such operating trade-offs.

\subsection{Computational performance and operating assessment}

The shorter Extended ICSOR search times are consistent with the mathematical structure of the two routes. The surrogate search varies only seven operating controls and obtains plant responses from second-order regression followed by a convex joint projection. Smooth NLP optimizes the controls together with reactor and clarifier states and must satisfy the smoothed mechanistic equations through a continuation sequence. The surrogate therefore performs a less expensive response calculation during the initial search.

That computational advantage is useful for repeated exploration. A fitted surrogate can evaluate many alternative aeration, recycle, wasting, and capacity combinations quickly, which is valuable when the objective is to screen candidate regions or compare responses across several influent conditions. The measured scenario mean search time was approximately three times shorter for Extended ICSOR than for Smooth NLP under the stated timing definition.

Search time, however, is only one part of the workflow. Development sampling, surrogate fitting, lower-problem projection audits, convergence qualification, conditional recovery, and independent mechanistic verification all add computational work. These terms were intentionally excluded from the reported initial-search interval and therefore cannot be inferred from Figure~\ref{fig:plant_time_results}. The complete value of the surrogate depends on whether the development cost can be amortized across repeated decisions and whether the resulting mechanistically verified decisions remain acceptable. This broader accounting follows the benchmarking principles discussed by \citet{Bliek2023}.

The convergence evidence also remains separate from decision quality. First-order qualification and finite polling assessed local numerical behavior of the surrogate optimization problem, while the direct route retained unresolved qualification in its center-start solutions. These statuses do not determine which operating decision is better. The original-model evaluation supplies that comparison. Conversely, a favorable mechanistic objective does not retroactively establish that a search was globally optimal or numerically certified.

Taken together, the results define a practical role for the connected surrogate. Extended ICSOR provides rapid candidate generation, explicit material accounting, and a fully auditable physical correction. Its search can screen operating alternatives substantially faster than the direct route under the recorded interval. The original mechanistic model then remains necessary to establish the treatment outcome and objective at the selected controls. The resulting decision aid therefore combines rapid exploration with explicit verification rather than treating computational speed as evidence of equivalent decision quality.

\subsection{Data Availability}
\label{sec:plant_data_availability}

The simulation data and numerical results supporting this study are publicly available at \url{https://github.com/egbertoselerio1997/surrogate-optimization}.